\documentclass[11pt]{article}
\usepackage[margin=1in]{geometry}
\usepackage[T1]{fontenc}
\usepackage{lmodern,amsmath,amssymb,amsthm,booktabs,array,graphicx,longtable}
\usepackage[colorlinks=true,linkcolor=blue,citecolor=blue,urlcolor=blue]{hyperref}
\usepackage{microtype}
\hypersetup{pdftitle={Polynomial depth: unrestricted Gaussian learning and numerical stability},pdfauthor={ChatGPT-6 Astra; Claude Fable 5 and Opus 5}}
\newtheorem{theorem}{Theorem}[section]
\newtheorem{proposition}[theorem]{Proposition}
\newtheorem{lemma}[theorem]{Lemma}
\newtheorem{corollary}[theorem]{Corollary}
\newtheorem{definition}[theorem]{Definition}
\newtheorem{remark}[theorem]{Remark}
\newcommand{\E}{\mathbb E}
\newcommand{\Pp}{\mathbb P}
\newcommand{\Prb}{\mathbb P}
\newcommand{\R}{\mathbb R}
\newcommand{\X}{\{-1,1\}^d}
\newcommand{\op}{\mathrm{op}}
\newcommand{\Lip}{\operatorname{Lip}}
\newcommand{\sgn}{\operatorname{sign}}
\newcommand{\poly}{\operatorname{poly}}
\title{Polynomial depth in Bayesian neural-network learning:\\unrestricted Gaussian priors, trained networks,\\and numerical stability}
\author{Research report prepared by ChatGPT-6 Astra and Claude (Fable 5 and Opus 5)\thanks{Authors are listed in alphabetical order.}\\for Kaarel H\"anni}
\date{Research report --- 17 October 2026}
\begin{document}
\maketitle
\begin{abstract}
We study the depth dependence of exact Bayesian learning of Boolean
functions with fully connected tanh networks and a linear readout.
At width $N$, the independent Gaussian weights have variances
$t^2/d$, $t^2/N$, and $t^2/N^2$ in the input, internal, and readout layers.
The scale is a prescribed schedule $t=t(d,L,N)$; in the power-law cases
$t=N^{(1-\gamma)/2}$. Learning conditions on noiseless labels.
Unless stated otherwise, the prior also requires an absolute output
margin greater than $1/2$ at every Boolean input. Without an additional robustness restriction,
a Gaussian-distance argument gives an all-width finite-variance theorem for
every fixed exponent $\gamma>1$. Its target-dependent cost is polynomial in
$Q_*=[(K+2)/L]^{L/2}$, where $K$ is the energy of a unit-margin target
witness. A zero-padding and scalar-transmission construction makes this cost
polynomial for every bounded-weight one-hidden-layer target, uniformly
over all learner depths and every width above the actual depth-specific
complexity. In particular, parity, symmetric Boolean functions, and
polynomial-size DNF/CNF have polynomial sample bounds for every fixed
$\gamma>1$. Any fixed intrinsic target depth is also covered, with
polynomial degree depending on that depth.
Linear thresholds with inverse-polynomial geometric margin
also have joint polynomial bounds. It does not establish a joint polynomial
bound for every small circuit or give an all-width mean-field theorem.
We also give an elementary polynomial-width result uniform in all
$\gamma\ge1$. At $\gamma=1$, a separate construction shows that absolute
prior-cost counting has exponential depth dependence. At a common depth
$O(d^2)$, every Boolean function even has critical rate bounded by a
universal constant. These are not simple-target sample lower bounds.
For alternative priors, numerical output stability removes
the all-input clean-margin requirement: a margin is required only on the
training examples, and the test guarantee holds at a smaller margin.
Trusted input regions and soft stability penalties give further variants.
Experiments on synthetic Boolean tasks and small handwritten-digit networks
separate near-boundary failures from stability of confident predictions.
Explicit gradient-descent basins show that all-input sign robustness against small internal disturbances
is not a universal consequence of successful training.
\end{abstract}

\paragraph{Status.}
This report contains mathematical proofs, independent adversarial proof
reviews by other AI agents, and reported experiments.
The empirical section states the limits of the publicly packaged reproduction
materials; this writing review did not rerun the training experiments. This report is not itself a Lean formalization. Pointwise sufficient certificates generally use floating-point arithmetic.
Twenty global positive certificates and selected negative witnesses are
additionally verified with interval arithmetic and exact rational checks. Bounds are qualitative polynomial bounds
with deliberately loose constants, not practical sample estimates.
The foundational posterior-compression proof and Gaussian target-cylinder
construction are reproduced in the
appendices, so the statistical statements can be read independently.
No novelty claim is made for standard finite-parameter counting,
product-norm Rademacher methods, Gaussian concentration, or the general
idea of learning through stability and compression.

\paragraph{Scope and later developments (29 October 2026).}
The report records the results and experiments of 17 October.
Later work gives an \href{https://github.com/kaarelh/AI-works/blob/main/neural-network-generalization/main-results/meanfield-all-width/paper.pdf}{all-width mean-field theorem at fixed depth} and a \href{https://github.com/kaarelh/AI-works/blob/main/neural-network-generalization/works-in-progress/conditional-nonpolynomial-learning/paper.pdf}{conditional obstruction to a universal joint-polynomial bound}. The latter assumes
$\mathsf P/\mathrm{poly}\ne\mathsf{TC}^1$; it does not prove an
exponential depth lower bound. These developments should be distinguished
from this report's unconditional obstructions to particular proof methods.

\paragraph{Reading guide.}
Section~\ref{sec:model} defines the architecture, posterior, complexity
conventions, and reusable estimates. The subsequent sections treat
unrestricted-prior upper bounds and rate-counting obstructions, then
trained-network examples and stability-restricted alternatives. The
appendices supply the critical variational formula, target-cylinder
construction, compression theorem, and stability certificate used throughout.

\section{Results and limitations}

Throughout, ``unrestricted'' means no additional internal robustness
restriction. It retains the all-input output-margin conditioning defined above
unless a theorem explicitly changes that convention.

\begin{enumerate}
\item \textbf{Unrestricted prior.}
If the learner width is polynomially bounded in the target parameters,
the sample bound is polynomial in depth for all $\gamma\ge1$.
At arbitrary widths, a stronger result holds for each fixed
$\gamma\ge1+\zeta$, $\zeta>0$: a unit-margin witness of energy $K$
gives a bound polynomial in $d,L,C,K,Q_*$ and the statistical parameters.
The degree may depend on $1/\zeta$. For every bounded-weight shallow
target, a new representation makes $Q_*$ polynomial uniformly over
learner depth. Consequently XOR, all symmetric Boolean targets, and
polynomial-size DNF/CNF have joint polynomial bounds at every learner
depth and every admissible width. Linear thresholds with
inverse-polynomial geometric margin also have such bounds. For an
arbitrary fixed deeper implementation, the bound is polynomial in added
depth but can still be exponential in the original implementation energy.
The result also extends to any fixed intrinsic target depth, with polynomial
degree depending on that depth. A single polynomial jointly in target
size and intrinsic depth, and arbitrary-width $\gamma=1$ learning,
are not established by this report; see the later developments above.
More quantitatively, if a target representation has depth $D_0$, width $C$,
and weight bound $R$, the unrestricted $\gamma>1$ sample bound is
$(2+d+C+R)^{O_\zeta(D_0)}$ times a polynomial in learner depth and the
statistical parameters, for every $L\ge D_0$. Logarithmic target depth
therefore gives a quasipolynomial target-size bound.
At $\gamma=1$ we additionally prove an obstruction to a common proof
strategy: at depth $L=20(d+1)+4$, every Boolean function on $d$ bits has
critical prior cost at most $100000L$, with a direct finite-width
Gaussian-mass bound. Thus counting all functions below an absolute
prior-cost threshold necessarily incurs an exponential-in-depth term.
By using small covariance changes instead of mean shifts, we further show
that at one common depth $L_d=O(d^2)$, every Boolean function has critical
rate at most a universal constant. Thus an absolute-rate counting argument
fails even at a fixed rate threshold.
These results do not compare competing functions with a simple target and
are not sample lower bounds.
\item \textbf{Gradient-trained networks.}
The all-cube sign-robustness condition can fail because of small clean outputs
away from the data, even when observed predictions are accurate and stable.
Positive finite experiments do not establish a universal training law.
We prove a stronger negative statement: there are open basins of gradient
descent converging to exact, low-energy, internally fragile fits.
Parameter flatness alone also does not imply internal activation stability.
\item \textbf{Other conditions.}
The numerical-stability condition controls the change in the numerical output
rather than requiring its sign to be stable everywhere. With positive
training margins, it gives a joint polynomial theorem for every
$\gamma\ge1$ and every width above a layered Boolean circuit's width,
without an all-input clean margin. It can be localized to a trusted input
region or imposed as a soft prior penalty. The region can be chosen and
validated using independent unlabeled samples. Networks may alternatively
have their own exceptional regions from a family of controlled input-set
complexity, at an explicit error cost. These are specified changes of
prior or likelihood, rather than conclusions about unchanged Gaussian
weights.
\end{enumerate}

A learning lower bound for the unrestricted prior must be target-specific. A large class
of low-energy functions is not by itself a lower bound for this Bayesian
learner: one must compare their Gaussian posterior mass with that of a
small target. Conversely, fragility of one implementation does not imply
that its Boolean function intrinsically needs many samples. The same
function may have a robust implementation. In fact the fragile trained
dictator construction below has energy at most $L+76$, a class already
controlled by a polynomial-depth product-norm bound. Internal fragility
alone is therefore insufficient evidence for a sample lower bound.

Related representation-cost work studies infinite-depth limits for
homogeneous networks~\cite{jacot2023} and depth separation in
norm-controlled ReLU learning~\cite{parkinson2024}. These results provide
useful comparisons, but their activation functions and learning rules
differ from the tanh Gaussian posterior here.

\section{Model, conventions, and reusable statistical ingredients}
\label{sec:model}

Let $d\ge1$, $L\ge2$, and $N\ge1$. Inputs lie in $\X$ and
$h^0=(1,x)$. There are $J=L-1$ hidden layers of width $N$, no biases after
the augmented input, and a linear readout:
\[
 h^s=\tanh(W^sh^{s-1})\quad(1\le s<L),\qquad z_w(x)=v^Th^{L-1}.
\]
The Gaussian energy and readout energy are
\[
 \kappa_N(w)=\frac dN\|W^1\|_F^2+
       \sum_{s=2}^{L-1}\|W^s\|_F^2+N\|v\|_2^2,
 \qquad E_v=N\|v\|_2^2.
\]
For a prescribed scale schedule $t=t(d,L,N)>0$, the raw prior
$\mu_{N,t}$ has independent entry variances
$t^2/d,t^2/N,t^2/N^2$, respectively. Its density in energy coordinates
is proportional to $\exp[-N\kappa_N/(2t^2)]$.
Put $\sigma^2=t^2/N$; exponent $\gamma$ means $t^2=N^{1-\gamma}$.
Our unified robustness theorems use $0<t\le1$, hence all $\gamma\ge1$.

The baseline prior is this Gaussian conditioned on
$|z_w(x)|>1/2$ at \emph{every} Boolean input. Given independent training
inputs $x_i\sim D$ with $y_i=f(x_i)$, condition on correct training signs.
The resulting posterior is written $Q_S$.
Its Gibbs error is the expected test sign error of a posterior draw.
If $p_S(x)$ is the posterior probability of the correct label, then
\[
 D\{x:p_S(x)\le q\}\le
 \frac{\operatorname{GibbsErr}(Q_S)}{1-q}\quad(0<q<1).
\]
Thus the upper bounds here control both confidence thresholds $2/3$ and
$1/2$, with factors three and two. Unlike certain lower-bound arguments,
these upper bounds allow arbitrary posterior mixtures.

For robustness, use additive postactivation and output disturbances
\[
 h^s_\eta=\tanh(W^sh^{s-1}_\eta)+\eta^s,
 \quad z_w(x;\eta)=v^Th^{L-1}_\eta+\eta^L,
 \qquad
 \|\eta\|_*^2=\sum_{s<L}\frac{\|\eta^s\|_2^2}{N}+|\eta^L|^2.
 \tag{1}\label{eq:noisenorm}
\]
The reference sign-robustness radius is
$r_{\rm ref}(w)=.01/\sqrt{1+JE_v(w)}$.
A second calibration uses
$S_N=(d/N)\|W^1\|_F^2+\sum_{2\le s<L}\|W^s\|_{\rm op}^2+E_v$
and $r_S=.01/\sqrt{1+S_N}$.
A network is sign-robust at an input if every disturbance in the stated
closed ball retains its clean sign strictly. Clean zeros fail this condition.
Set $\operatorname{sign}(0)=+1$ for auxiliary classifiers only.
Write $\mathcal R_{r,N,f}$ for width-$N$ networks with clean signed
margin greater than $1/2$ for $f$ and sign robustness at radius $r$
on the entire cube.

\paragraph{Circuit terminology.}
For the robustness theorems, a width-$C$ layered Boolean circuit has
$L-1$ stages over fan-in-two AND, OR, NOT, COPY, and constant gates,
including a constant $+1$ wire at every stage, including the final
hidden stage. The width counts this wire and any gate-compilation overhead. This is distinct from the minimum width of
an arbitrary analog tanh network at exactly the same depth. Whenever the
latter is used, the needed margin and weight bound are stated explicitly.
On inputs in $\{-1,1\}$, the preactivations
$u+v-1$, $u+v+1$, $-u$, and $u$ implement the four nonconstant gates with
signed margin at least one. A constant wire supplies the apparent biases.
Multiplying these gate preactivations by a common gain and applying tanh
gives the saturated target construction in Appendix~\ref{sec:cylinder}.

\paragraph{Posterior compression and target mass.}
Appendix~\ref{sec:robust-entropy} proves that a law $Q$ with
$D_{\rm KL}(Q\Vert\mu_{N,t})\le NK/t^2$ can be approximated at each input
by a common law on a finite network class with
$\log|\mathcal H|\le\poly(d,L,K,r^{-1},\alpha^{-1})$, provided its
networks are sign-robust at radius $r$. The coupling can depend on the
input; its marginals do not. Before robustness is used, the candidate
output is exactly the original network under a small simultaneous
postactivation perturbation except on probability at most $\alpha$.
This stronger construction is what permits the pointwise and numerical
variants below. Compression through noise stability and all-layer margins
has precedents in~\cite{arora,weima2020}; the exact normalizations and
Gaussian conditioning conventions here matter.

For every layered digital target, Appendix~\ref{sec:cylinder} supplies a
raw Gaussian target cylinder of mass $\exp[-ND_*/t^2]$, with
$D_*=\poly(d,L,C)$, uniformly at all $N\ge C$ and $0<t\le1$.
Its networks have clean target margin greater than one and internal
robustness radius inverse polynomial in $C$.
If a posterior is the raw Gaussian conditioned on an event containing
this cylinder, its KL is at most $ND_*/t^2$. Also
$E_v=(t^2/N)\chi_N^2$ under the raw prior. For
\[
 B=\max\{8,8[D_*+\log(1/\beta)]\},
\]
its posterior mass on $E_v>B$ is at most $\beta$.
The retained conditional posterior has KL at most $N(D_*+1)/t^2$
for $\beta\le1/2$. These facts follow by the chi-square exponential
moment and division by the target probability bound.
They explain why a radius calibrated to readout energy gives uniform
all-width learning without imposing a hard energy cutoff.

\section{Unrestricted learning at bounded widths}
\label{sec:finitewidth-unrestricted}

A robustness restriction is unnecessary when the learner's width is itself
polynomially bounded. The following elementary argument retains the original
all-input clean margin and Gaussian conditioning conventions, and is uniform
in the Gaussian scale. It also covers analog target networks at their original
depth, without digital simulation.

Let $P_N=N(d+2)+(L-2)N^2$ be the parameter count, and let $\mu_{N,t}$ denote
the raw Gaussian with variances $t^2/d,t^2/N,t^2/N^2$, respectively. Thus
$\kappa_N=(t^2/N)\chi^2_{P_N}$ under $\mu_{N,t}$. Fix integers
$1\le C\le M$ and suppose, for some $D\ge1$, that
\begin{equation}\label{eq:finitewidth-targetmass}
 \mu_{N,t}\{w:f(x)z_w(x)>1/2\text{ for every }x\}
 \ge \exp(-DN/t^2).
\end{equation}
This holds for $C\le N\le M$ and $0<t\le1$. Only this target-mass
hypothesis is needed; posterior networks may be fragile.

\begin{theorem}\label{thm:finitewidth-unrestricted}
For $0<\varepsilon,\delta<1$, set
\begin{align*}
 B&=4D+2[d+2+(L-2)M]\log2+4\log(2/\varepsilon),\\
 A&=\sqrt{MB},\qquad T=d+1+M(L-1),\qquad
 q=[4T(MA)^{L-1}]^{-1}.
\end{align*}
If
\[
 n\ge\frac2\varepsilon\left[
 \log M+P_M\log(5A/q)+\log(1/\delta)\right],
\]
then, with probability at least $1-\delta$ over independent examples from any
input distribution, the unrestricted Gaussian posterior has Gibbs error at
most $\varepsilon$, simultaneously for every $C\le N\le M$ and $0<t\le1$.
Here ``unrestricted'' retains the original conditioning on
$|z_w(x)|>1/2$ for all inputs, but imposes no internal robustness condition.
The same sample event works after any number of additional correctly labeled
examples. In particular, if $D$ and $M$ are polynomial in $d,L,C$, the sample
bound is polynomial in $d,L,C,\varepsilon^{-1},\log(1/\delta)$, uniformly for
all $\gamma\ge1$ with $t^2=N^{1-\gamma}$.
\end{theorem}

\begin{proof}
For every sample, the event used to condition the raw Gaussian contains the
target event in~\eqref{eq:finitewidth-targetmass}. The chi-square exponential
moment therefore gives
\[
 Q_{S,N,t}\{\kappa_N>B\}
 \le\exp\left(\frac{DN}{t^2}+\frac{P_N}2\log2-
                       \frac{NB}{4t^2}\right)
 \le\varepsilon/2.
\]
The last inequality uses $P_N/N\le d+2+(L-2)M$, $N\ge1$, and $t\le1$.
This bound holds for every sample, without a probabilistic qualification.

If $\kappa_N\le B$, every physical parameter has magnitude at most $A$.
Round each to the nearest multiple of $q$, using at most $5A/q$ grid values
per coordinate. For the maximum coordinatewise hidden-state discrepancies,
boundedness and the Lipschitz property of tanh give
\[
 e_1\le(d+1)q,\qquad e_s\le NAe_{s-1}+Nq,\qquad
 |z-z_{\rm grid}|\le NAe_{L-1}+Nq.
\]
Consequently, uniformly over Boolean inputs,
\[
 |z-z_{\rm grid}|\le q[d+1+M(L-1)](MA)^{L-1}=1/4.
\]
Rounding thus preserves the entire Boolean function of every posterior
network with energy at most $B$.

The union of these grid classifiers over $C\le N\le M$ has logarithmic
cardinality at most $\log M+P_M\log(5A/q)$. A consistency union bound gives
one sample event on which every consistent grid classifier has true error
at most $\varepsilon/2$. Every good-energy posterior draw has the same
function as such a classifier; the remaining posterior mass is at most
$\varepsilon/2$. This proves the Gibbs bound. The finite class is independent
of $t$, and consistency only becomes more restrictive as examples are added.
\end{proof}

The logarithm of the grid size satisfies
\[
 \log(5A/q)=\log(20AT)+(L-1)\log(MA).
\]
The exponentially small parameter precision costs only a polynomial in
$L$ in this logarithm. The unresolved unrestricted problem is therefore
uniformity at learner widths \emph{larger than every polynomial in the
specified target complexity and depth}; this theorem does not show those
larger widths actually require more samples.

\begin{corollary}[Same-depth bounded-weight analog targets]
Suppose a width-$C$, depth-$L$ tanh network realizes $f$ with
$f(x)z(x)\ge1$ on every Boolean input and weights of magnitude at most
$R\ge1$. Fix $M\ge C$ and set
\begin{align*}
 K&=2[d(d+1)+(L-1)C^2]R^2,\\
 a&=[4(d+1+M(L-1))(MR)^{L-1}]^{-1},\\
 D&=\max\{1,K+[d+2+(L-2)M][2+\log(1/a)]\}.
\end{align*}
Then~\eqref{eq:finitewidth-targetmass} holds with this $D$.
Thus Theorem~\ref{thm:finitewidth-unrestricted} supplies a polynomial sample
bound at the original depth for every $C\le N\le M$, with dependence on
$R$ only through $\log R$ in the final sample bound. No robustness condition
on the target implementation or prior support is needed. A target with the
original strict margin $>1/2$ satisfies the displayed hypothesis after
doubling its readout.
\end{corollary}

\begin{proof}
Clone the $C$ neurons into balanced nonempty groups of sizes $n_i$ summing
to $N$. Copy first-layer rows and replace internal and readout weights by
$W_{ij}/n_j$ and $v_i/n_i$. This preserves the exact function. Since
$\max_i n_i/\min_i n_i\le2$, its energy is at most $K$, and its weights
remain bounded by $R$. A physical-coordinate perturbation of magnitude
at most $a$ changes every output by at most $1/4$, by the recurrence above.

Use isotropic coordinates
$y=(\sqrt d\,W^1,\sqrt N\,W^{\rm internal},Nv)$, so
$y\sim\mathcal N(0,t^2I)$ and $\|y_*\|^2\le NK$ for the cloned center.
The standard-normal-coordinate box
$|y_j/t-y_{*,j}/t|\le a$ lies within the physical box, because all coordinate
scalings are at least one and $t\le1$. Its probability is at least
\begin{align*}
 &\left(\frac{2a}{\sqrt{2\pi}}\right)^{P_N}
 \exp\left[-\frac12\sum_j(|y_{*,j}|/t+a)^2\right]\\
 &\qquad\ge\exp[-NK/t^2-P_N(2+\log(1/a))]
 \ge\exp(-DN/t^2).
\end{align*}
Here $(u+a)^2\le2u^2+2a^2$ and $a\le1$ suffice. Every network in this box
has target margin at least $3/4$.
\end{proof}

This is an elementary finite-parameter learning result, rather than a new
claim about the general VC dimension of neural networks. Its role is to
isolate the width regime not covered by this finite-parameter argument.

\section{A width-independent tanh bound without robustness}
\label{sec:productnorm}

A second unrestricted result addresses arbitrary learner width, but with an
exponential dependence on the energy needed by a fixed target circuit.
Its depth dependence is polynomial. The following Rademacher argument
accounts for the fact that tanh is not positively homogeneous.

The normalized-activation and exponential-moment method is closely related
to Nakhleh--Nowak's treatment of deep neural variation spaces
\cite{nakhleh2026}, which explicitly uses $\tanh(st)/s$ and slope
discretization, and to the exponential peeling argument of
Golowich--Rakhlin--Shamir~\cite{golowich2018}. We give a self-contained
product-Frobenius version for the present Gaussian-energy conventions;
we do not claim the underlying proof technique is new.

\begin{theorem}[Product-Frobenius bound for tanh]\label{thm:tanh-product}
Let $\mathcal F_{L,P}$ contain all scalar-output tanh networks with $L-1$
hidden layers, arbitrary finite widths, a linear readout, and product of
layer Frobenius norms at most $P$. Biases are absent except through an
augmented input. On every sample $x_1,\ldots,x_n$ with $\|x_i\|\le R$,
\[
 \widehat{\mathfrak R}_n(\mathcal F_{L,P})
 \le \frac{PR}{\sqrt n}
 \left[2+\sqrt{2(L-1)\log\bigl(3+n(L-1)^2\bigr)}\right].
\]
The bound is independent of every hidden width and of the individual layer
norms apart from their product.
\end{theorem}

\begin{proof}
If $R=0$, every sampled input and every network output on the sample is
zero, so the bound is immediate. Assume $R>0$ below.
Define $\phi_s(u)=\tanh(su)/s$ for $s>0$, $\phi_0(u)=u$, and
$\phi_\infty(u)=0$. These functions are odd, $1$-Lipschitz, and zero at zero.
For finite $s,t\ge0$ and $r\ge0$,
\[
 \phi_s(ru)=r\phi_{sr}(u),\qquad
 |\phi_s(u)-\phi_t(u)|\le u^2|s-t|,
 \qquad |\phi_s(u)|\le1/s\quad(s>0).
\]
The derivative bound follows by differentiating
$\phi_s(u)=u\int_0^1\operatorname{sech}^2(suv)\,dv$.

Let $\mathcal G_0$ contain the identity feature map. Inductively allow
arbitrary finite-dimensional features
\[
 g_j(x)=\phi_{s_j}(a_j^T h(x)),\quad
 h\in\mathcal G_{\ell-1},\quad \sum_j\|a_j\|^2\le1,
 \quad s_j\in[0,\infty].
\]
All these feature vectors have norm at most $R$ on the sample. Put
$B_\ell(\epsilon)=\sup_{g\in\mathcal G_\ell}
 \|\sum_i\epsilon_i g(x_i)\|$ for independent Rademacher signs.
Writing $a_j=r_j u_j$ and using the scaling identity shows
\[
 B_\ell\le
 \sup_{h,\,\|u\|\le1,\,s\in[0,\infty]}
 \left|\sum_i\epsilon_i\phi_s(u^T h(x_i))\right|.
\]
For any $\eta>0$, uniformly on $[-R,R]$ the slope family admits an
$\eta/n$ approximation using at most
$J=3+n^2R^2/\eta^2$ choices: truncate $s>n/\eta$ to the zero function,
and round slopes below that threshold down at mesh $\eta/(nR^2)$.
The error in the Rademacher sum is at most $\eta$.

We use the convex contraction principle
\[
 \mathbb E G\!\left(\sup_{t\in T}\sum_i\epsilon_i\psi_i(t_i)\right)
 \le
 \mathbb E G\!\left(\sup_{t\in T}\sum_i\epsilon_i t_i\right)
\]
for convex nondecreasing $G$ and $1$-Lipschitz $\psi_i$ vanishing at zero.
For completeness, the single-coordinate proof compares
$X_\pm=\sup_t(a_t\pm\psi(t_k))$ and
$Y_\pm=\sup_t(a_t\pm t_k)$. One has
$\max(Y_+,Y_-)\ge\max(X_+,X_-)$ by $|\psi(t)|\le|t|$.
An approximately maximizing pair for $X_+,X_-$ also gives
$Y_++Y_-\ge X_++X_-$, by the Lipschitz property.
Convex monotonicity therefore gives
$G(Y_+)+G(Y_-)\ge G(X_+)+G(X_-)$; iterate over coordinates.
Bounded separable classes follow by approximation from finite classes.
In our application, $u$ ranges symmetrically and $\phi_s$ is odd, so the
absolute supremum is the same as the one-sided supremum.

Apply this principle with $G(t)=e^{\lambda t}$ to each grid slope, then
bound the exponential of a maximum by the sum of the exponentials:
\[
 \mathbb E e^{\lambda B_\ell}
 \le e^{\lambda\eta}J\,\mathbb E e^{\lambda B_{\ell-1}}.
\]
Furthermore, $\mathbb E B_0\le R\sqrt n$, and bounded differences gives
$\mathbb E e^{\lambda B_0}
 \le e^{\lambda R\sqrt n+\lambda^2nR^2/2}$.
Writing $H=L-1$, iteration and Jensen yield
\[
 \mathbb E B_H\le H\eta+R\sqrt n+
             \frac{H\log J}{\lambda}+\frac{\lambda nR^2}{2}.
\]
Take $\eta=R\sqrt n/H$ and optimize $\lambda$ to obtain
$\mathbb E B_H\le R\sqrt n[2+\sqrt{2H\log(3+nH^2)}]$.

Finally, for a particular original network, write
$M_\ell=\|W^\ell\|_F$ and $S_\ell=\prod_{j\le\ell}M_j$.
Dividing its $\ell$th hidden vector by $S_\ell$ changes its recursion to
$\phi_{S_\ell}((W^\ell/M_\ell)g^{\ell-1})$.
Thus its normalized hidden features belong to the enlarged classes above,
and its scalar output is at most $P$ times a norm-one linear readout of
$\mathcal G_H$. Zero-norm layers give the zero function and require no
normalization. This proves the theorem. All suprema are measurable by
countable approximation of finite-dimensional parameter spaces and then
union over widths.
\end{proof}

For the Gaussian energy in our model, the product of the $L$ weighted
squared layer norms is $dP^2$. Consequently
\[
 P\le d^{-1/2}(\kappa_N/L)^{L/2}
 \le d^{-1/2}e^{B/2}\qquad\text{if }\kappa_N\le L+B.
\]
Applying the theorem to the ramp loss
$\ell(yz)=\min\{1,\max\{0,1-2yz\}\}$ gives the following explicit bound.

\begin{corollary}[All widths at bounded excess energy]\label{cor:excess-risk}
Fix $B\ge0$. With probability at least $1-\delta$, every sample-consistent
depth-$L$ network, at any finite width, with training margins at least $1/2$
and energy $\kappa_N\le L+B$ has sign error at most
\[
 \frac{4\sqrt2 e^{B/2}}{\sqrt n}
 \left[2+\sqrt{2(L-1)\log\bigl(3+n(L-1)^2\bigr)}\right]
 +\sqrt{\frac{\log(1/\delta)}{2n}}.
\]
No internal robustness assumption is imposed.
\end{corollary}

\begin{proof}
Augmented Boolean inputs have norm $\sqrt{d+1}$, so $PR\le\sqrt2e^{B/2}$.
The ramp loss is bounded by one and is $2$-Lipschitz. Symmetrization,
scalar contraction, and bounded differences give population loss at most
empirical loss plus four times the network Rademacher bound and the
stated confidence term. Empirical loss is zero and population loss
upper-bounds sign error.
\end{proof}

\begin{lemma}[A scalar transmission construction]\label{lem:scalar-transmission}
For every real $D>0$ and integer $m\ge2$, there are positive scalar weights
$\alpha,b_2,\ldots,b_{m-1},c$ such that
\[
 r_1=\tanh\alpha,\quad r_s=\tanh(b_s r_{s-1}),\quad
 cr_{m-1}=1/2,
\]
and, writing $D'=\max\{D,4\}$,
\[
 D\alpha^2+\sum_{s=2}^{m-1}b_s^2+c^2
 \le(m+12)(D'/4)^{1/m},\qquad c^2=(D'/4)^{1/m}.
\]
\end{lemma}

\begin{proof}
Set $q=(D'/4)^{1/(2m)}$, $k_s=\min\{s,m-s\}$, and
$r_s=q^s/(\sqrt{D'}k_s)$. Then $0<r_s\le1/(2qk_s)\le1/2$ and
$r_{m-1}=1/(2q)$. Choose
$\alpha=\operatorname{atanh}r_1$,
$b_s=\operatorname{atanh}(r_s)/r_{s-1}$ and $c=q$.
For $0<u\le1/2$ the power series gives
$(\operatorname{atanh}(u)/u)^2\le1+u^2$.
Writing $R_s=k_{s-1}/k_s$, the reciprocal pairs
$k/(k+1),(k+1)/k$ satisfy
\[
 \left(\frac{k+1}{k}\right)^2+
 \left(\frac{k}{k+1}\right)^2-2\le4/k^2.
\]
Consequently $\sum_{s=2}^{m-1}R_s^2\le m-2+7$,
$R_s\le2$, and $\sum_{s=1}^{m-1}k_s^{-2}<4$.
Using $q^{2m}=D'/4$ therefore gives
\[
 D\alpha^2\le q^2+1/4,\qquad
 \sum_{s=2}^{m-1}b_s^2\le q^2(m-2+7)+4,\qquad c^2=q^2.
\]
Their sum is at most $(m+12)q^2$. Empty internal sums at $m=2$
satisfy the same estimates.
\end{proof}

\begin{corollary}[Polynomial cost of additional depth in the strong-prior limit]
\label{cor:padded-strongprior}
Suppose a width-$C$, depth-$L_0$ network implements $f$ with global margin
at least $1/2$ and Gaussian energy $E$. Put
\[
 A=\max\{2E,1\},\qquad B=2E+13+7A.
\]
For each learner depth $L\ge L_0$, Corollary~\ref{cor:excess-risk}, with
$B$ replaced by $B+1$, gives a sample bound for the original unrestricted
Gaussian strong-prior learner at every width $N\ge C$.
Specifically, choose $n$ so that the displayed risk bound is at most
$\varepsilon/2$. On one sample event of probability $1-\delta$, for each
$N\ge C$ the posterior Gibbs risk is at most $\varepsilon$ for all
sufficiently small Gaussian scales. The scale threshold may depend on
$N,L,f,\varepsilon$; no uniform threshold is asserted.
The sample bound is polynomial in $L$ and exponential in $E$.
\end{corollary}

\begin{proof}
Balanced cloning into width $N\ge C$ preserves the target function and
increases energy by at most a factor of two. Write $c^2=N\|v\|^2\le2E$ for its readout
cost and $E_{\rm pre}\le2E$ for the other costs.
Apply Lemma~\ref{lem:scalar-transmission} with $D=4c^2$ and $m=t+1$.
Replace the old readout by a hidden layer whose every row is $2\alpha v^T$,
then by the lemma's scalar channels and readout. Its first new preactivation
is $2\alpha z(x)$, whose signed magnitude is at least $\alpha$.
Oddness and monotonicity preserve the sign and final margin. The new first
block has energy $4\alpha^2c^2$, and the remaining costs are exactly the
scalar ones. This gives, for every $t\ge1$, a same-width realization at
depth $L_0+t$ with energy
\[
 E_{\rm pre}+(t+13)\max\{c^2,1\}^{1/(t+1)}.
\]
Since $A^a\le1+a(A-1)$ for $0\le a\le1$ and
$(t+13)/(t+1)\le7$, the target energy is at most $L+B$.
This also covers $t=0$ by cloning directly.

For fixed $N,L$, Gaussian small-variance asymptotics imply that prior
mass on functions of infimum energy greater than $L+B+1$, divided by
prior mass on $f$, tends to zero. There are finitely many Boolean functions,
and $f$ has infimum energy at most $L+B$. This ratio bounds the posterior
mass of those functions for every consistent sample. Relaxed minimum-energy
representatives exist at fixed width, have margin at least $1/2$, and have
the same infimum as strict-margin representatives. Therefore every remaining
consistent function has a representative covered by
Corollary~\ref{cor:excess-risk}. The sample event is common to all widths,
and adding at most $\varepsilon/2$ exceptional posterior mass proves the claim.
\end{proof}

This removes exponential dependence on \emph{added learner depth} for a
fixed target witness, even at arbitrarily large learner widths. It does not
give a bound jointly polynomial in circuit size and depth, and it does not
yet cover unrestricted arbitrary-width posteriors at every nonzero scale.

\subsection{A nontrivial unrestricted family: linear threshold targets}

The preceding result is not limited to padding a fixed network. It gives a
bound polynomial in depth and input dimension for every linear threshold
whose geometric margin on the Boolean cube is inverse polynomial.

\begin{corollary}[Linear thresholds at every learner width]
\label{cor:linear-threshold}
Suppose
\[
 f(x)=\operatorname{sign}(a^T(1,x)),\qquad
 \mu=\min_{x\in\{-1,1\}^d}\frac{|a^T(1,x)|}{\|a\|}>0,
 \qquad D_* =\max\{d/\mu^2,4\}.
\]
For every depth $L\ge2$ and width $N\ge1$, the original unrestricted
Gaussian strong-prior learner admits the following common sample event.
With probability at least $1-\delta$, its Gibbs risk is at most
$\varepsilon$ for all sufficiently small Gaussian scales, provided
\[
 \frac{4 e^{13/2}}{\sqrt n}
 \sqrt{\frac{(d+1)D_*}{4d}}
 \left[2+\sqrt{2(L-1)\log\bigl(3+n(L-1)^2\bigr)}\right]
 +\sqrt{\frac{\log(1/\delta)}{2n}}
 \le\frac\varepsilon2.
\]
The sample event is simultaneous over all widths; the small-scale threshold
may depend on width and depth. In particular the sample bound is polynomial
in $d,L,\mu^{-1},\varepsilon^{-1},\log(1/\delta)$.
\end{corollary}

\begin{proof}
Use the scalar construction with first-layer energy coefficient
$D=d/\mu^2$. If its first scalar weight is $\alpha$, give every first-layer
row the vector $\alpha a/(\mu\|a\|)$. Every first preactivation has the
correct sign and magnitude at least $\alpha$. Put $b_s/N$ in every internal
entry and $c/N$ in every readout entry, where $(\alpha,b,c)$ is the positive
scalar tuple. Oddness and monotonicity preserve the target sign and the
scalar output lower bound. The energy is exactly
$D\alpha^2+\sum_s b_s^2+c^2$, at every width.
Lemma~\ref{lem:scalar-transmission} therefore gives the following upper bound on target infimum energy:
\[
 K_L=(L+12)(D_*/4)^{1/L}.
\]
Functions with infimum energy at most $K_L+1$ have representatives with
product-Frobenius norm at most
\[
 P\le\frac1{\sqrt d}
       \left(\frac{K_L+1}{L}\right)^{L/2}
 \le e^{13/2}\sqrt{\frac{D_*}{4d}}.
\]
Here $(D_*/4)^{1/L}\ge1$ and $(1+13/L)^L\le e^{13}$.
Theorem~\ref{thm:tanh-product}, the ramp-loss argument, and the
small-variance comparison with the target prove the claim exactly as in
Corollary~\ref{cor:padded-strongprior}. The elementary inequality
$\mu\le1$ follows by averaging $(a^T(1,x))^2$ over the cube; hence
$D_*\le4d/\mu^2$, making the asserted polynomial dependence explicit.
\end{proof}

For the finite-variance extension, the same construction has a
\emph{margin-one} center of energy
\[
 K_L^{(1)}\le(L+15)(D_*/4)^{1/L}.
\]
Indeed the scalar construction has readout cost $(D_*/4)^{1/L}$;
doubling that readout adds three times this cost. A posterior cutoff
$K_L^{(1)}+1$ therefore has product bound
$P\le e^8\sqrt{D_*/(4d)}$. For Theorem~\ref{thm:unrestricted-above-one},
\[
 Q_*^2=\left(\frac{K_L^{(1)}+2}{L}\right)^L
 \le e^{17}D_*/4.
\]
Thus that theorem gives a bound jointly polynomial in
$d,L,\mu^{-1},\varepsilon^{-1},\log(1/\delta)$ at \emph{every} learner width
and every exponent $\gamma\ge1+\zeta$, for each fixed $\zeta>0$.
No robustness restriction or infinite-width limit is used. The polynomial
degree may depend on $1/\zeta$; the mean-field exponent $\gamma=1$
remains outside this argument.

\section{Unrestricted finite-variance learning for \texorpdfstring{$\gamma>1$}{gamma greater than one}}
\label{sec:gaussian-distance}

There is a stronger bridge than concentration of the full Frobenius energy.
A typical Gaussian bulk has small operator norms when $\gamma>1$, even when
its total Frobenius energy is large. Gaussian isoperimetry shows that the
posterior lies near such a bulk, with the remaining signal having energy
close to a target witness's. Removing the bulk then preserves all Boolean
signs under an explicit product-norm condition. No additional restriction of prior
support is used.

Write $\sigma^2=t^2/N$, so the original Gaussian has independent entry
variances $N\sigma^2/d,\sigma^2,\sigma^2/N$. Its energy coordinates are
\[
 A_1=\sqrt{d/N}\,W^1,\qquad A_s=W^s\ (2\le s<L),
 \qquad A_L=\sqrt N\,v,
\]
and all their entries are independent $\mathcal N(0,\sigma^2)$.
Their squared Euclidean norm is exactly $\kappa_N$.

\begin{theorem}[Unrestricted learning above exponent one]
\label{thm:unrestricted-above-one}
Fix integers $d\ge1$, $L\ge2$, $C\ge1$, and a real $K\ge L$.
Use tanh hidden activations and a linear readout. Suppose that for every
$N\ge C$ there is a width-$N$, depth-$L$ target witness satisfying
\[
 f(x)z(x)\ge1\quad\hbox{for every Boolean input},\qquad \kappa_N\le K.
\]
Put
\[
 Q_* = \left(\frac{K+2}{L}\right)^{L/2}.
\]
For every fixed $\zeta>0$, every $0<\varepsilon,\delta<1$, and every
distribution $D$ on the Boolean cube, the original Gaussian learner, conditioned only
on the original global clean margin and noiseless training consistency,
has a sample bound polynomial in
\[
 d,L,C,K,Q_*,\varepsilon^{-1},\log(1/\delta)
\]
which, with probability at least $1-\delta$ over independent training
inputs drawn from $D$ and labeled by $f$, gives Gibbs error at most
$\varepsilon$, simultaneously for all
$N\ge C$ and all $0<\sigma^2\le N^{-(1+\zeta)}$.
Thus it covers all $\gamma\ge1+\zeta$ at their actual finite Gaussian scales.
The polynomial degree may depend on $\zeta$; no uniform claim as
$\zeta\downarrow0$ is made.
\end{theorem}

We prove the statement with explicit width thresholds and sample inequalities.
The hypothesis $K\ge L$ costs nothing: averaging over the cube and using
AM--GM shows that a global unit-margin witness has
$\kappa_N\ge Ld^{1/L}\ge L$.

\paragraph{A perturbation estimate using operator norms.}
Normalize hidden states by $\sqrt N$ and the augmented input by $\sqrt d$.
The input norm is at most $\sqrt2$, all coordinate activations remain
$1$-Lipschitz and zero-preserving, and the layer maps are the $A_s$ above.
If $U$ has energy at most $K+1$, and an additive weight perturbation $B$
has operator norm at most $\eta$ in every normalized layer, telescoping
one layer at a time gives
\begin{equation}\label{eq:operator-perturbation}
 \sup_x|z_{U+B}(x)-z_U(x)|
 \le\sqrt2\,L\eta
 \left[\frac{K+1+2\eta\sqrt{L(K+1)}+L\eta^2}{L-1}\right]^{(L-1)/2}.
\end{equation}
Indeed, each summand is at most $\sqrt2\eta$ times a product of $L-1$
other layer norms. The sum of squares of all perturbed layer-norm bounds
is at most the numerator, and AM--GM bounds each such product.
Set
\[
 A=\left(\frac{K+2}{L-1}\right)^{(L-1)/2},\qquad
 \eta=\min\left\{1,\frac1{4\sqrt{L(K+1)}},
                         \frac1{8\sqrt2\,LA}\right\}.
\]
The numerator in~\eqref{eq:operator-perturbation} is at most $K+2$, so the
output change is at most $1/8$. Furthermore
$A\le\sqrt e\,Q_*$, hence $\eta^{-1}$ is polynomial in $L,K,Q_*$.

\paragraph{A symmetric typical Gaussian set.}
For completely explicit constants, put
\[
 a_d^2=2(d+1)\log2+4\log8,\qquad
 b_L^2=16\log9+8\log(8L),\qquad b=\max\{a_d,b_L\}.
\]
Let $T$ consist of bulk matrices whose first and last energy-coordinate
Frobenius norms are at most $ta_d$, and whose internal operator norms are
at most $tb_L$. Then $T$ is closed, convex and centrally symmetric, and
\begin{equation}\label{eq:typical-half}
 \mu_{N,t}(T)\ge1/2.
\end{equation}
For the first block, its squared Frobenius norm is
$(t^2/N)\chi^2_{N(d+1)}$; the readout is the same expression with $d+1$
replaced by one. The exponential moment of $\chi^2$ bounds each failure
probability by $1/8$. For an internal matrix $tG/\sqrt N$, a $1/4$-net of
the unit sphere has at most $9^N$ points, and the two-net scalar Gaussian
bound gives
\[
 \Pr\{\|tG/\sqrt N\|_{\rm op}>tb_L\}
 \le2\exp(2N\log9-Nb_L^2/8)\le\frac1{4L}.
\]
Union bounding over the $L-2$ internal blocks proves~\eqref{eq:typical-half}.
Every normalized block of a member of $T$ has operator norm at most $tb$.

\paragraph{Target mass by a Gaussian shift.}
Suppose $tb\le\eta$. For a target witness $w_*$, the shifted set $w_*+T$
has target margin at least $7/8$, by~\eqref{eq:operator-perturbation}.
Central symmetry and the Gaussian change-of-mean identity yield
\begin{align*}
 \mu(w_*+T)
 &=e^{-\kappa_N(w_*)/(2\sigma^2)}
   \int_T e^{-\langle w_*,B\rangle_\kappa/\sigma^2}\,d\mu(B)\\
 &\ge e^{-K/(2\sigma^2)}\mu(T)
 \ge\tfrac12 e^{-K/(2\sigma^2)}.
\end{align*}
The inequality follows by pairing $B$ with $-B$ and using $\cosh u\ge1$.
This lower bound holds for the raw probability of the conditioning event
for every correctly labeled sample.

\paragraph{The posterior is close to the typical set.}
Let $d_\kappa(w,T)$ be Euclidean distance in energy coordinates. Gaussian
isoperimetry~\cite{borell1975} and~\eqref{eq:typical-half} imply
\[
 \mu\{d_\kappa(w,T)>r\}\le e^{-r^2/(2\sigma^2)}.
\]
Dividing by the preceding target-mass lower bound gives, for every sample,
\begin{equation}\label{eq:posterior-distance-tail}
 Q_s\{d_\kappa(w,T)^2>K+1\}\le2e^{-1/(2\sigma^2)}.
\end{equation}
On the complement, let $B$ be the metric projection of $w$ onto $T$ and
write $U=w-B$. The projection is well defined and continuous, since $T$
is nonempty, closed and convex. Then $\kappa_N(U)\le K+1$ and
$\sup_x|z_w(x)-z_U(x)|\le1/8$. The original posterior has clean margin
$>1/2$ everywhere, so $U$ has \emph{exactly the same Boolean function}
as $w$ and training margins $>3/8$. It is enough to use a margin-$1/4$
ramp loss below.

\paragraph{A polynomial width split.}
Take $\beta=\varepsilon/2$ and
\[
 M=\left\lceil\max\left\{C,2,(b/\eta)^{2/\zeta},
                                      2\log(2/\beta)\right\}\right\rceil.
\]
For every $N\ge M$ and $\sigma^2\le N^{-(1+\zeta)}$, one has
$t\le N^{-\zeta/2}$, hence $tb\le\eta$, and
$2e^{-1/(2\sigma^2)}\le\beta$.
The product norm of $U$ is at most
$d^{-1/2}[(K+1)/L]^{L/2}$. Write
\[
 G=\sqrt{(d+1)/d}\left(\frac{K+1}{L}\right)^{L/2},\qquad H=L-1.
\]
The product-norm Rademacher bound, now with the $4$-Lipschitz ramp loss
$\min\{1,\max\{0,1-4yz\}\}$, gives one sample event of probability at
least $1-\delta/2$ on which every such consistent residual network has
error at most
\begin{equation}\label{eq:distance-sample-bound}
 F_n=\frac{8G}{\sqrt n}\left[2+\sqrt{2H\log(3+nH^2)}\right]
                     +\sqrt{\frac{\log(2/\delta)}{2n}}.
\end{equation}
This event already includes all finite widths. When $F_n\le\varepsilon/2$,
adding the posterior exceptional mass $\beta$ proves the required bound
for every $N\ge M$ and every stated scale.

For $C\le N\le M$, use Theorem~\ref{thm:finitewidth-unrestricted}.
Here is a target-mass constant requiring only the present witness assumption.
Every witness weight is at most $R=\sqrt{MK}$. Put
\begin{align*}
 a&=[4(d+1+M(L-1))(MR)^{L-1}]^{-1},\\
 D_M&=\max\{1,K+[d+2+(L-2)M][2+\log(1/a)]\}.
\end{align*}
The same Gaussian-coordinate box calculation as in the finite-width proof
has physical perturbation at most $a$, retains target margin at least $3/4$,
and has mass at least $\exp(-ND_M/t^2)$.
Applying that theorem with error $\varepsilon$ and failure probability
$\delta/2$ controls all these widths and all $t\le1$ simultaneously.
Its finite-grid logarithm is polynomial in $M,d,L,\log K$.
The union of the two sample events proves the theorem.
Finally, $G\le\sqrt2 Q_*$, $b^2=O(d+\log L)$ and
$A\le\sqrt e Q_*$ show that $M$ and a sample size satisfying all the
displayed inequalities are polynomial in the stated arguments for fixed
$\zeta$. This proves the asserted quantitative dependence.

\paragraph{Scope.}
This argument does not apply at $\gamma=1$: the Gaussian bulk's operator
norm does not tend to zero with width there. It also does not make $Q_*$
polynomial for every small circuit. If a fixed base witness can be padded
with energy $K=L+B$, where $B$ is independent of the added depth, then
$Q_*\le e^{(B+2)/2}$. The theorem therefore removes exponential dependence
on \emph{additional depth}, at every width and every fixed exponent above
one, while retaining a potentially exponential dependence on base-witness
energy. Stronger target-specific energy estimates can make $Q_*$ itself
polynomial.

\section{Polynomial unrestricted learning of shallow targets at every depth}
\label{sec:shallow-unrestricted}

A shallow target can have a small first-layer Gaussian energy at a
polynomial base width simply by adding zero neurons. Its readout energy
is then polynomial. Distributing that readout cost over the learner's
remaining layers makes the product-norm parameter polynomial uniformly
in depth. This gives a family for which neither arbitrary width nor
arbitrary added depth requires a robustness restriction.

\begin{theorem}[Shallow targets, all depths and all admissible widths]
\label{thm:shallow-all-depths}
Suppose a Boolean target $f:\{-1,1\}^d\to\{-1,1\}$ has a one-hidden-layer
representation
\[
 z(x)=\sum_{j=1}^{C}v_j\tanh(a_j^T\bar x),\qquad
 \bar x=(1,x),\qquad f(x)z(x)\ge1
\]
with every coordinate of $a_j$ and every $v_j$ bounded in absolute value
by $R\ge1$. Let $C_L(f)$ denote its bounded-weight network complexity
at depth $L$: the least width $N$ admitting a realization with all
weights at most $N$ in absolute value and signed output greater than
$1/2$ everywhere.

For each fixed $\zeta>0$, the original unrestricted Gaussian learner has
sample complexity polynomial in
\[
 d,C,R,L,\varepsilon^{-1},\log(1/\delta),
\]
for every specified learner depth $L\ge2$ and every input distribution,
simultaneously for every width $N\ge C_L(f)$ and every
$\gamma\ge1+\zeta$. With probability at least
$1-\delta$ over a sample of that size, posterior Gibbs error is at most
$\varepsilon$ simultaneously at all these widths and all prior scales
$0<\sigma^2\le N^{-(1+\zeta)}$. The polynomial's degree may depend on
$\zeta$. The prior is conditioned only on the original global clean
margin and the training labels.
\end{theorem}

\begin{proof}
Let $A$ have rows $a_j^T$, and set
\[
 M=\left\lceil\max\{4,C,2R,d\|A\|_F^2\}\right\rceil,
 \qquad T=\max\{1,M\|v\|_2^2\}.
\]
Both are polynomial in $d,C,R$. Add zero hidden neurons to obtain width
$M$. This preserves the function, and its first-layer and readout
energies satisfy
\[
 E_{\rm pre}=\frac dM\|A\|_F^2\le1,
 \qquad c^2=M\|v\|_2^2\le T.
\]
In fact $M\|v\|_2^2\ge1$, since unit margin and $|\tanh|\le1$ imply
$\|v\|_2\ge1/\sqrt C$, but the maximum makes this unneeded.

For any $N\ge M$, balanced cloning of this shallow network preserves its
output exactly and gives
\[
 E_{\rm pre}\le2,\qquad c^2=N\|v_N\|_2^2\le2T.
\]
Explicitly, assign each of the $M$ original neurons either
$\lfloor N/M\rfloor$ or $\lceil N/M\rceil$ clones, copy its incoming
weights, and divide its readout weight by its number of clones.
The first cost increases by at most $M\lceil N/M\rceil/N\le2$,
and the readout cost by at most $N/(M\lfloor N/M\rfloor)\le2$.
For learner depth $L=2$ we already have a unit-margin witness with
\[
 K_2=2+2T,\qquad
 Q_*^2=\left(\frac{K_2+2}{2}\right)^2=(T+2)^2.
\]

For $L\ge3$, apply Lemma~\ref{lem:scalar-transmission} with
$D=4c^2$ and $m=L-1$. Replace the old readout by a layer whose every
row is $2\alpha v_N^T$. Follow it with the scalar transmission layers,
implemented by equal coordinates with dense matrix entries $b_s/N$,
and final readout entries $q/N$, where
\[
 q^2=\max\{c^2,1\}^{1/(L-1)}\le(2T)^{1/(L-1)}.
\]
The target had margin at least one, hence in particular at least
$1/2$, so the scalar construction gives margin at least $1/2$.
Double its final readout to restore margin at least one. This adds
$3q^2$ to the energy. Thus every width $N\ge M$ has a depth-$L$
unit-margin witness with
\[
 \kappa_N\le K_L:=2+(L+14)(2T)^{1/(L-1)}.
\]
Writing $q_*^2=(2T)^{1/(L-1)}\ge1$, we have
$K_L+2\le(L+18)q_*^2$. Consequently
\begin{equation}\label{eq:shallow-product-polynomial}
 Q_*^2=\left(\frac{K_L+2}{L}\right)^L
 \le e^{18}(2T)^{L/(L-1)}
 \le e^{18}(2T)^{3/2}\qquad(L\ge3).
\end{equation}
The exponent $L/(L-1)$ is at most $3/2$ for $L\ge3$.
Theorem~\ref{thm:unrestricted-above-one}, with base width $M$, therefore
provides a sample bound polynomial in the stated quantities for all
$N\ge M$.

We verify admissible smaller widths rather than silently identifying
$M$ with the target's minimum complexity. First, $C_L(f)\le M$ for all
$L\ge2$. At width $M$, the original incoming weights have magnitude
at most $R\le M/2$. In the added scalar channel,
$\alpha\le\operatorname{atanh}(1/2)<0.55$, so its first weights
$2\alpha v_j$ have magnitude at most $1.1R<M$.
The scalar construction has $b_s\le2.2q$ and final weights $2q/M$.
Since $T\le MC R^2\le M^4$ and $L\ge3$,
\[
 q\le(2T)^{1/4}\le2^{1/4}M,
 \qquad b_s/M<3\le M,\qquad 2q/M<3\le M.
\]
Thus this depth-$L$, width-$M$ circuit satisfies the bounded-weight
convention and has margin at least one. The case $L=2$ is immediate.

For $C_L(f)\le N<M$, start from a defining width-$C_L(f)$ circuit,
clone it to width $N$, and double its readout. It has margin greater
than one, all weights at most $2M$, and Gaussian energy bounded by a
polynomial in $d,L,M$. The Gaussian box estimate and
Theorem~\ref{thm:finitewidth-unrestricted} consequently give a sample
bound polynomial in $d,L,M,\varepsilon^{-1},\log(1/\delta)$,
simultaneously for all these widths and all $t\le1$.
Use failure probability $\delta/2$ for this event and $\delta/2$ for
the preceding event for $N\ge M$. Their intersection proves the claim.
\end{proof}

The complexity parameter controlling this theorem's sample size is the
size and weight bound of a \emph{shallow} representation. Its width
admissibility is the actual depth-$L$ complexity $C_L(f)$. The theorem
does not identify shallow complexity with the minimum complexity of an
arbitrary deep circuit, and therefore does not settle the unrestricted
problem for all small deep circuits. It also makes no claim at $\gamma=1$.

\begin{corollary}[Parity at arbitrary depth and admissible width]
\label{cor:parity-full-polynomial}
For every fixed $\zeta>0$, parity of any $k\ge1$ chosen coordinates in
$\{-1,1\}^d$ has unrestricted Gaussian posterior sample complexity
polynomial in $d,k,L,\varepsilon^{-1},\log(1/\delta)$, for every
$L\ge2$, every $N\ge C_L(f)$, and every $\gamma\ge1+\zeta$.
The bound holds for every input distribution and for the original
finite-variance posterior.
\end{corollary}

\begin{proof}
Let $q=(k-\sum_{i\in S}x_i)/2\in\{0,\ldots,k\}$ count negative
coordinates, and put $a=\log(16k)$. Define
\[
 \tau_j(x)=\tanh\bigl(a(q-j+1/2)\bigr),\qquad 1\le j\le k.
\]
Every argument has magnitude at least $a/2$, so
$|\tau_j-\operatorname{sign}(q-j+1/2)|\le2e^{-a}=1/(8k)$.
Writing $c_k=1$ for even $k$ and $c_k=0$ for odd $k$, the exact step
identity is
\[
 (-1)^q=c_k+\sum_{j=1}^k(-1)^j
                  \operatorname{sign}(q-j+1/2).
\]
Hence $F=c_k+\sum_j(-1)^j\tau_j$ satisfies
$|F-(-1)^q|\le1/8$. Represent the constant by one extra hidden neuron
with augmented-input weight $\operatorname{atanh}(1/2)$ and no other
incoming weights. A linear readout equal to $2F$ has signed margin
at least $7/4$. Its threshold-neuron readout weights are $2(-1)^j$,
and its constant-neuron weight is $4c_k$.
Its incoming weights have magnitude at most $a(k+1)$.
Thus $C=k+1$ and $R=\max\{4,a(k+1)\}$ are polynomial in $k$.
Apply Theorem~\ref{thm:shallow-all-depths}.
\end{proof}

\begin{corollary}[Symmetric Boolean targets]
The conclusion of Corollary~\ref{cor:parity-full-polynomial} holds for
any Boolean function of the Hamming weight of $k$ selected coordinates.
\end{corollary}

\begin{proof}
If its values are $h(0),\ldots,h(k)\in\{-1,1\}$, then
\[
 h(q)=\frac{h(0)+h(k)}2+
       \sum_{j=1}^k\frac{h(j)-h(j-1)}2
                    \operatorname{sign}(q-j+1/2).
\]
Each coefficient has magnitude at most one. Replacing the signs by
the same $\tau_j$ and doubling the output gives margin at least $7/4$,
using the same width and weight bounds as in the parity proof.
\end{proof}


\begin{corollary}[Polynomial-size DNF and CNF targets]
A Boolean function on $d$ bits represented by a DNF or CNF with $s\ge1$
clauses has unrestricted Gaussian posterior sample complexity polynomial
in $d,s,L,\varepsilon^{-1},\log(1/\delta)$, for every specified $L\ge2$,
every $N\ge C_L(f)$, and every $\gamma\ge1+\zeta$, with $\zeta>0$ fixed.
\end{corollary}

\begin{proof}
For a DNF, let $q_j(x)$ be the number of false literals in clause $j$,
and put $a=\log(16s)$ and
$\tau_j(x)=\tanh(a(1/2-q_j(x)))$.
The exact clause indicator is
$[1+\operatorname{sign}(1/2-q_j)]/2$.
Thus
\[
 H(x)=s-1+\sum_{j=1}^s\operatorname{sign}(1/2-q_j(x))
\]
is $-1$ when no clause is satisfied and at least $1$ otherwise.
Replacing the signs by $\tau_j$ changes $H$ by at most $1/8$.
Doubling the resulting expression gives global signed margin at least
$7/4$. It is a shallow tanh network with $s+1$ features: the extra
feature is the constant $1/2$, with readout weight $4(s-1)$.
After removing duplicate literals and discarding contradictory clauses,
each clause has at most $d$ literals. Each $q_j$ is affine in the bits,
so the remaining incoming and readout weights have magnitude at most
$R=\max\{4s,a(d+1)\}$. The shallow-target theorem applies.
A CNF is the negation of a DNF in negated literals, and negating the
readout preserves every bound. Constant functions are immediate.
\end{proof}

\subsection{Targets of any fixed intrinsic depth}
\label{sec:fixed-intrinsic-depth}

The shallow-target argument extends to any fixed target depth by using
different numbers of clones at its successive hidden layers. The degree
of the resulting polynomial may depend on the target depth, but not on
the learner's additional depth.

\begin{corollary}[Fixed target depth, arbitrary learner depth]
\label{cor:fixed-intrinsic-depth}
Fix an integer $D_0\ge2$ and $\zeta>0$. Suppose $f$ has a width-$C$,
depth-$D_0$ tanh realization in the present architecture, with signed
output at least one on the Boolean cube and all weights bounded by
$R\ge1$. For every specified learner depth $L\ge D_0$, the original
unrestricted Gaussian learner has sample complexity polynomial in
\[
 d,C,R,L,\varepsilon^{-1},\log(1/\delta).
\]
The degree and constants may depend on $D_0$ and $\zeta$.
For every input distribution, one sample event of probability at least
$1-\delta$ gives Gibbs error at most $\varepsilon$ simultaneously for
every $N\ge C_L(f)$ and every
$0<\sigma^2\le N^{-(1+\zeta)}$.
Here $C_L(f)$ is the actual bounded-weight network complexity defined in
Theorem~\ref{thm:shallow-all-depths}. The prior retains the original
global clean-margin conditioning and sign-only training likelihood.
\end{corollary}

\begin{proof}
Put $H=D_0-1$. Pad the target's hidden layers to width $C$ with zero
neurons if necessary, and write $W^1,\ldots,W^H,v$ for its weights.
Choose an integer
\[
 B=\left\lceil\max\left\{4,2R,\frac dC\|W^1\|_F^2,
                 \|W^2\|_F^2,\ldots,\|W^H\|_F^2\right\}\right\rceil,
 \qquad M=CB^H,
\]
omitting the internal-layer terms when $H=1$.
Since $\|W^1\|_F^2\le C(d+1)R^2$ and
$\|W^s\|_F^2\le C^2R^2$ internally, $B$ is polynomial in $d,C,R$,
and so is $M$ for fixed $D_0$.

At hidden layer $s$, give each original neuron exactly
\[
 n_s=B^{H-s}
\]
active copies, and fill the other positions up to width $M$ with zero
neurons. Copy first-layer rows. For $s\ge2$, give every edge from a copy
of source $j$ to a copy of target $i$ weight $W^s_{ij}/n_{s-1}$.
All unused edges are zero. Since $n_H=1$, retain the original readout
on the final active neurons. This preserves the function exactly.
The first-layer and internal energy terms are respectively
\[
 \frac dM n_1\|W^1\|_F^2
      =\frac d{CB}\|W^1\|_F^2\le1,
 \qquad
 \frac{n_s}{n_{s-1}}\|W^s\|_F^2
      =\frac1B\|W^s\|_F^2\le1.
\]
Thus the energy before the readout is at most $H$. Put
\[
 T=\max\{1,M\|v\|_2^2\}\le MC R^2.
\]
It is polynomial in $d,C,R$ for fixed $D_0$.

Balanced cloning of this width-$M$ network into any $N\ge M$ preserves
the function and gives
\[
 E_{\rm pre}\le2H,\qquad c^2=N\|v_N\|_2^2\le2T.
\]
If $L=D_0$, these are already unit-margin witnesses with
$K=2H+2T$, and
\[
 Q_*^2=\left(\frac{K+2}{D_0}\right)^{D_0}
       \le(4T)^{D_0}.
\]

If $L>D_0$, set $m=L-H\ge2$. Replace the old readout by the scalar
transmission from Lemma~\ref{lem:scalar-transmission}, using $D=4c^2$,
and double the final scalar readout to obtain unit margin. Its readout
cost before doubling is
$q^2=\max\{c^2,1\}^{1/m}\le(2T)^{1/m}$.
The total energy is at most
\[
 K_L=2H+(m+15)(2T)^{1/m}.
\]
This is the desired depth: the $H$ old hidden layers and $m-1$ new hidden
layers total $L-1$. Since $(2T)^{1/m}\ge1$,
\begin{align*}
 Q_*^2
 &=\left(\frac{K_L+2}{L}\right)^L\\
 &\le\left(1+\frac{H+17}{L}\right)^L(2T)^{L/(L-H)}\\
 &\le e^{H+17}(2T)^{(D_0+1)/2}.
\end{align*}
The last step uses $L/(L-H)=1+H/m\le1+H/2$.
For $L\ge2D_0$, the sharper upper bound
$Q_*^2\le e^{H+17}(2T)^2$ also holds.
Both $K_L$ and $Q_*$ are polynomial in the stated parameters when
$D_0$ is fixed, uniformly over the learner depth.
Theorem~\ref{thm:unrestricted-above-one} now controls every $N\ge M$.

It remains to check admissible widths below $M$.
First $C_L(f)\le M$ at every $L\ge D_0$.
The tapered construction has all original weights at most $R\le M/2$.
In a padded scalar channel, $\alpha\le\operatorname{atanh}(1/2)<0.55$,
so its first added weights $2\alpha v_j$ have magnitude less than $M$.
The scalar construction has $b_s<2.2q$, and
\[
 T\le MC R^2\le M^4,\qquad
 q\le(2T)^{1/4}\le2^{1/4}M.
\]
Its dense internal entries $b_s/M$ and final readout entries $2q/M$
are therefore less than $3\le M$.
This gives a bounded-weight width-$M$ witness with unit margin at each
learner depth; for $L=D_0$ no scalar channel is needed.

For $C_L(f)\le N<M$, clone a defining width-$C_L(f)$ realization to
width $N$ and double its readout. Its margin is greater than one and
all weights are at most $2M$. The Gaussian box estimate and
Theorem~\ref{thm:finitewidth-unrestricted} give a bound polynomial in
$d,L,M,\varepsilon^{-1},\log(1/\delta)$ for all these widths and all
$t\le1$. Allocate failure probability $\delta/2$ to this event and to
the preceding event for $N\ge M$. Their intersection proves the claim.
\end{proof}

Thus polynomial dependence on learner depth holds for targets of any
fixed intrinsic depth with polynomial target-size dependence. The width
$M=CB^{D_0-1}$ and the exponents above explain why this does not give a
single polynomial jointly in the intrinsic target depth and target size.
The mean-field exponent $\gamma=1$ is also still outside this result.

\begin{corollary}[Quantitative dependence on intrinsic depth]
\label{cor:intrinsic-depth-quantitative}
Fix $\zeta>0$. There are constants $a_\zeta,b_\zeta>0$, independent of
the intrinsic depth, with the following property. Suppose $f$ has a
width-$C$, depth-$D_0$ tanh realization with unit signed margin and
weights bounded by $R\ge1$, where $D_0\ge2$, and set
$S=2+d+C+R$. For every specified learner depth $L\ge D_0$, the sample
bound
\[
 n\ge\left\lceil S^{a_\zeta D_0}
       \bigl(L+\varepsilon^{-1}+\log(2/\delta)\bigr)^{b_\zeta}
       \right\rceil
\]
suffices for posterior Gibbs error at most $\varepsilon$, with probability
at least $1-\delta$ for every input distribution, simultaneously for all
$N\ge C_L(f)$ and $0<\sigma^2\le N^{-(1+\zeta)}$.
The prior and likelihood are those of
Corollary~\ref{cor:fixed-intrinsic-depth}.

Thus the dependence on intrinsic depth is
$\exp[O_\zeta(D_0\log S)]$, while the dependence on learner depth remains
polynomial with degree independent of $D_0$. In particular, if
$D_0\le c\log S$ for a fixed constant $c$, the target-parameter factor
is at most $\exp[O_{\zeta,c}((\log S)^2)]$.
This is a quasipolynomial bound in $S$, not a single polynomial jointly
in intrinsic depth and target size.
\end{corollary}

\begin{proof}
First suppose $L\ge2D_0$. Use the tapered construction in the preceding proof, with
$H=D_0-1$, $M=CB^H$, and $T=\max\{1,M\|v\|_2^2\}$.
Its explicit definition gives $B\le S^5$, so
$M\le S^{5D_0}$ and $T\le MC R^2\le S^{7D_0}$.
Put $m=L-H$. Since $L\ge2D_0$, one has $L/m\le2$ and
\[
 K_L=2H+(m+15)(2T)^{1/m}=O(LT),\qquad
 Q_*^2\le e^{H+17}(2T)^{L/m}
          \le e^{H+17}(2T)^2=S^{O(D_0)}.
\]

For the remaining range $D_0\le L<2D_0$, do not taper the target.
Instead pad its original realization to
$M_0=\lceil\max\{4,C,2R\}\rceil\le S^2$, and put
$T_0=\max\{1,M_0\|v\|_2^2\}\le S^5$.
Balanced cloning to every $N\ge M_0$ gives
\[
 E_{\rm pre}\le2H S^4,\qquad c^2\le2T_0.
\]
At $L=D_0$, retain the readout. For $L>D_0$, append the same scalar
transmission with $m=L-H\ge2$. These constructions give, respectively,
\[
 K_L\le2H S^4+2T_0,
 \qquad
 K_L\le2H S^4+(m+15)(2T_0)^{1/m}.
\]
In either case $K_L\le19LS^5$. Therefore
\[
 Q_*^2\le(20S^5)^L\le S^{7L}\le S^{14D_0}.
\]
The width-$M_0$ witness has weights bounded by $M_0$: the original
weights satisfy $R\le M_0/2$, and the scalar-channel check from the
preceding proof applies because $T_0\le M_0^4$.
Thus $C_L(f)\le M_0$ in this range.

In both depth ranges, the polynomial in
Theorem~\ref{thm:unrestricted-above-one} has degree and constants
depending only on $\zeta$. Substituting the corresponding bounds gives
the stated sample form above the respective base width $M$ or $M_0$.
For smaller admissible widths, the defining-witness and Gaussian-box
argument in the preceding proof gives a sample bound polynomial in
$d,L,M,\varepsilon^{-1},\log(1/\delta)$ (or with $M_0$ in place of $M$),
with degree independent of $D_0$. Substitution gives the same sample form.
Split the failure probability between the two width ranges and enlarge
$a_\zeta,b_\zeta$ to finish.
\end{proof}

\section{Limits of cloning and absolute-energy counting}
\label{sec:representation-limits}

These observations delimit two proof strategies. Neither is a lower bound
on the sample complexity of the original Gaussian learner.

\paragraph{Cloning preserves an absolute path quantity.}
For a depth-$D_0$ realization with $H=D_0-1$ hidden layers, define
\[
 A=|v|^T|W_H|\cdots|W_1|\mathbf1_{d+1},\qquad
 P_F=\|v\|_2\prod_{s=1}^H\|W_s\|_F.
\]
Absolute values are entrywise. By the operator norm bound,
\begin{equation}\label{eq:absolute-path-product-bound}
 A\le\sqrt{d+1}\,P_F.
\end{equation}
Replacing each neuron by any positive integer number of exact clones,
and dividing each outgoing weight equally among its source clones,
preserves $A$. Indeed, the multiplicity of each original path cancels
exactly against the product of the source-clone divisors on that path.
Adding zero neurons also preserves $A$.

Suppose the original output attains absolute value one, and append a
zero-bias tanh network that takes this scalar output as its input and
restores unit signed margin. Write its scalar response as $F$. Oddness
implies $F(1)\ge1$. Since $|\tanh u|\le|u|$, its own absolute path sum
$A_F$ satisfies $1\le|F(1)|\le A_F$. The appended realization has
absolute path sum $AA_F\ge A$. Thus these two operations cannot reduce
the lower bound in \eqref{eq:absolute-path-product-bound}.

For any resulting width-$N$, depth-$L$ network, its $L$ Gaussian layer
energies have product exactly $dP_F^2$. Consequently
\begin{equation}\label{eq:cloning-product-obstruction}
 \left(\frac{\kappa_N+2}{L}\right)^L
 \ge dP_F^2\ge\frac{d}{d+1}A^2.
\end{equation}
In particular, cloning followed by scalar output transmission cannot
always make the product-norm parameter polynomial in the size of the
\emph{given realization}.

A simple example makes the qualification essential. On one bit, let
$h_0(x)=x$, $h_s(x)=\tanh(2h_{s-1}(x))$ for $1\le s\le H$, and use
readout $z(x)=h_H(x)/h_H(1)$. This width-one, depth-$H+1$ network
represents $f(x)=x$ with exact unit margin. Inductively $h_H(1)>0.9$,
so all weights are at most two, whereas $A=2^H/h_H(1)\ge2^H$.
The right side of \eqref{eq:cloning-product-obstruction} is therefore
at least $4^H/2$. The same function has an inexpensive shallow
representation. This is an obstruction to the specified representation
operations, not to finding a different representation of the function.

For uniform-per-layer clone counts $n_s$ and common padded width $M$,
one has the stronger exact identity
\[
 a_1=\frac{dn_1}{M}\|W_1\|_F^2,\quad
 a_s=\frac{n_s}{n_{s-1}}\|W_s\|_F^2\ (2\le s\le H),\quad
 a_{H+1}=\frac{M}{n_H}\|v\|_2^2,\qquad
 \prod_{s=1}^{H+1}a_s=dP_F^2.
\]
If the hidden-layer energy is at most $H+B$ with $B\ge0$, AM--GM
gives $\prod_{s=1}^H a_s\le(1+B/H)^H\le e^B$, and hence
$a_{H+1}\ge dP_F^2e^{-B}$. This explains why tapering can leave an
exponentially expensive readout when intrinsic depth grows.

\begin{proposition}[An absolute energy sublevel contains every Boolean function]
\label{prop:all-functions-linear-energy}
For every integer $d\ge1$, put
\[
 L=20(d+1)+4,\qquad M=64\,2^d(d+1)^6.
\]
For every Boolean function $f$ on $d$ bits and every width $N\ge M$,
there is a depth-$L$, width-$N$ tanh network with unit signed margin
on the entire cube and Gaussian energy $\kappa_N\le3L$.
In particular, the logarithm of the number of Boolean functions in this
energy sublevel is $2^d\log2=\exp(\Omega(L))$.
\end{proposition}

\begin{proof}
Write $m=2^d$ and $b=\log(16m)$. For each cube point $a$, let
\[
 h_a(x)=\tanh\bigl(b(a\cdot x-d+1/2)\bigr),\qquad
 s_a(x)=2\mathbf1\{x=a\}-1.
\]
The preactivation has magnitude at least $b/2$, with the correct sign,
so $|h_a-s_a|\le2e^{-b}=1/(8m)$. The exact identity
\[
 f(x)=\frac12\sum_a f(a)\bigl(s_a(x)+1\bigr)
\]
therefore gives a shallow tanh realization
\[
 z(x)=\sum_a f(a)h_a(x)+\sum_a 2f(a)\,(1/2)
\]
with signed margin at least $2-1/8>1$. Implement the second sum by
$m$ constant neurons, each with first-layer bias
$\operatorname{atanh}(1/2)$ and all other incoming weights zero.
There are $C=2m$ hidden neurons and $\|v\|_2^2=5m$.
All weights have magnitude at most $R=4(d+1)^2$.

Our chosen $M=32C(d+1)^6$ satisfies
$M\ge\max\{4,C,2R,d\|W_1\|_F^2\}$, since
$d\|W_1\|_F^2\le dC(d+1)R^2\le16C(d+1)^6$.
Thus the zero-padding and balanced-cloning construction in
Theorem~\ref{thm:shallow-all-depths} applies at every $N\ge M$ with
\[
 T=M\|v\|_2^2=320m^2(d+1)^6,
 \qquad \kappa_N\le2+(L+14)(2T)^{1/(L-1)}.
\]
Since $d+1\le2^d$ and $640<2^{10}$,
$2T\le2^{8d+10}\le2^{L-1}$. Therefore
$\kappa_N\le2L+30\le3L$.
\end{proof}

This proposition concerns representations and energy, independently of
the prior variance exponent. It rules out a bound polynomial in $d,L,K$
for the logarithm of the number of all Boolean functions with energy
at most $K=3L$. The exhibited shallow widths are exponential in $d$;
the proposition does not say these are all small-circuit targets.
It also does not compare competing posterior masses relative to a
particular simple target, and gives no sample-complexity lower bound.
The product-norm parameter can still be exponential when $K=3L$.

\section{Two partial results at the critical exponent}
\label{sec:critical-padding-new}

The preceding unrestricted learning theorem stops at $\gamma>1$.
At $\gamma=1$, polynomial-depth target-mass estimates alone are already
available from robust neighborhood constructions; they do not settle the
posterior's generalization. The next estimate supplies a particularly
simple stochastic padding construction whose additional depth costs only
a universal constant per layer. It also identifies precisely which part
of the argument is still missing.

\begin{proposition}[Finite-width critical target mass by stochastic padding]
\label{prop:critical-stochastic-padding}
Let $f$ have the shallow unit-margin representation from
Theorem~\ref{thm:shallow-all-depths}, with $C$ features and all weights
bounded by $R\ge1$. Set
\[
 \rho=(8\sqrt2 CR)^{-1},\quad
 B=\frac12\left[d\max_j\|a_j\|^2+
                (d+1)(\rho^2-1-\log\rho^2)\right],
 \qquad K=B+8C^2R^2+18L+2.
\]
At the exact critical prior scale $\gamma=1$, for every $L\ge2$ and
\[
 N\ge512C^2R^2\bigl[(d+1)\log2+\log(8L)\bigr],
\]
the raw Gaussian probability of the target's global clean-margin event is
\[
 \mu_{N,1}(A_{N,f})\ge\tfrac14e^{-2NK}.
\]
No support restriction is imposed on the original prior.
\end{proposition}

\begin{proof}
We construct an auxiliary law $\mathbb Q$ on network weights and show
that it has relative entropy at most $NK$ from the original Gaussian
law, while assigning the target event probability at least $1/2$.
Introduce independent uniform labels $J_i\in\{1,\ldots,C\}$, also under
the reference law; forgetting these labels can only decrease entropy.
Under $\mathbb Q$, the first-layer row is
\[
 W_i^1=a_{J_i}+\rho G_i/\sqrt d,
 \qquad G_i\sim\mathcal N(0,I_{d+1}).
\]
Its conditional entropy from the reference $\mathcal N(0,I_{d+1}/d)$
is at most $B$, so the entire first layer costs at most $NB$.
Let $h_i(x)=\tanh(W_i^1\bar x)$ and
\[
 u_i=Cv_{J_i}/N,\qquad z_N(x)=\sum_i u_i h_i(x).
\]
The expected value of $z_N$ differs from the original shallow output
by at most $CR\sqrt2\rho=1/8$. Its summands before division by $N$
are bounded in absolute value by $CR$. Hoeffding's inequality and a
union bound over the cube imply, except with probability at most $1/8$,
\[
 f(x)z_N(x)\ge3/4\qquad\hbox{for every }x.
\]
Indeed the probability that the uniform approximation to its expectation
exceeds $1/8$ is at most $2^{d+1}e^{-N/(128C^2R^2)}$.

For $L\ge3$, give each row of the next weight matrix the mean $4u$
and retain its reference covariance $I/N$. The conditional preactivation
mean is $4z_N(x)$, and its variance is at most one.
If a Gaussian has signed mean at least $3$ and variance at most one,
its expected signed tanh is greater than $2/3$. For example, its signed
value is at least $1$ except on an event of probability at most
$\Phi(-2)<1/32$, and $\tanh1>3/4$, giving the lower bound $89/128>2/3$.
The layer's conditional entropy cost is
\[
 \frac{N^2}{2}\|4u\|^2\le8NC^2R^2.
\]
Its empirical mean activation therefore has signed value at least
$7/12$ at every input, except on an event of conditional probability
at most $2^d e^{-N/288}$.

For every remaining hidden layer, give each incoming weight mean $6/N$
and retain variance $1/N$. Its preactivation mean is six times the
preceding empirical mean. Thus the same induction preserves the signed
empirical mean bound $7/12$: its next preactivation mean is at least
$7/2>3$, and the same concentration estimate applies.
Each such layer costs exactly $18N$ in conditional relative entropy.
The union of the bad events over all hidden layers has probability
less than $1/8$ under the displayed width assumption.

Finally, give each readout weight mean $2/N$ and retain variance $1/N^2$.
This costs $2N$. Conditional on the hidden layers, its mean has signed
value at least $7/6$ and its noise variance is at most $1/N$ at each
input. A Gaussian tail bound and the union bound show that its noise
is uniformly at most $1/2$, except with probability at most
$2^{d+1}e^{-N/8}<1/8$. Consequently the output has signed margin at
least $2/3$, strictly above the required $1/2$.
When $L=2$, omit the extra hidden layers and instead give the readout
mean $2u$, at entropy cost at most $2NC^2R^2$. The same final Gaussian
noise estimate applies, and its mean has signed margin at least $3/2$.

The relative entropy chain rule gives
$D(\mathbb Q\Vert\mu_{N,1})\le NK$, and the preceding good events have
joint probability at least $1/2$. Binary data processing yields, for
$p=\mu_{N,1}(A_{N,f})$ and $q=\mathbb Q(A_{N,f})\ge1/2$,
\[
 NK\ge q\log(1/p)-\log2\ge\tfrac12\log(1/p)-\log2.
\]
Rearranging proves the stated target-mass lower bound.
\end{proof}

The estimate is deliberately elementary. In particular, it proves
neither polynomial counting of all low critical-rate Boolean functions
nor polynomial sample complexity at arbitrary learner width. The
existing critical-rate compression argument propagates tilt bounds by
$g\mapsto3g+2g^2$; a linear-in-depth target rate does not prevent that
recursion from growing very rapidly. The unresolved issue is control of
\emph{competing posterior functions}, rather than construction of a
likely enough target.

\begin{lemma}[Polynomial attenuation in the unperturbed critical bulk]
\label{lem:critical-bulk-attenuation}
Fix $q_0>0$, and define
\[
 q_{s+1}=\mathbb E\tanh^2(\sqrt{q_s}Z),\qquad
 a_s=\mathbb E\operatorname{sech}^2(\sqrt{q_s}Z),
 \qquad Z\sim\mathcal N(0,1).
\]
Then, as $s\to\infty$,
\[
 q_s=\frac1{2s}+O\!\left(\frac{\log s}{s^2}\right),\qquad
 \prod_{s=0}^{L-1}a_s=c(q_0)L^{-1/2}(1+o(1)),
\]
where $0<c(q_0)<\infty$. Moreover
\[
 \sum_{s=0}^{L-1}(a_s^{-2}-1)=\log L+O(1).
\]
The constants may depend on $q_0$.
\end{lemma}

\begin{proof}
The strict inequality $\tanh^2 u<u^2$ for $u\ne0$ shows that $q_s$
decreases to zero; continuity excludes a positive limit.
Taylor expansion, with globally bounded derivatives and Gaussian
moments controlling the remainder, gives
\[
 F(q):=\mathbb E\tanh^2(\sqrt qZ)=q-2q^2+O(q^3).
\]
For all sufficiently large $s$ this implies
$1/q_{s+1}-1/q_s=2+O(q_s)$, first giving $q_s=\Theta(1/s)$ and then,
by summation, $1/q_s=2s+O(\log s)$. This proves the variance estimate.
The identity $a_s=1-q_{s+1}$ is exact. Hence
\[
 \log a_s=-q_{s+1}+O(q_{s+1}^2),\qquad
 a_s^{-2}-1=2q_{s+1}+O(q_{s+1}^2).
\]
Since $q_s-1/(2s)$ and $q_s^2$ are summable, summing the first identity
shows that $\log\prod_{s<L}a_s+(1/2)\log L$ converges to a finite real
constant. Exponentiating gives the positive constant $c(q_0)$.
Summing the second identity proves the final assertion.
\end{proof}

The factor $a_s$ is the derivative at zero of
$m\mapsto\mathbb E\tanh(m+\sqrt{q_s}Z)$. Thus a weak coherent signal
in the \emph{unperturbed} critical bulk experiences polynomial, rather
than exponential, total attenuation. Compensating each infinitesimal
attenuation with a gain $a_s^{-1}$ would add only $\log L+O(1)$ to the
sum of squared scalar gains. This motivates a possible
$L/2+O(\log L)$ entropy padding cost. It is not a proof of such a
posterior result: a finite coherent signal changes the covariance
recursion, and the conditioned posterior need not retain the
unperturbed Gaussian feature laws.

\section{An obstruction to a uniform critical-rate counting theorem}
\label{sec:critical-rate-obstruction}

There is a stronger reason to avoid treating the critical exponent as
only a technical extension of the preceding proof. Absolute critical-rate
sublevel sets cannot have a counting bound jointly polynomial in input
dimension, depth, and rate. This is a statement about the actual
mean-field rate functional, although it is still not a sample-complexity
lower bound for a simple target.

We use the critical finite-domain variational formula proved in
Appendix~\ref{sec:critical-variational}.
For a feature law $\nu$ on $[-1,1]^{\mathcal X}$, let
$Q(\nu)=\mathbb E_\nu hh^T$, and let $\mathsf M_Q$ be the law of
coordinatewise $\tanh(G)$ for $G\sim\mathcal N(0,Q)$.
The critical rate of a global-margin Boolean function is
\[
 J_L(f)=\inf\left\{
 D(\nu_1\Vert\mathsf M_{Q_0})+
 \sum_{s=2}^{L-1}D(\nu_s\Vert\mathsf M_{Q(\nu_{s-1})})+
 \frac12z^TQ(\nu_{L-1})^+z:
 f(x)z(x)\ge\frac12\ \forall x\right\},
\]
where the final term is infinite outside the covariance range and
$Q_0(x,y)=\bar x\cdot\bar y/d$.

\begin{theorem}[Every Boolean function has polynomial critical rate at a polynomial depth]
\label{thm:critical-counting-obstruction}
For each $d\ge2$, set
\[
 r=20(d+1)^5,\qquad L=r+4.
\]
Every Boolean function on $\{-1,1\}^d$ satisfies
\[
 J_L(f)\le100000L.
\]
Consequently no polynomial $P$ can satisfy
\[
 \log\#\{f:J_L(f)\le K\}\le P(d,L,K)
\]
for every $d,L,K$.
\end{theorem}

\begin{proof}
Write $m=2^d$. First we obtain a crude quantitative lower bound on the
smallest eigenvalue of the unperturbed first hidden covariance
\[
 Q=\mathbb E_w h(w)h(w)^T,\qquad
 h(w)_x=\tanh(w\cdot\bar x),\qquad
 w\sim\mathcal N(0,I_{d+1}/d).
\]
For each $a\in\{-1,1\}^d$, choose the center
\[
 w_a=\bigl(-b(d-1/2),ba_1,\ldots,ba_d\bigr),
 \qquad b=\log(16m).
\]
At $x=a$ its preactivation is $b/2$, whereas at every other cube point
it is at most $-3b/2$. Thus its feature vector is within $1/(8m)$
coordinatewise of the vector that is $+1$ at $a$ and $-1$ elsewhere.
Let $E_a$ be the coordinate box of radius
$\rho=1/[8m(d+1)]$ around $w_a$. These boxes are disjoint, and changing
weights within a box changes every activation by at most $1/(8m)$.
The conditional mean vectors
$\bar h_a=\mathbb E[h(w)\mid w\in E_a]$ therefore form a matrix
$H$ within operator norm $1/4$ of $2I-\boldsymbol1\boldsymbol1^T$.
Since $m\ge4$, that matrix has smallest singular value at least two;
hence $HH^T\succeq I$.

Every box has Gaussian probability at least
\[
 (2\rho)^{d+1}\left(\frac d{2\pi}\right)^{(d+1)/2}
 \exp\left[-\frac d2
       (\|w_a\|+\sqrt{d+1}\rho)^2\right]
 \ge p:=e^{-20(d+1)^5}.
\]
For the last crude estimate, use
$b\le3(d+1)$, $\|w_a\|+\sqrt{d+1}\rho\le4(d+1)^2$,
$(d+1)\log(1/(2\rho))\le3(d+1)^2$, and
$(d/(2\pi))^{(d+1)/2}\ge e^{-(d+1)}$.
Conditional Jensen's inequality, together with disjointness, gives
\[
 Q\succeq\sum_a\Pr(E_a)\bar h_a\bar h_a^T\succeq pI.
\]

Fix any Boolean vector $f\in\{-1,1\}^m$ and put
$\eta=\sqrt{p/m}\le1$. Take the first feature law to be its
unperturbed law, at zero entropy cost. For the second feature law,
shift its Gaussian preactivation mean to $\eta f$, retaining covariance
$Q$. Gaussian relative entropy and data processing give cost at most
\[
 \frac12\eta^2 f^TQ^{-1}f\le\frac12.
\]
All diagonal entries of $Q$ coincide, so the new feature mean is
exactly $m_0f$ for some $m_0>0$. It satisfies $m_0\ge\eta/32$:
for $0\le t\le1$ and variance $q\le1$,
\[
 \mathbb E\operatorname{sech}^2(t+\sqrt qZ)
 \ge\Pr\{|Z|\le1\}\operatorname{sech}^2(2)>1/32,
\]
and integration from zero to $\eta$ proves the claim.

For any feature law with mean vector $u$, its second moment obeys
$Q\succeq uu^T$, hence $u$ lies in its range and $u^TQ^+u\le1$.
Therefore shifting the next Gaussian preactivation mean to $g u$
costs at most $g^2/2$ in entropy. All feature variances remain at most
one. In our construction all diagonal second moments remain identical,
so mean vectors remain of the form $m_sf$.
Choose $g=100$ for $r$ successive layers. If $0<m_s\le1/100$, the
preceding derivative estimate gives
\[
 m_{s+1}\ge100m_s/32>3m_s.
\]
If $m_s\ge1/100$, monotonicity gives $m_{s+1}>1/32>1/100$.
Since
\[
 \log(1/\eta)=10(d+1)^5+(d/2)\log2<11(d+1)^5,
\]
the choice $r=20(d+1)^5$ guarantees $m_r\ge1/100$.
Each amplification layer costs at most $5000$.

Use one last mean shift with $g=400$. Its mean preactivation has
signed magnitude at least four. With Gaussian variance at most one,
the expected signed tanh is greater than $2/3$, by the Gaussian tail
estimate used in Proposition~\ref{prop:critical-stochastic-padding}.
This layer costs at most $80000$.
Take the final output $z$ equal to this last feature mean. It has
signed margin greater than $2/3$, and its Gaussian readout cost is
at most $1/2$, again by $Q\succeq zz^T$.
The number of hidden layers is $1+1+r+1=r+3$, so the depth is
$L=r+4$. The total rate is at most
\[
 \tfrac12+5000r+80000+\tfrac12\le100000L.
\]
Finally, at these polynomial values of $L$ and $K=100000L$, the rate
sublevel contains every one of the $2^{2^d}$ Boolean functions. Its
logarithmic cardinality is $2^d\log2$, which exceeds every polynomial
in $d,L,K$ for sufficiently large $d$.
\end{proof}

This does not prove that a simple target is hard for the critical
posterior. Such a lower bound would require comparison with the target's
\emph{own optimal rate}, sufficiently accurately to control which
functions win after conditioning. The theorem instead rules out a
specific proposed upper-bound route: a universal counting theorem based
only on absolute critical rate, with polynomial dependence on depth.
It is compatible with more precise comparison to the target, or with
additional structural information about posterior minimizers.


\begin{proposition}[The obstruction also holds directly at finite width]
\label{prop:critical-finitewidth-capacity}
With $d,L$ as in Theorem~\ref{thm:critical-counting-obstruction}, for
all widths $N\ge\lceil e^{100(d+1)^5}\rceil$ and every Boolean $f$,
\[
 \mu_{N,1}(A_{N,f})\ge\tfrac14
                  \exp\{-200000NL\}.
\]
This is a direct statement about the original finite-width Gaussian
prior and requires no appeal to a large-deviation limit.
\end{proposition}

\begin{proof}
Retain $m=2^d$, $p=e^{-20(d+1)^5}$ and $r=20(d+1)^5$, and now put
$\eta=\sqrt{p/(4m)}$. Sample the first layer from its original law.
Let $H$ be its $N\times m$ matrix of activations and
$\widehat Q=H^TH/N$. The preceding covariance bound and entrywise
Hoeffding imply
\[
 \Pr\{\lambda_{\min}(\widehat Q)<p/2\}
 \le2m^2e^{-Np^2/(8m^2)}\le1/8.
\]
Indeed entrywise error at most $p/(2m)$ implies operator error at most
$p/2$. On the good event, set
\[
 u=\frac\eta N H\widehat Q^{-1}f;
\]
on its complement set $u=0$. Give every row of the next matrix this
mean vector $u$, preserving reference covariance $I/N$.
On the good event,
\[
 H^Tu=\eta f,\qquad
 \|u\|^2=\frac{\eta^2}{N}f^T\widehat Q^{-1}f
             \le\frac1{2N}.
\]
This changed layer has conditional relative entropy at most $N/4$.
Every conditional variance is at most one, so its signed expected
activation is at least $\eta/32$ on the good event.

At each of the next $r$ layers give every incoming weight mean $100/N$,
and at one last hidden layer use mean $400/N$, always retaining
variance $1/N$. Simultaneously over all cube inputs and all these
$r+2$ shifted hidden layers, conditional Hoeffding gives absolute
empirical-mean error at most $\eta/128$, with total failure probability
at most
\[
 2m(r+2)e^{-N\eta^2/32768}\le1/8.
\]
On this event, the first shifted empirical mean has signed value at
least $\eta/64$ everywhere. As long as its minimum signed value is
at most $1/100$, a gain-$100$ layer increases it by at least a factor
of two: its expected increase is at least a factor of three, and
$\eta/128$ is at most half the current minimum. Once the minimum
reaches $1/100$ it stays above that level, by monotonicity and the
bound $1/32-\eta/128>1/100$.
The choice of $r$ ensures $2^r\eta/64\ge1/100$.
The final gain-$400$ layer then has signed expected mean greater than
$2/3$ and signed empirical mean greater than $5/8$.

Give the final readout mean $2/N$ per entry and preserve variance
$1/N^2$. Its signed mean exceeds $5/4$, and its Gaussian noise is
uniformly at most $1/2$ except with probability
$2m e^{-N/8}\le1/8$. Its clean margin is consequently greater than
$3/4$. The proposed weight law assigns $A_{N,f}$ probability at least
$5/8$, hence at least $1/2$.

The relative entropy chain rule bounds this law's entropy from the
original Gaussian by
\[
 N/4+5000Nr+80000N+2N\le100000NL.
\]
Binary data processing as in
Proposition~\ref{prop:critical-stochastic-padding} proves the mass bound.
For completeness, the displayed width assumption dominates each of
\[
 \frac{8m^2}{p^2}\log(16m^2),\qquad
 \frac{32768}{\eta^2}\log(16m(r+2)),\qquad 8\log(16m),
\]
which are the three concentration thresholds used above. This follows
directly by substituting $m=2^d$, $p=e^{-20(d+1)^5}$ and $d\ge2$.
All adaptive means are measurable functions of preceding weights;
the inverse is used only on the event with smallest eigenvalue at
least $p/2$.
\end{proof}

\subsection{An exponential lower bound for the first tanh kernel eigenvalue}
\label{sec:critical-first-spectrum}

The preceding box argument can be sharpened considerably by exploiting
the positive Hermite expansion of the tanh kernel.

\begin{lemma}[A quantitative first-layer spectrum bound]
\label{lem:first-tanh-spectrum}
Let $d\ge1$, $D=d+1$, $m=2^d$, and let $Q$ be the $m\times m$ matrix
\[
 Q(x,y)=\mathbb E_w\tanh(w\cdot\bar x)\tanh(w\cdot\bar y),
 \qquad w\sim\mathcal N(0,I_{d+1}/d).
\]
Then
\[
 \lambda_{\min}(Q)\ge e^{-13D}.
\]
\end{lemma}

\begin{proof}
Write $q=D/d\in[1,2]$, and let $h_n$ be the coefficient of the
normalized probabilists' Hermite polynomial $\operatorname{He}_n/\sqrt{n!}$
in $z\mapsto\tanh(\sqrt q z)$. The coefficients vanish for even $n$.
For every odd $n\ge1$, the Fourier sine identity
\[
 \tanh(\sqrt q z)=\frac1{\sqrt q}\int_0^\infty
       \frac{\sin(tz)}{\sinh(\pi t/(2\sqrt q))}\,dt
\]
and $\mathbb E[\operatorname{He}_n(Z)e^{itZ}]=(it)^ne^{-t^2/2}$ give
\[
 |h_n|=\frac1{\sqrt{q n!}}\int_0^\infty
       \frac{t^ne^{-t^2/2}}{\sinh(\pi t/(2\sqrt q))}\,dt.
\]
The sine identity follows, for example, by expanding
$1/\sinh(\pi t/2)=2\sum_{j\ge0}e^{-\pi(j+1/2)t}$ and using the
logarithmic derivative of the Euler product for $\cosh$.
Fubini is valid in the displayed Hermite integral: near zero use
$|\sin(tZ)|\le t|Z|$, and at infinity the reciprocal hyperbolic sine
has exponential decay.

Restrict the positive integral to $[\sqrt n,\sqrt n+1]$.
Using $n!\le n^n$, $1/\sinh u\ge2e^{-u}$, and $1\le q\le2$ gives
\[
 |h_n|\ge\sqrt2\exp\left[-\frac n2
               -\left(1+\frac\pi2\right)\sqrt n
               -\frac{1+\pi}{2}\right]\ge e^{-6n}.
\]
The last inequality follows from $\sqrt n\le n$, $n\ge1$ and $\pi<4$.

Mehler's identity now yields the nonnegative power series
\[
 Q(x,y)=\sum_{n\ \mathrm{odd}}h_n^2
       \left(\frac{1+\sum_{i=1}^d x_i y_i}{D}\right)^n.
\]
The Boolean characters diagonalize this translation-invariant kernel.
For a character of degree $r$, choose the odd integer $n=r$ if $r$
is odd and $n=r+1$ if $r$ is even. Then $1\le n\le D$.
In the multinomial expansion of the $n$th power, choosing every
selected coordinate exactly once, and choosing the constant coordinate
once when $r$ is even, contributes $n!/D^n$ to this character's
normalized eigenvalue. Every other contribution is nonnegative.
Therefore its eigenvalue for $Q/m$ is at least
\[
 h_n^2\frac{n!}{D^n}\ge e^{-12n-D}\ge e^{-13D}.
\]
Here $n!\ge(n/e)^n$, and
$n[1+\log(D/n)]\le D$, prove $n!/D^n\ge e^{-D}$.
Multiplying by $m\ge1$ proves the stated lower bound for $Q$.
\end{proof}

\begin{corollary}[An exponential-in-depth critical-rate capacity obstruction]
\label{cor:critical-linear-depth-obstruction}
For every $d\ge1$, set $r=20(d+1)$ and $L=r+4$.
Then every Boolean function satisfies $J_L(f)\le100000L$.
Moreover, at every $N\ge\lceil e^{100(d+1)}\rceil$,
\[
 \mu_{N,1}(A_{N,f})\ge\tfrac14e^{-200000NL}
 \qquad\hbox{for every Boolean }f.
\]
Thus at $K=100000L$ the logarithm of the critical-rate sublevel
cardinality is $2^d\log2=\exp(\Omega(L))$.
\end{corollary}

\begin{proof}
Repeat the mean-shift constructions from
Theorem~\ref{thm:critical-counting-obstruction} and
Proposition~\ref{prop:critical-finitewidth-capacity}, replacing
$p=e^{-20(d+1)^5}$ by $p=e^{-20(d+1)}$ and
$r=20(d+1)^5$ by $r=20(d+1)$.
Lemma~\ref{lem:first-tanh-spectrum} gives the required covariance lower
bound for every $d\ge1$. The same estimates
$\log(1/\eta)<11(d+1)$ in the variational construction and
$\log(1/\eta)<12(d+1)$ in the finite-width construction ensure that
$r$ amplification layers suffice. The three finite-width concentration
thresholds remain below $e^{100(d+1)}$ upon the same substitutions.
All entropy costs and layer counts are unchanged.
\end{proof}

% Covariance-control result: independently audited PASS; see critical-constant-rate-fresh-audit.txt.
\section{Bounded critical rate for all functions at quadratic depth}
\label{sec:critical-constant-rate}

Covariance changes are cheaper than coherent mean shifts when the
relative change is small. The following sharper capacity obstruction
uses that fact. It concerns the critical variational rate, and does not
compare the optimal rates of competing functions closely enough to
prove a posterior sample-complexity lower bound.

\begin{theorem}[Universal bounded rate at quadratic depth]
\label{thm:critical-constant-rate}
There are universal finite constants $A,K_0>0$ such that, for every
$d\ge1$, there is an integer depth $L_d\le A(d+1)^2$ for which
\[
 J_{L_d}(f)\le K_0
 \qquad\hbox{for every Boolean }f:\{-1,1\}^d\to\{-1,1\}.
\]
The depth $L_d$ is the same for all functions on that cube.
\end{theorem}

We first record exact scalar formulas for a controlled feature law.
Put
\[
 F(q)=\mathbb E\tanh^2(\sqrt q Z),\qquad
 a(q)=1-F(q),\qquad
 R(q)=\frac{q[1-F(q)]^2}{F(q)}\quad(q>0).
\]
The value $R(0)=1$ makes it continuous at zero.

\begin{lemma}[Rank-one covariance control]
\label{lem:rankone-critical-control}
Suppose a feature covariance $Q$ has every diagonal entry equal to $q$
and $Q\succeq vff^T$ for $v>0$, where $f$ is Boolean.
For any $\epsilon\ge0$, let the next feature law be coordinatewise tanh
of a centered Gaussian with covariance $Q+\epsilon vff^T$.
Its entropy relative to the ordinary next-layer law $\mathsf M_Q$ is
at most $\epsilon^2/4$. Its second moment $Q'$ has constant diagonal
$q'$ and obeys $Q'\succeq v'ff^T$, where
\begin{equation}\label{eq:critical-covariance-control}
 u=q+\epsilon v,\qquad
 q'=F(u),\qquad v'=(1+\epsilon)v[1-F(u)]^2.
\end{equation}
Consequently, for $w=v/q$ and $w'=v'/q'$,
\[
 w'=w\frac{1+\epsilon}{1+\epsilon w}R(u).
\]
All these scalar quantities depend only on $q,v,\epsilon$, not on the
other entries of $Q$ or on the choice of $f$.
\end{lemma}

\begin{proof}
The Loewner bound implies $f\in\operatorname{range}(Q)$ and
$v f^TQ^+f\le1$. The two Gaussian laws have the same support, and the
rank-one determinant identity gives their relative entropy as
\[
 \frac12\{t-\log(1+t)\},\qquad
 t=\epsilon v f^TQ^+f\le\epsilon.
\]
This is at most $t^2/4\le\epsilon^2/4$. Applying coordinatewise tanh
can only decrease relative entropy.
Write the changed Gaussian as
$\sqrt{(1+\epsilon)v}\,fZ+G$, where $Z$ is standard Gaussian and
$G$ is an independent centered Gaussian of covariance $Q-vff^T$.
Its conditional feature mean given $Z$ is $f\psi(Z)$, since the
variances of the coordinates of $G$ all equal $q-v$.
Thus $Q'\succeq\mathbb E\psi(Z)^2 ff^T$.
Gaussian integration by parts and projection onto the first Hermite
polynomial give
\[
 \mathbb E\psi(Z)^2\ge(\mathbb E Z\psi(Z))^2
     =(1+\epsilon)v\,[\mathbb E\operatorname{sech}^2(\sqrt u Z)]^2.
\]
This is the claimed $v'$, and the diagonal formula and ratio identity
follow directly. In particular, these formulas automatically retain
$0<v'\le q'$, since they arise from a valid covariance lower bound.
\end{proof}

\begin{proof}[Proof of Theorem~\ref{thm:critical-constant-rate}]
Write $D=d+1$ and $m=2^d$. Begin with the unperturbed first feature law.
Its covariance $Q$ has constant diagonal
$q_0=F(D/d)\in[F(1),F(2)]$, and
Lemma~\ref{lem:first-tanh-spectrum} implies
\[
 Q\succeq e^{-20D}I\succeq v_0ff^T,
 \qquad v_0=e^{-20D}/m.
\]
This first layer costs zero entropy. We will use
Lemma~\ref{lem:rankone-critical-control} at each subsequent layer,
tracking its guaranteed scalar component $v$, not the largest possible
rank-one component of the covariance.

We collect universal scalar estimates. The function $F$ is increasing,
$0<F(q)<q$ for $q>0$, and Taylor expansion gives
\[
 F(q)=q-2q^2+O(q^3),\qquad R(q)=1+O(q^2).
\]
There consequently exist universal $u_0>0$ and $C\ge1$ such that,
for $0<u\le u_0$,
\begin{equation}\label{eq:critical-local-estimates}
 F(u)\le u-u^2,\qquad R(u)\ge1-Cu^2.
\end{equation}
We may and do take $u_0\le1/4$.
A useful global lower bound is
\begin{equation}\label{eq:critical-variance-lower}
 F(q)\ge\frac q{1+3q}.
\end{equation}
Indeed $\tanh^2 t\ge t^2/(1+t^2)$ follows from $|\sinh t|\ge|t|$.
Cauchy--Schwarz applied to $V=qZ^2$ gives
$\mathbb E[V/(1+V)]\ge(\mathbb EV)^2/\mathbb E[V(1+V)]
=q/(1+3q)$.

Choose a sufficiently small universal $v_*>0$ satisfying
\[
 v_*\le\min\{u_0/4,1/64,1/(100C)\}.
\]
Then choose a positive universal $c$ so small that
\[
 c\le\min\{v_*/4,u_0/3,1/3,1/(144C)\}.
\]
Put $\delta=c/D$.

\emph{Free cooling.}
Apply the unmodified next-layer law, that is, $\epsilon=0$, until
$q\le\delta$ for the first time. This costs no entropy.
There is a universal $c_F>0$ with
$q-F(q)\ge c_Fq^2$ for $0<q\le F(2)$: the ratio is positive and
continuous, and tends to two at zero. Hence the number of cooling
layers is at most a universal constant times $1/\delta$.
The final cooled variance $q_c$ lies in $[\delta/2,\delta]$, by
monotonicity and~\eqref{eq:critical-variance-lower}.

During cooling the ratio $w=v/q$ is multiplied by $R(q)$ at each step.
This loses at most a fixed positive factor, uniformly in the number
of cooling layers. To see this directly, on $[0,F(2)]$ one has
$|\log R(q)|\le C_R q^2$ for a universal $C_R$, by continuity,
positivity and the expansion at zero. Also
$\sum q_s^2\le c_F^{-1}\sum(q_s-q_{s+1})\le F(2)/c_F$.
Thus the cooled ratio satisfies
\[
 w_c\ge c_0v_0,
 \qquad \log(1/w_c)\le A_0D
\]
for universal $c_0>0,A_0<\infty$.

\emph{Amplifying the relative component.}
Now set $\epsilon=\delta$ at each layer until $w\ge1/2$.
Throughout this phase $q\le2\delta$. Indeed
$u=q+\delta v\le(1+\delta)q$, and the increasing function $u-u^2$
on the relevant interval gives
\[
 q'\le2\delta(1+\delta)-[2\delta(1+\delta)]^2\le2\delta.
\]
While $w\le1/2$, the exact ratio identity and
$u\le3\delta$ imply
\[
 \log(w'/w)
 \ge\frac\delta4-18C\delta^2\ge\frac\delta8.
\]
Here we used $\log(1-Cu^2)\ge-2Cu^2$, valid by the chosen constants.
The number $\tau$ of these layers is therefore at most
$8A_0D/\delta+1=O(D^2)$. Their total entropy is at most
\[
 \tau\delta^2/4\le2A_0D\delta+\delta^2/4=O(1).
\]
Furthermore, since $u\ge q$, monotonicity and
\eqref{eq:critical-variance-lower} imply
$1/q'\le1/q+3$ at every layer. Thus, at the end of this phase,
\[
 q\ge\frac1{2/\delta+3\tau}\ge\frac{c_1}{D^2},
 \qquad v\ge\frac q2\ge\frac{c_1}{2D^2}
\]
for some universal $c_1>0$.
At the same time $v\le q\le2\delta<v_*$.

\emph{Growing the absolute component.}
Set $\epsilon=16v$ until $v\ge v_*$.
The invariant $w\ge1/3$, or equivalently $q\le3v$, is preserved.
To check this, while $v<v_*$ one has
$u=q+16v^2\le4v$. If $w\ge1/2$, the ratio identity gives
$w'\ge\frac12(1-16Cv^2)\ge1/3$.
If $1/3\le w\le1/2$, it gives
\[
 \log(w'/w)\ge4v-32Cv^2\ge0.
\]
The initial ratio is at least $1/2$, so induction proves the invariant.
Using $F(u)\le u$ also gives
\[
 v'\ge v(1+16v)(1-2u)
       \ge v(1+16v)(1-8v)\ge v+4v^2,
 \qquad v'\le v+16v^2.
\]
Thus $1/v-1/v'\ge2$ until the stopping time, which is $O(D^2)$.
The final value is between $v_*$ and $2v_*$, and the total entropy
of this phase telescopes:
\[
 \sum\frac{(16v)^2}{4}
   =64\sum v^2\le16\sum(v'-v)\le32v_*.
\]

Finally choose output $z=f$. Its signed margin is one, and
$Q\succeq v_*ff^T$ bounds its Gaussian readout cost by $1/(2v_*)$.
All entropy costs plus this readout cost are bounded by a universal
constant $K_0$.
The number of hidden layers, and therefore the depth, is at most a
universal constant times $D^2$.
Crucially, $q_0,v_0$, every subsequent scalar recursion, every control,
and every stopping rule depend only on $d$. Although the full
covariance matrices depend on $f$, the resulting depth $L_d$ is common
to all Boolean functions. This proves the theorem.
\end{proof}

This construction sharpens an absolute-rate capacity obstruction. It
still does not show that the optimal rate of a complex competitor is
below, or sufficiently close to, the optimal rate of a specified simple
target. That comparison remains essential for a posterior lower bound.

\section{What ordinary training and other stability notions do not guarantee}
\label{sec:naturalness-obstructions}

The following results delimit the empirical question. They do not imply
that fragile outcomes are common under standard initialization, or that
their Boolean functions are statistically difficult.

\begin{proposition}[An open gradient-descent basin of fragile exact fits]
For every even $L\ge256$, there is a width-one, depth-$L$ tanh network
on $x\in\{-1,1\}$ with output $z(x)=x$, energy $\kappa_1\le L+76$,
and a postactivation sign-changing disturbance of size at most
\[
 \frac2L e^{-\sqrt L/4}.
\]
This disturbance is smaller than both calibrated support radii
\[
 r_{\rm ref}=\frac{.01}{\sqrt{1+(L-1)E_v}},\qquad
 r_S=\frac{.01}{\sqrt{1+S}}
\]
defined in Section~\ref{sec:model}.
For sufficiently small fixed step size, full-batch gradient descent on
the unregularized squared loss for the two labeled inputs has an
\emph{open set} of initializations converging to nearby exact fits that
still violate both support conditions.
\end{proposition}

\begin{proof}
Put $k_s=\min(s,L-s)$ and
\[
 r_s=\frac{\exp[-(k_s-1)/\sqrt L]}{2k_s},\quad
 a=\operatorname{atanh}(1/2),\quad
 b_s=\frac{\operatorname{atanh}(r_s)}{r_{s-1}},\quad v=2.
\]
Use first-layer input coefficient $a$ and bias zero, followed by scalar
weights $b_s$. Then $h_s(x)=r_sx$ and $z(x)=x$.
The reciprocal-ratio estimate from the bounded-excess construction in~\cite{previous} gives $\kappa_1\le L+76$; replacing its constant input by
$x$ preserves the computation and reduces the first-layer energy charge.
Also $b_s\ge e^{-1/2}/2>0.30$.

At $x=1$, the middle-layer disturbance $-2r_{L/2}$ changes the entire
remaining computation to its negative and gives final output $-1$.
Its norm is at most $2L^{-1}e^{-\sqrt L/4}$. Here $E_v=4$, so the ratio
of this upper bound to $r_{\rm ref}$ is at most
$400e^{-\sqrt L/4}/\sqrt L<1$ for $L\ge256$.
Since $S=\kappa_1\le L+76$ and $L+77\le4L-3$, it is also below $r_S$.

For the training claim, let $F(\theta)=(z_\theta(1),z_\theta(-1))$ and
$\ell(\theta)=\frac12\|F(\theta)-(1,-1)\|^2$.
At the constructed point $\theta_*$, the columns of $DF$ corresponding
to the first input coefficient and bias are $(c,-c)$ and $(c,c)$,
respectively, with $c>0$. Thus $DF$ has full row rank.
In a sufficiently small closed ball about $\theta_*$, its least singular
value is at least $\mu>0$, its norm is at most $G$, and
$\|\nabla^2\ell\|\le M$. The clean margins remain $>1/2$, and the
\emph{fixed} middle-layer attack remains strictly inside both calibrated
balls and gives strictly negative output.

For $0<\eta\le1/M$, the descent lemma gives, as long as the iterates
remain in this neighborhood,
\[
 \ell(\theta_{j+1})\le(1-\eta\mu^2)\ell(\theta_j).
\]
Their total movement is at most
\[
 \frac{\eta G\sqrt{2\ell(\theta_0)}}{1-\sqrt{1-\eta\mu^2}}
 \le \frac{2G\sqrt{2\ell(\theta_0)}}{\mu^2}.
\]
Starting in a sufficiently small open ball makes this displacement smaller
than the distance to the boundary, proving by induction that the iterates
stay inside. The summable increments give convergence, and the geometric
loss decay gives exact interpolation at the limit. Every limit retains the
strict fragility witness.
\end{proof}

A nondegenerate Gaussian initialization gives this open basin positive
probability. There is no useful lower bound here on that probability;
the result rules out a universal guarantee, not a typical-case guarantee.
It concerns unregularized full-batch squared-loss GD, rather than arbitrary
weight decay, Adam, or a prescribed finite training budget.

The same construction also violates the numerical-stability rule with
the radius in~\eqref{eq:numeric-radius} and permitted output change $1/16$.
For even $L\ge1024$, the attack-to-radius ratio is at most
$4000e^{-\sqrt L/4}/\sqrt L<1$, and the attack changes the output by two.
The same open-basin argument applies. Thus that alternative condition is
also not a universal consequence of ordinary gradient descent.

\begin{proposition}[Parameter flatness need not give activation robustness]
The networks just constructed preserve their correct signs under every
simultaneous Euclidean parameter perturbation of norm at most $1/4$.
At the exact fit their squared-loss Hessian has operator norm $O(L)$.
Nevertheless their absolute internal-noise radius is at most
$2L^{-1}e^{-\sqrt L/4}$.
\end{proposition}
\begin{proof}
Every internal scalar weight and the readout stay positive under the stated
parameter perturbation. At both inputs the first signed preactivation is
at least $\operatorname{atanh}(1/2)-\sqrt2/4>0$. Odd monotonic tanh
layers therefore retain the signs throughout the parameter ball.

For $t>0$, $t\tanh'(t)/\tanh(t)\le1$, since
$\sinh(t)\cosh(t)\ge t$. Each internal log-weight elasticity of the
positive output is consequently at most one. All scalar weights are
bounded below by a positive numerical constant, so every parameter
derivative at either input is bounded by a numerical constant.
There are $L+1$ parameters, and at interpolation
$\nabla^2\ell=DF^\top DF$; its norm is therefore $O(L)$.
\end{proof}

This is a separation between stability notions. Parameter flatness could
still support a different learning theorem, but it cannot be substituted
for absolute activation robustness in the existing proof.

\begin{proposition}[Eventual membership after readout amplification]
Every fixed network with nonzero outputs on the finite Boolean cube belongs
to either calibrated support after sufficiently large positive readout
amplification.
\end{proposition}
\begin{proof}
Continuity and finiteness of the input set give $\rho,m>0$ such that every
hidden disturbance of total normalized norm at most $\rho$ preserves each
clean sign and leaves unscaled signed output at least $m$.
Multiplying the readout by $A$ makes either calibrated radius $r_A$ tend
to zero. For large enough $A$, $r_A<\rho$ and $Am>r_A$, so even adding
the permitted scalar output noise cannot change a sign. Increasing $A$
also makes the clean output margin exceed $1/2$.
\end{proof}

This does not remove the Gaussian cost: the readout energy grows as $A^2$,
and the required $A$ may be exponentially large. It does mean that empirical
membership before and after output rescaling answers different questions.
The calibrated condition is a joint restriction on robustness and scale.

\begin{proposition}[Global and typical margins can differ exponentially]
There is a one-hidden-unit network with bounded input weight norm and
Boolean function $x_1$ whose minimum output margin is $\tanh(2^{-d})$.
Enforcing the original global clean margin $>1/2$ by readout amplification
requires a multiplier at least $2^{d-1}$, while permitting an exceptional
fraction of inputs requires only polynomial amplification in the reciprocal
exceptional fraction.
\end{proposition}
\begin{proof}
Take
\[
 a_d(x)=\sum_{j=1}^d2^{-j}x_j,\qquad z_d(x)=\tanh(a_d(x)).
\]
The squared input weight norm is less than $1/3$. The first coefficient
exceeds the sum of the others, so the sign is $x_1$, and the minimum
absolute preactivation is $2^{-d}$. Thus the global clean-margin multiplier
must exceed $[2\tanh(2^{-d})]^{-1}\ge2^{d-1}$.
Under the uniform cube, $a_d$ is uniform on the odd multiples of $2^{-d}$
in $(-1,1)$, giving
\[
 \Pr(|a_d(X)|\le u)\le u+2^{-d}\qquad(0\le u\le1).
\]
Consequently for $A\ge1$,
\[
 \Pr(|Az_d(X)|\le1/2)
 \le\operatorname{atanh}(1/(2A))+2^{-d}
 \le A^{-1}+2^{-d}.
\]
For example, $A\ge2/\epsilon$ suffices when $2^{-d}\le\epsilon/2$.
\end{proof}

The last example concerns the inherited global clean-margin convention,
not an additional failure of internal robustness. It motivates distinguishing
training, held-out, and worst-input quantities, and investigating support
conditions that explicitly tolerate a small exceptional input mass.

\begin{proposition}[A generic Gaussian halfspace has a tiny cube margin]
Let $d\ge64$ and $w\sim\mathcal N(0,I_d/d)$. Then
\[
 \Pr\left\{\min_{x\in\{-1,1\}^d}|w\cdot x|
                    \le d^2 2^{-d}\right\}\ge1-\frac6d.
\]
On this event, the one-hidden-unit output $A\tanh(w\cdot x)$ needs
$A>2^{d-1}/d^2$ to satisfy the global clean margin $>1/2$.
This is a statement about Gaussian coefficients, not the distribution of
weights after training and not a sample-complexity lower bound.
\end{proposition}
\begin{proof}
For $0<\tau\le d^{-1/2}$ put
$Z=2^{-d}\sum_x\mathbf1\{|w\cdot x|\le\tau\}$ and
$p=\mathbb EZ=\Pr(|G|\le\tau)$ for $G\sim\mathcal N(0,1)$.
For independent uniform cube points $x,x'$, their Gaussian evaluations
have correlation $\rho=d^{-1}\sum_jx_jx'_j$.
When $\rho\ne\pm1$, the bivariate Gaussian density bound gives
\[
 \frac{\Pr(|w\cdot x|\le\tau,\ |w\cdot x'|\le\tau)}{p^2}
 \le\frac{e^{\tau^2}}{\sqrt{1-\rho^2}},
\]
using $p\ge\sqrt{2/\pi}\,\tau e^{-\tau^2/2}$.
The antipodal or identical pairs have total probability $2^{1-d}$.
For $|\rho|\le1/2$, $(1-\rho^2)^{-1/2}\le1+\rho^2$.
Otherwise, excluding the two degenerate correlations, this factor is at
most $\sqrt{d/2}$. Since $\mathbb E\rho^2=1/d$ and
$\Pr(|\rho|>1/2)\le2e^{-d/8}$,
\[
 \mathbb E\!\left[
   \frac{\mathbf1\{\rho\ne\pm1\}}{\sqrt{1-\rho^2}}\right]
 \le1+\frac1d+\sqrt{2d}\,e^{-d/8}\le1+\frac2d.
\]
Here the last inequality holds for $d\ge64$.
Using $p\ge\tau/2$ and $e^{\tau^2}\le1+2/d$, Chebyshev gives
\[
 \Pr(Z=0)\le\frac{\operatorname{Var}Z}{p^2}
 \le\frac{2^{2-d}}\tau+\frac5d.
\]
Take $\tau=d^22^{-d}$; this is at most $d^{-1/2}$ for $d\ge64$.
The right-hand side is at most $6/d$. The output-scaling claim follows
from $|\tanh u|\le|u|$.
\end{proof}

The same observation includes a Gaussian first-layer bias by passing to
the $(d+1)$-dimensional cube: absolute affine values on the original cube
are exactly absolute homogeneous values on one representative from every
antipodal pair of the augmented cube.

\IfFileExists{empirical-results.tex}{\section{Expanded experiments on trained networks}
\label{sec:expanded-experiments}

The strongest positive observation is that twenty unmodified trained digit
networks satisfy the new numerical-stability condition at \emph{every real
input}, with interval-verified matrix inequalities and an exact rational
final bound. All twenty also fit every training example with signed output
greater than $1/2$. The condition is nevertheless not universal: a separate
well-performing trained network has an interval-verified violation even
at an actual held-out image. The old sign condition and the new numerical
condition must be kept separate when interpreting these findings.

\subsection{Protocol and what is being measured}

We trained 168 networks in five retained cohorts: 80 initial synthetic
models, 32 synthetic extensions, 32 main digit models, 16 deeper digit
models, and eight smaller-step-size controls. Configurations were fixed
before training each cohort; later cohorts were chosen after some earlier
results were visible. These are exploratory experiments, not a publicly
preregistered confirmatory study. Failed fits were retained. No Bayesian
posterior was sampled: the experiments test properties of trained
parameters, not posterior prediction probabilities or statistical sample
complexity. The new numerical condition was proposed after training, and
its evaluation is a post-training audit of the resulting networks.

The architecture is $h_0=(x,1)$,
$h_s=\tanh(h_{s-1}W_s)$, and $z=h_Jv$, with no biases beyond the input
augmentation. Synthetic widths are $32$ or $64$, dimensions $8,10,12$,
and hidden depths $4,8,16,32$. Tasks include majority, parity, random
Boolean labels, and labels from a fixed small tanh teacher. Training
uses either the whole cube or a fixed subset, with all cube inputs
subsequently audited. Image models have width $64$ and hidden depths
$4,16,32,64$. Initialization uses Gaussian first and final weights and
orthogonal internal matrices, with gains $0.7$, $1.1$, or $1.5$ in the
synthetic cohorts and $1.1$ for images.

Both optimizers run for 6000 steps with $10^{-6}$ weight decay.
The code's label \texttt{sgd} means \emph{full-batch vanilla gradient
descent}, without momentum or minibatch noise. Its learning rate is
$0.05$. Adam uses learning rate $0.003$, except in the separately
declared $0.0003$ control; rates are multiplied by $0.3$ after $80\%$
of training. Synthetic losses are half the squared error, while image
losses are binary logistic. No test-driven rescaling is used for images.
Exact saved floating-point weights, data splits, protocols, source, and
detailed records accompany the report.

For clarity, write $E_v=N\|v\|_2^2$ and
\[
 \|\eta\|_*^2=\sum_{s=1}^J\frac{\|\eta_s\|_2^2}{N}
                    +|\eta_{\rm out}|^2.
\]
There are four distinct measurements:
\begin{enumerate}
\item correctness relative to the target labels;
\item the inherited global clean-margin requirement $|z(x)|>1/2$;
\item preservation of the network's own clean sign throughout radius
      $r_{\rm sign}=10^{-2}/\sqrt{1+JE_v}$;
\item numerical output change at most $1/16$ throughout
      $r_{\rm num}=10^{-3}/\sqrt{1+JE_v}$, as in
      Section~\ref{sec:numerical-stability}.
\end{enumerate}
The last condition imposes no all-input clean margin. Its likelihood still
requires $y_i z(x_i)>1/2$ for \emph{every training example}; we audit this
separately on the raw trained networks. A constant failed fit can be stable
without being eligible under this likelihood.

Positive pointwise tests use the finite-radius clean-transport certificate,
including nonlinear Taylor remainder bounds; they are not just tests of
the clean Jacobian. Unless explicitly marked interval verified, the
sufficient inequalities are evaluated in float64. Failure to certify is
reported as unresolved. Negative witnesses use explicit perturbations
strictly inside the normalized noise ball. Selected witnesses and the
twenty global positive results receive the stronger checks described below.

\subsection{Full-cube synthetic audits}

Of 112 synthetic runs, 105 fit every training label, and all 105 have
correct signed training outputs greater than $1/2$. Under the old global
sign-support condition, 52 pass a sufficient certificate on their entire
cube, 59 fail the clean-margin requirement at an explicit cube point,
and one remains unresolved. Seventy networks certify their own sign on
the whole cube, independently of the clean-margin requirement; 35 have
explicit, strictly budget-checked sign-flipping perturbations.

\begin{center}
\begin{tabular}{lrrrrr}
\toprule
Cohort & Runs & Fits & Old support & All-cube sign & Sign attacks\\
 & & & in/out/? & certified & found\\
\midrule
Full cube, $d=8$ & 24 & 23 & $23/1/0$ & 24 & 0\\
Half cube, $d=8$ & 24 & 23 & $12/12/0$ & 17 & 6\\
Initialization gains, $d=8$ & 16 & 16 & $7/9/0$ & 10 & 5\\
Subset, $d=10$ & 16 & 15 & $6/10/0$ & 10 & 5\\
Subset, $d=12$ & 8 & 8 & $2/6/0$ & 5 & 2\\
Small teacher, $d=10$ & 16 & 13 & $0/15/1$ & 0 & 14\\
Depth $32$, $d=10$ & 8 & 7 & $2/6/0$ & 4 & 3\\
\bottomrule
\end{tabular}
\end{center}
Here ``out'' means a witnessed clean-margin failure; ``?'' is not a
nonrobustness claim. Counts include failed fits.

The apparent fragility is concentrated near the decision boundary.
Across all 105 successful fits, \emph{every} cube point whose clean
absolute output exceeds $1/2$ passes the finite-radius sign certificate:
80,603 input--model pairs in total, with a nonempty qualifying set for
every model. Seven networks classify their entire held-out set correctly
but fail the global clean margin; all seven certify sign stability on
their whole cube. Two depth-16, dimension-12 full-batch-GD models achieve
held-out accuracies $99.55\%$ and $99.44\%$ yet admit sign attacks at rare
inputs with clean margins about $0.0072$ and $0.0021$.

Rescaling the readout using the global minimum margin makes 108 of the
112 networks pass the old support certificate, leaving four unresolved.
This is a changed-network diagnostic, not a raw training outcome:
rescaling raises the margin and simultaneously shrinks the required noise
radius, and may incur a large Gaussian energy cost. On a finite cube,
eventual membership under sufficiently large rescaling is possible for
every network with no zero clean outputs, so this statistic alone provides
little evidence of a natural training bias.

At the new numerical radius, 108 of 112 raw networks pass on the entire
cube, including \emph{all 105} training-margin-eligible fits. Of 90,112
input--model pairs, 90,108 pass; the four unresolved points belong to
failed training runs. These positive counts use floating-point evaluation
of sufficient bounds, not interval arithmetic.

\subsection{Actual images, depth, and failed optimization}

We use the locally bundled \texttt{sklearn} handwritten-digit dataset:
1797 images of $8\times8$ pixels, with tasks ``3 versus 8'' and
``even versus odd''. A $70/30$ split is stratified by the original digit
labels. Boolean features threshold pixels at $8$; continuous features
map a pixel value $p$ to $p/8-1$. These are distinct experiments. The
continuous feature results extend the numerical audit, not the Boolean
input premise of the statistical theorem. The main cohort crosses both
tasks, both feature maps, depths $4,16$, both optimizers, and two seeds.
All 32 models satisfy the training-margin likelihood.

\begin{center}
\begin{tabular}{llrrr}
\toprule
Task & Features & Test accuracy & Own-sign certificate & Margin and sign\\
\midrule
3 versus 8 & Boolean & $95.37$--$97.22\%$ & $100\%$ & $98.15$--$100\%$\\
3 versus 8 & Continuous & $97.22$--$100\%$ & $100\%$ & $99.07$--$100\%$\\
Even/odd & Boolean & $93.33$--$95.00\%$ & $99.63$--$100\%$ & $98.89$--$100\%$\\
Even/odd & Continuous & $97.04$--$98.15\%$ & $99.81$--$100\%$ & $98.89$--$100\%$\\
\bottomrule
\end{tabular}
\end{center}
Each row contains eight models; ranges are across those models. The last
two columns concern \emph{held-out image inputs} and the old radius,
not all possible inputs. Every held-out point with $|z|>1/2$ passes
the sign certificate in all 32 models. Thresholding can give identical
feature vectors in different images: at most 23 held-out vectors per
model occur in the training set. Counts and accuracy restricted to novel
feature vectors are recorded separately; these are descriptive robustness
experiments, not a benchmark generalization claim.

The separate depth-$32/64$ cohort uses only the 3-versus-8 task. All eight
full-batch-GD models fit the training likelihood and achieve
$96.30$--$100\%$ test accuracy. All eight standard-rate Adam runs fail
to fit: depth-32 accuracies range from $50.93\%$ to $75.93\%$, and all
four depth-64 runs collapse to $50.93\%$. The separately declared
smaller-rate Adam cohort is retained alongside these failures: seven of
eight fit every training margin and achieve $96.30$--$100\%$ test accuracy;
one depth-64 Boolean-feature run still collapses to $50.93\%$.
Stable constant predictions in failed runs do not support a claim about
successful learning. No failed optimization run was replaced by a
successful rerun in the reported cohort.

\subsection{Off-data sign failures and global numerical stability}

For each of the sixteen Boolean-feature models in the main digit cohort,
we sampled 10,000 uniform 64-bit inputs. All sixteen have explicit
clean-margin violations and sign-flipping perturbations at such inputs;
all sixteen attacks still work after deleting scalar output noise.
These inputs generally do not resemble digit images. A separate attack
on 32 sampled points just above clean margin $1/2$ per model found no
flips; a failed attack is not a certificate.

Three hidden-only sign witnesses are checked using exact rational noise
budgets and 60-digit interval forward evaluation of the saved binary
floating-point weights and perturbations. Their outputs change as follows:
\[
 \begin{array}{c|rr}
 \text{Model ID}&\text{clean output}&\text{perturbed output}\\\hline
 203&+0.00164981&-0.0274141\\
 204&+0.0287959&-0.364879\\
 205&-0.0192534&+0.300013
 \end{array}
\]
All three already fail the inherited $1/2$ margin. They therefore do not
show that added internal robustness excludes a parameterization which
otherwise meets the original global-margin support restriction.

The new numerical condition gives a substantially better positive result.
All 10,368 held-out input--model pairs in the main digit cohort pass its
sufficient inequality. Twenty of these 32 raw trained models---all sixteen
depth-4 models and four depth-16 models---have a \emph{global} positive
certificate, covering every real input, not only images or the cube.
Ten of these twenty use Boolean features and ten use continuous features.
These twenty results are additionally interval verified. For each
hidden-to-hidden matrix, an exact dyadic-integer Gram matrix and a
90-digit interval $LDL^\top$ decomposition prove
$W^\top W\prec g^2I$ for an explicitly stored rational $g^2$.
Cauchy--Schwarz and global Lipschitzness of tanh bound the simultaneous
noise gain. The final comparison to $(1/16)^2$ is an exact rational
inequality. The resulting output-change upper bounds range from
$0.00453$ to $0.04950$, all below $1/16$. This is a numerical proof
trusting the interval-arithmetic implementation, not a Lean formalization.

The radius calibration matters. Reusing these same verified operator bounds
and rescaling the final inequality with exact rational arithmetic, sixteen,
twelve, and four of these twenty models also certify output change at most
$1/16$ when the radius coefficient $0.001$ is increased to $0.002$, $0.005$,
and $0.01$, respectively. The last coefficient equals the old sign-radius
coefficient, although numerical stability and sign stability remain
different properties. A failed sufficient inequality at a larger radius
does not establish a violation. The rescaling code and records are supplied
separately from the frozen experiment archive.

\begin{center}
\small\begin{tabular}{lrrrr}
\toprule
Cohort & Models & Training-margin & Numerical certificate & Full tested set\\
 & & eligible & input--model pairs & certified models\\
\midrule
Synthetic, full cubes & 112 & 105 & $90{,}108/90{,}112$ & 108\\
Main digits, held out & 32 & 32 & $10{,}368/10{,}368$ & 32\\
Deep digits, held out & 16 & 8 & $1{,}728/1{,}728$ & 16\\
Smaller-rate digits, held out & 8 & 7 & $863/864$ & 7\\
\bottomrule
\end{tabular}
\end{center}
The deep numerical certificates in this table concern held-out images,
not global input support. Three failed synthetic fits and eight failed
standard-rate deep image fits pass their entire tested set; this illustrates
why the likelihood eligibility column is essential. The one unresolved
smaller-rate image point has the negative witness described next.
Pooled input--model counts reuse inputs across models and are not counts
of independent test examples for statistical inference.

Model 203 illustrates the distinction particularly clearly. It has the
old hidden-only sign flip displayed above, but the new global certificate
bounds output changes by $0.004744$ throughout the smaller numerical
radius, for every real input. Numerical stability intentionally allows
small-margin decisions to change sign; it need not rule out otherwise
useful trained networks because of a decision boundary away from the data.

There are also genuine counterexamples to automatic numerical stability.
The smaller-rate Adam model 400, with depth 32 and Boolean pixels, fits
every training margin and achieves $98.15\%$ test accuracy. At held-out
image index 25, an exactly budget-checked perturbation changes its output
from about $+0.188308$ to $-0.006960$, an absolute change of $0.195268$.
A 120-digit interval calculation proves that this exceeds $1/16$.
The true label is $-1$, so this particular perturbation \emph{corrects}
an existing classification error; it is a counterexample to numerical
stability, not an induced label error. Other verified violations occur
on uniform-cube inputs. One changes $5.296417$ to $5.218443$, which violates
raw numerical stability while leaving both clipped scores equal to one.
This supports considering the bounded-score condition of
Section~\ref{sec:clipped-stability}; it does not establish that every
trained model satisfies that weaker condition. In fact, clipping leaves
model 400's small-output held-out violation unchanged, so that trained
model also violates the bounded-score condition at the same radius.

\subsection{Other diagnostics and limits of the evidence}

The product of layer Frobenius norms can be enormous even when the new
stability test succeeds. For raw main digit models, the median
$\log_{10}(\sqrt d\prod_s\|W_s\|_F)$, including the readout as its last
factor, is $4.80$ at depth 4 and $15.74$ at depth 16. In the standard
deeper cohort it is $31.34$ at depth 32 and $58.96$ at depth 64. Thus a
small product-Frobenius condition can be much more restrictive on these
trained parameterizations than finite-noise output stability. These finite
depths do not establish an asymptotic exponential law, and the experiment
does not exclude a lower-product implementation of the same function.

A separate, fixed activation-rounding diagnostic rounds every hidden
activation and the final output to $8$, $12$, or $16$ binary fractional
bits, without rounding weights. Across the main digit models, 8-bit
rounding changes two held-out signs; 12- and 16-bit rounding change none.
Only 16 bits uniformly place the worst-case rounding budget within the
old radius for all 32 models. Success of one deterministic rounding rule
is weaker than robustness against every allowed perturbation.

These are modest tanh multilayer perceptrons, not evidence about large
contemporary ResNets or Transformers. The strongest positive global
certificates occur at depths 4 and 16; favorable on-image behavior at
depths 32 and 64 does not prove all-input membership. None of these finite
experiments proves a polynomial-depth scaling law or a universal property
of gradient descent. They support the naturalness of numerical stability
for some successful trained networks, expose its exceptions, and show why
training correctness, confident-point stability, global margin, and global
numerical stability should be reported separately.
}{}
\section{A posterior certificate using robustness only where it is needed}
\label{sec:pointwise}

The earlier compression proof uses robustness only at its final prediction
comparison. This permits a version that does not assume robust support.
Fix a radius $0<r\le1$. Define
\[
 b_w(x;r)=\mathbf1\{\text{some perturbation of total RMS norm at most }r
 \text{ makes }\sgn z_w(x)\,z_w(x;\eta)\le0\}.
\]
Clean zeros are declared bad. For a parameter law $Q$, write
$b_Q(x;r)=\E_Q b_w(x;r)$, and let $b_D(Q;r)$ and $b_s(Q;r)$ be its
population and empirical input averages. The reference measure $D$ here is
the input distribution, not a target-mass exponent.

\begin{proposition}[Compression with a pointwise defect]
For $D_{\rm KL}(Q\Vert\mu_{N,t})\le NK/t^2$, the same finite class
$\mathcal H(d,L,K,r,\alpha)$ from the earlier robust entropy-compression
theorem admits a probability law $\Gamma_Q$ with
\[
 \left|\E_Q\mathbf1\{\sgn z_w(x)\ne f(x)\}
       -\E_{h\sim\Gamma_Q}\mathbf1\{h(x)\ne f(x)\}\right|
 \le\alpha+b_Q(x;r)\qquad(x\in\{-1,1\}^d).
\]
The construction and the class do not depend on the target $f$.
\end{proposition}
\begin{proof}
Run the earlier construction for the unrestricted law $Q$. Its entropy
budgets, clipping bounds, moment estimates, and common candidate marginal
do not use robustness. At each fixed input, its comparison coupling has
the original full-network marginal $Q$ and, except with probability less
than $\alpha/4$, produces an exact postactivation perturbation of norm
less than $r/2$. When $b_w(x;r)=0$, the candidate has the original sign
and margin greater than $r/2$: otherwise an additional output injection
within the remaining norm budget could reach zero. The same finite-grid
rounding preserves that sign. In the small-width branch, good parameter
clipping likewise preserves the original network until rounding; local
robustness gives its output margin greater than $r$. Adding the probability
$b_Q(x;r)$ to the original exceptional probability proves the assertion.
No conditioning on a robustness event is performed, so no extra entropy
cost or inconsistency between pointwise couplings is introduced.
\end{proof}

\begin{theorem}[Robustness-defect certificate for a consistent posterior]
With probability at least $1-\delta$ over $n$ iid labeled examples,
simultaneously for all $N,t,Q$ with the stated entropy bound,
\[
 \operatorname{GibbsErr}_D(Q)
 \le 2\operatorname{GibbsErr}_s(Q)+3\alpha+b_D(Q;r)+2b_s(Q;r)
       +\frac{8[\log|\mathcal H|+\log(1/\delta)]}{n}.
\]
In particular the empirical Gibbs error term vanishes for the original
noiseless-consistent Gaussian posterior. This statement imposes no
robustness restriction on that prior.
\end{theorem}
\begin{proof}
On the usual finite-class relative-Chernoff event,
$\operatorname{err}_D(h)\le2\operatorname{err}_s(h)
+8[\log|\mathcal H|+\log(1/\delta)]/n$ for every $h\in\mathcal H$.
Average under the common law $\Gamma_Q$ and use the preceding proposition
once on the population and once on the empirical inputs.
\end{proof}

For a fixed posterior selected without using an independent unlabeled
sample $U$ of size $m$, Hoeffding also gives
\[
 b_D(Q;r)\le b_U(Q;r)+\sqrt{\log(1/\delta')/(2m)}
\]
with probability at least $1-\delta'$. Sound upper certificates for
$b_w(x;r)$ may replace its exact value. This is a mathematical posterior
certificate; one trained network is not automatically a draw from the
specified Bayesian posterior. The training defect cannot simply be omitted:
the posterior depends on the training inputs, so small population defect
alone does not establish a small empirical defect.

\subsection{A distributional prior with robust training likelihood}
There is also a genuine learning theorem under a weaker prior restriction.
Suppose the input distribution $D$ is known before the labels, set
$J=L-1$, $E_v=N\|v\|^2$, and put
\[
 r_{\rm ref}(w)=\frac{1}{100\sqrt{1+J E_v}}.
\]
For fixed $0\le\rho<1$, condition the raw Gaussian only on
\[
 D\{x:b_w(x;r_{\rm ref}(w))=1\}\le\rho.
\]
Unlike the earlier prior, this condition imposes no all-input clean margin.
For every observed example $(x_i,y_i)$, use the likelihood requiring
$y_i z_w(x_i;\eta)>0$ for every perturbation of norm at most
$r_{\rm ref}(w)$. Thus the likelihood enforces a robust answer on training
inputs. This is a modified learning rule, not the original sign-only
likelihood with a silently omitted empirical defect.

\begin{corollary}[Polynomial-depth learning up to a distributional defect]
For targets with a width-$C$, depth-$(L-1)$ layered Boolean circuit, this
posterior has Gibbs error at most $\rho+\varepsilon$ after
$\operatorname{poly}(d,L,C,\varepsilon^{-1},\log(1/\delta))$ examples,
simultaneously for all $N\ge C$ and all $0<t\le1$. The polynomial has
degree independent of depth, width, and Gaussian scale.
\end{corollary}
\begin{proof}
The earlier robust target cylinder belongs to this prior support and
satisfies the robust training likelihood for every sample. If its raw
mass is at least $\exp(-ND_*/t^2)$, posterior KL is at most $ND_*/t^2$.
With $B=\max\{8,8[D_*+\log(1/\beta)]\}$, the same Gaussian readout tail
bound gives posterior mass at most $\beta$ to $E_v>B$. On the complementary
event the required radius is at least
$r_0=1/[100\sqrt{1+JB}]$. The conditional posterior has entropy bound
$N(D_*+1)/t^2$, population robustness defect at most $\rho$, and empirical
defect zero. Apply the preceding theorem and add the discarded mass:
\[
 \operatorname{GibbsErr}_D(Q_s)
 \le\beta+3\alpha+\rho
       +8[\log|\mathcal H|+\log(1/\delta)]/n.
\]
Choosing $\beta=\varepsilon/4$, $\alpha=\varepsilon/12$ and the earlier
polynomial sample size proves the claim. The finite class is independent
of $N,t$, which gives one simultaneous sample event.
\end{proof}

If a fixed reference input distribution $\nu$ is used instead, and
$D\le M\nu$, replace the population defect $\rho$ by $M\rho$. This makes
explicit what distributional information enters the prior. The result
does not assert that the support can be estimated from arbitrary unlabeled
data without an additional uniform-convergence argument.

\section{Numerical output stability without an all-input margin}
\label{sec:numerical-stability}

The preceding sign-robustness conditions can be replaced by a condition on
\emph{numerical changes in the output}. This permits the clean output to be
arbitrarily close to zero on unobserved inputs. Only the observed examples
are required to have the original positive training margin.

Let $J=L-1$, $E_v(w)=N\|v\|_2^2$, and set
\begin{equation}\label{eq:numeric-radius}
 r_{\rm num}(w)=\frac{1}{1000\sqrt{1+J E_v(w)}}.
\end{equation}
Use the same additive postactivation and scalar output disturbances as
before. Say that $w$ is \emph{numerically stable} if
\begin{equation}\label{eq:numeric-support}
 |z_w(x;\eta)-z_w(x)|\le\frac1{16}
 \quad\text{for all }x\in\{-1,1\}^d
 \text{ and all }\|\eta\|_*\le r_{\rm num}(w).
\end{equation}
The prior is the raw Gaussian conditioned on this event. There is
\emph{no all-input clean-margin condition}. For an observed labeled sample,
condition further on
\begin{equation}\label{eq:numeric-likelihood}
 y_i z_w(x_i)>\frac12\qquad(1\le i\le n).
\end{equation}
Thus the original output-margin convention is imposed only at training
inputs. With no global margin restriction, this likelihood is stronger
than sign-only consistency; that distinction is part of the theorem.

The condition has a simple sufficient example. If every hidden-to-hidden
matrix has operator norm at most one, then
\[
 |z_w(x;\eta)-z_w(x)|
 \le\sqrt{1+J E_v(w)}\,\|\eta\|_*.
\]
Consequently every such network satisfies~\eqref{eq:numeric-support},
regardless of how small its clean outputs are. The first layer is
unrestricted because there is no input perturbation, and the final readout
may have arbitrary gain. Numerical stability can also hold through
saturation and cancellation without nonexpansiveness.

\begin{lemma}[Numerical posterior compression]
\label{lem:numeric-compression}
Let $0<r\le1/6$, and suppose
$D_{\rm KL}(Q\Vert\mu_{N,t})\le NK/t^2$. Suppose every network in the
support of $Q$ obeys
\[
 |z_w(x;\eta)-z_w(x)|\le1/16\qquad(\|\eta\|_*\le r)
\]
at every input $x$ in a specified set $G$. There is a finite class
$\mathcal H$ of \emph{real-valued} depth-$L$ network outputs with
\[
 \log|\mathcal H|\le\operatorname{poly}(d,L,K,r^{-1},\alpha^{-1}),
\]
independent of $N,t,Q$, and a law $\Gamma_Q$ on $\mathcal H$, such that
at each $x\in G$ a coupling with marginals $Q,\Gamma_Q$ satisfies
\[
 \mathbb P\{|h(x)-z_w(x)|>1/12\}\le\alpha.
\]
The candidate marginal is the same at every input.
\end{lemma}
\begin{proof}
Run the earlier entropy compressor before using its final sign-robustness
step. Its entropy bounds and construction do not require robustness.
Except with probability less than $\alpha/4$, its unrounded candidate
output at a fixed input is exactly the original network's output under
a simultaneous postactivation disturbance of norm less than $r/2$.
The numerical stability assumption bounds the change by $1/16$.
The same fixed parameter-grid rounding changes outputs by at most $r/8$,
uniformly on the cube. Since $r\le1/6$,
$1/16+r/8\le1/12$. In the small-width branch the unrounded candidate is
the original network except on its bounded-weight failure event, and the
same rounding bound suffices. Keep the rounded real outputs, rather than
only their signs. The finite grid and the candidate marginal are otherwise
unchanged. No clean output margin has been used.
\end{proof}

\begin{theorem}[Polynomial-depth learning without a global clean margin]
\label{thm:numeric-prior}
For the prior~\eqref{eq:numeric-support} and training
likelihood~\eqref{eq:numeric-likelihood}, a target computed by a width-$C$
layered Boolean circuit has a sample bound
\[
 n\ge\operatorname{poly}(d,L,C,\varepsilon^{-1},\log(1/\delta))
\]
such that, with probability at least $1-\delta$,
\begin{equation}\label{eq:numeric-risk}
 \mathbb E_{w\sim Q_S,\,X\sim D}
       \mathbf1\{f(X)z_w(X)\le1/4\}\le\varepsilon.
\end{equation}
One sample event gives this guarantee simultaneously for every $N\ge C$
and $0<t\le1$, hence for every $\gamma\ge1$. The polynomial's degree is
independent of depth, width, and Gaussian scale. In particular, it bounds
Gibbs sign error, and gives the usual factors three and two for posterior
confidence thresholds $2/3$ and $1/2$, respectively.
\end{theorem}

\begin{proof}[Statistical proof, given the target cylinder below]
The target cylinder established below has raw Gaussian mass at least
$\exp(-ND_*/t^2)$ with $D_*=\operatorname{poly}(d,L,C)$, is numerically
stable, and has correct signed clean output greater than one on the full
cube. It therefore lies in the posterior conditioning event for every
sample. As before, posterior KL is at most $ND_*/t^2$.
The Gaussian readout tail and this partition lower bound allow discarding
posterior mass at most $\beta$ while retaining $E_v\le B$, where
$B=\operatorname{poly}(D_*,\log(1/\beta))$.
The retained law $Q'$ has KL at most $N(D_*+1)/t^2$ and numerical stability
at the common radius
$r_0=[1000\sqrt{1+JB}]^{-1}<1/6$.

Apply Lemma~\ref{lem:numeric-compression}, and define the bounded ramp
\[
 \psi(u)=
 \begin{cases}
 1,&u\le1/3,\\
 5-12u,&1/3<u<5/12,\\
 0,&u\ge5/12.
 \end{cases}
\]
On a successful comparison at a training point, the candidate has
$y_i h(x_i)>1/2-1/12=5/12$, so its ramp loss is zero. Thus the empirical
ramp risk averaged under $\Gamma_{Q'}$ is at most $\alpha$.
On a successful comparison at a test point, if $f(x)z_w(x)\le1/4$, then
$f(x)h(x)\le1/4+1/12=1/3$, so the candidate's ramp loss is one.
Therefore the original test margin-failure probability is at most the
candidate's population ramp risk plus $\alpha$.

The finite-class relative Chernoff inequality applies to these $[0,1]$
losses, and gives simultaneously for all $h\in\mathcal H$
\[
 R_\psi(h)\le2\widehat R_\psi(h)
             +\frac{8[\log|\mathcal H|+\log(1/\delta)]}{n}.
\]
For completeness, the bounded-loss version follows from the same
exponential argument as for indicators, using
$e^{-s u}\le1-(1-e^{-s})u$ for $0\le u\le1$.
Averaging under the candidate law and adding the discarded posterior mass
gives
\[
 \mathbb E_{Q_S,D}\mathbf1\{f(X)z_w(X)\le1/4\}
 \le\beta+3\alpha+
       \frac{8[\log|\mathcal H|+\log(1/\delta)]}{n}.
\]
Choose $\beta=\varepsilon/4$, $\alpha=\varepsilon/12$, and $n$ large enough
that the last term is at most $\varepsilon/2$.
The class is independent of $N,t$, proving simultaneity.
Finally, if the posterior gives probability at most $2/3$ to a correct
signed output greater than $1/4$, its conditional margin-failure probability
is at least $1/3$; Markov gives the factor three. The sign event is larger,
and the corresponding factor at $1/2$ is two.
\end{proof}

\subsection{A numerically stable target cylinder}

The robust digital-cylinder proof from the preceding report can be made
numerically stable by tightening its hidden approximation constants. Use
balanced clone groups, write $P$ for the group-constant projection, and
write $b^s(x)\in\{-1,1\}^N$ for the ideal repeated Boolean wires. Replace
that construction's parameters by
\begin{align*}
 e_*&=\frac1{128\sqrt{2C}},&
 \tau_0&=e_*/4,\\
 b_L&=\sqrt{16\log9+8\log(4L)},&
 A&=\max\{32b_L,\,2\log(8/e_*)\},\\
 \zeta&=\min\{1/4,\,A/[4(d+1)],\,Ae_*/(64\sqrt C)\}.
\end{align*}
Here $\zeta$ is the cylinder tolerance, unrelated to an input-distribution
exception mass. Constrain the standardized first-layer coordinates, active
internal group coordinates, and readout coordinates within $\zeta t$ of
the cloned saturated circuit, exactly as in that construction. On its
free-bulk event,
\[
 W^s=W_*^s+R^s+E^s,\quad
 R^s=R^sP,\quad \|R^s\|_{\rm op}\le\zeta t\sqrt C,\quad
 E^sP=0,\quad \|E^s\|_{\rm op}\le b_Lt.
\]
The free-bulk event has conditional probability at least $1/2$.

The old induction, with these tightened constants, proves
\begin{equation}\label{eq:numeric-tube}
 \frac{\|h^s_\eta(x)-b^s(x)\|_2}{\sqrt N}\le e_*
\end{equation}
for every input, at every hidden layer, whenever each hidden injection has
RMS at most $\tau_0$ and output noise has magnitude at most $\tau_0$.
To check the constants directly, an error of RMS $e_*$ makes every clone
group's mean error at most $\sqrt{2C}e_*=1/128$. Thus every ideal gate's
signed preactivation exceeds $A/2$. The global secant estimate
\[
 |\tanh(a+u)-b|\le2e^{-A/2}+(8/A)|u|
 \quad\text{when }b\in\{-1,1\},\ ba\ge A/2
\]
gives the recurrence
\[
 e_s\le2e^{-A/2}+(8/A)\zeta\sqrt C(1+\tau_0)
               +(8b_L/A)e_{s-1}+\tau_0\le e_*.
\]
Each of the four terms is at most $e_*/4$.
The first-layer perturbation changes each preactivation by at most
$\zeta(d+1)\le A/4$, so its signed preactivation is at least $3A/4$;
its activation error plus injection is at most $e_*/2$.
This initializes the recurrence. In particular the clean computation
also satisfies~\eqref{eq:numeric-tube}.

The cloned readout obeys $\sqrt N\|v_*\|\le\sqrt{8C}$, while its
parameter perturbation satisfies $\sqrt N\|v-v_*\|\le\zeta$.
Comparing the noisy and clean computations for the \emph{same} weights
therefore yields
\begin{align*}
 |z_w(x;\eta)-z_w(x)|
 &\le2e_*(\sqrt{8C}+\zeta)+\tau_0\\
 &\le\frac1{32}+\frac34 e_*<\frac1{16}.
\end{align*}
Also $\sqrt N\|v_*\|\ge\sqrt{2C}$, so
\[
 E_v(w)\ge(\sqrt{2C}-1/4)^2\ge C.
\]
Consequently
\[
 r_{\rm num}(w)\le\frac{1}{1000\sqrt C}
       <\frac1{512\sqrt{2C}}=\tau_0.
\]
Thus every draw in the cylinder and free-bulk event satisfies the numerical
support rule. Its correct signed clean output is at least
$2-\sqrt{8C}e_*-\zeta>1$, which also establishes the training likelihood.

There are still only $N[d+2+(L-2)C]$ constrained Gaussian coordinates,
and the standardized center energy is at most
$N[9dA^2+18(L-2)CA^2+8C]$. The preceding Gaussian box calculation therefore
gives mass at least $\exp(-ND_*/t^2)$ for
$D_*=\operatorname{poly}(d,L,C)$. Tightening $e_*$ by a constant factor
changes only constants and logarithmic factors; it introduces no exponential
dependence on depth. This proves the target-cylinder assertion used above.

\subsection{A trusted-region numerical version}

If the numerical support condition is imposed only on a fixed trusted
region $G$ with $D(G^c)\le\rho$, while the training-margin likelihood is
still imposed on \emph{all} training points, the same proof uses masked
ramp losses $\mathbf1_G(x)\psi(f(x)h(x))$. It gives margin-failure risk
at most $\rho+\varepsilon$ with the same polynomial bound. A finite catalog
of trusted regions can be handled by the same mixture construction as in
Corollary~\ref{cor:trusted-catalog}, at an additional $\log M$ cost.
The global target cylinder belongs to every component.

\subsection{A soft numerical-stability prior with full support}

The hard numerical condition can also be replaced by a soft penalty. Define
\[
 \rho_{\rm num}(w)=\min_x\inf\bigl\{\|\eta\|_*:
       |z_w(x;\eta)-z_w(x)|\ge1/16\bigr\}.
\]
For every finite parameter array, continuity on the finite input cube
gives $\rho_{\rm num}(w)>0$, including when some clean outputs vanish.
Output-only noise gives $\rho_{\rm num}(w)\le1/16$.
For fixed $p,\lambda>0$, use the prior density
\[
 d\nu_{N,t}(w)\ \propto\
       \exp\bigl[-(N/t^2)\lambda\rho_{\rm num}(w)^{-p}\bigr]
       \,d\mu_{N,t}(w),
\]
and retain the training-margin likelihood~\eqref{eq:numeric-likelihood}.
This density is positive at every finite parameter array; it retains the
support of the raw Gaussian and needs no clean-margin exclusion anywhere.

The soft-prior calculation from Theorem~\ref{thm:soft-margin-prior} applies
with $\rho_{\rm num}$ in place of the sign margin. The target cylinder has
$\rho_{\rm num}\ge\tau_0$ and raw mass $e^{-ND_*/t^2}$.
Set $B_*=D_*+\lambda\tau_0^{-p}$. The training partition is at least
$e^{-NB_*/t^2}$; posterior KL is at most $NB_*/t^2$; and
\[
 Q_S\{\rho_{\rm num}^{-p}>u\}
   \le\exp\{(N/t^2)(B_*-\lambda u)\}.
\]
Choosing $u=[B_*+\log(1/\beta)]/\lambda$ discards mass at most $\beta$.
The retained law has scaled KL parameter at most $B_*+1$ and numerical
stability at radius $r=\tfrac12u^{-1/p}\le1/32$.
Lemma~\ref{lem:numeric-compression} and the same ramp-loss argument prove
the conclusion of Theorem~\ref{thm:numeric-prior}, with a polynomial bound
for fixed $p,\lambda$. The prior probabilities have changed; this is not a
theorem for unchanged independent Gaussian weights.
One may also minimize only over a trusted region $G$, obtaining the
corresponding error bound $D(G^c)+\varepsilon$.

This theorem does not identify sign-only conditioning with positive-margin
conditioning. It shows that a global clean-output lower bound is unnecessary
when numerical internal stability is imposed and the observed labels retain
a positive output margin. Numerical stability is not a universal consequence
of ordinary gradient descent: the earlier fragile scalar-chain construction
also violates this condition for sufficiently large depth. For example, at
even $L\ge1024$, its displayed attack bound divided by the present radius
is at most $4000e^{-\sqrt L/4}/\sqrt L<1$, while the attack changes the output
from $+1$ to $-1$. The same strict inequalities persist in its open attraction
neighborhood for gradient descent.

\section{Finite-noise certificates that permit expansion and contraction}
\label{sec:numeric-certificates}

Numerical stability need not be enforced by making every layer nonexpansive
in the ordinary Euclidean norm. Diagonal quadratic metrics provide a
global, finite-noise certificate that allows intermediate expansion and
subsequent contraction. They also give a concrete optimization objective.

Write $H=L-1$ for the number of hidden layers. Let $M_s$ be diagonal
positive semidefinite matrices, and suppose
\begin{equation}\label{eq:diagonal-storage}
 M_H\succeq vv^\top,\qquad
 (W^s)^\top M_sW^s\preceq M_{s-1}\quad(2\le s\le H).
\end{equation}
Define
\[
 Q_M=N\sum_{s=1}^H\lambda_{\max}(M_s).
\]

\begin{proposition}[A diagonal-metric finite-noise certificate]
\label{prop:diagonal-numeric}
Under~\eqref{eq:diagonal-storage}, for every input and every finite
simultaneous postactivation and output disturbance,
\[
 |z_w(x;\eta)-z_w(x)|\le\sqrt{1+Q_M}\,\|\eta\|_*.
\]
In particular this is a global finite-radius statement, not a bound only
on the clean Jacobian.
\end{proposition}
\begin{proof}
Coordinatewise $1$-Lipschitzness of $\tanh$ makes it nonexpansive in every
diagonal seminorm $\|u\|_{M_s}=\sqrt{u^\top M_su}$.
If $\Delta h_s$ denotes the difference of the noisy and clean states,
\eqref{eq:diagonal-storage} gives
\[
 \|\Delta h_s\|_{M_s}
 \le\|\Delta h_{s-1}\|_{M_{s-1}}+\|\eta_s\|_{M_s}.
\]
The first layer has $\Delta h_1=\eta_1$, since there is no input noise.
Therefore
\[
 |\Delta z|\le
   \sum_{s=1}^H\sqrt{\lambda_{\max}(M_s)}\,\|\eta_s\|_2
       +|\eta_{\rm out}|.
\]
Cauchy--Schwarz in the RMS-normalized disturbance coordinates gives the
claim. No assumption on activation signs, margins, or smallness of the
perturbations has been used.
\end{proof}

If all internal operator norms are at most one, taking
$M_s=\|v\|^2I$ recovers $Q_M=H E_v$. But the metric condition is weaker.
For example, consider a scalar hidden stack with alternating internal
weights $2,1/2,2,1/2,\ldots$. Let
$p_s=\prod_{j=s+1}^H|W^j|$ and take $M_s=v^2p_s^2$.
Every matrix inequality in~\eqref{eq:diagonal-storage} is an equality,
every $p_s\le2$, and $Q_M\le4H E_v$. The noise gain is therefore only
$O(\sqrt H)$ despite expansive layers at a positive fraction of depths.
The product of the factors $\max\{1,|W^j|\}$ would instead be exponential
in depth. The certificate accounts for the compensating contractions.

The certificate is also output-specific. For instance, with two hidden
coordinates, readout $v=(1,0)^\top$, and every internal matrix
$W^s=\operatorname{diag}(1,10)$, take
$M_s=\operatorname{diag}(1,0)$. The second coordinate may have a very large
local gain but never affects the output; the certificate gives
$Q_M=H E_v$ nonetheless. Thus irrelevant expansive directions need not
be suppressed.

\paragraph{Computing and training the certificate.}
For fixed weights, minimizing $\sum_s b_s$ subject to
\eqref{eq:diagonal-storage} and $0\preceq M_s\preceq b_sI$, with diagonal
$M_s$, is a semidefinite program: every displayed constraint is affine in
the diagonal metric entries and $b_s$. Its value certifies a finite-noise
gain. One can penalize this certified gain during training, or use the
parameterization
\[
 W^s=D_s^{-1}A_sD_{s-1},\qquad \|A_s\|_{\rm op}\le1,
 \qquad v=D_Hu,\quad\|u\|_2\le1,
\]
with positive diagonal $D_s$ and $M_s=D_s^2$.
Spectral normalization or projection enforces the constraints on $A_s$;
the diagonal scales allow expansion and contraction in the original
coordinates. This is a proposed training constraint or regularizer, not
an assertion that unconstrained gradient descent automatically satisfies it.

For the numerical radius~\eqref{eq:numeric-radius}, a certificate satisfying
\[
 \frac{1+Q_M}{1+H E_v}\le\left(\frac{1000}{16}\right)^2
\]
proves the prior's $1/16$ output-stability condition. More generally, let
$R(L)\ge1$ be any fixed polynomial and replace that radius by
\begin{equation}\label{eq:poly-transient-radius}
 \frac{1}{1000R(L)\sqrt{1+H E_v}}.
\end{equation}
A certificate $Q_M\le R(L)^2H E_v$ then guarantees output change at most
$1/1000$ throughout the ball. The polynomial-depth learning theorem remains
valid: the target cylinder still belongs to the support because the radius
is smaller, and the common compression radius after the readout-tail
truncation is still inverse polynomial in depth.
This permits polynomial intermediate amplification, not just contraction.
For example, internal maps with operator norms at most $1+B/L$ have suffix
gain at most $e^B$; if $B=O(\log L)$, this is polynomial and fits such a
calibration. Applying this observation to a different residual-network
architecture would still require matching that architecture's prior and
target-mass calculation.

The metric inequalities are sufficient certificates for membership in the
numerical prior. If one instead restricts the prior to these particular
certificates, the existing digital target cylinder is not automatically
known to remain inside that stricter support; that would be an additional
expressivity and probability question.

\subsection{A finite-tube version that uses saturation}

The global diagonal-metric test can be conservative for saturated units.
There is a rigorous finite-tube variant. Fix an input, a proposed radius
$r$, and positive definite diagonal $M_s$. Let $a_s=W^sh_{s-1}$ be the
clean preactivations. Put
\[
 B_s=N\sum_{j=1}^s\lambda_{\max}(M_j),\qquad
 R_{s,i}=r\sqrt{B_{s-1}}
       \sqrt{(W^sM_{s-1}^{-1}(W^s)^\top)_{ii}}
       \quad(s\ge2),
\]
and define diagonal matrices $U_s$ by
\[
 (U_s)_{ii}
  =\sup_{|q|\le R_{s,i}}\operatorname{sech}^2(a_{s,i}+q)
  =\operatorname{sech}^2\bigl((|a_{s,i}|-R_{s,i})_+\bigr).
\]
Suppose
\begin{equation}\label{eq:tube-metric}
 M_H\succeq vv^\top,\qquad
 (W^s)^\top U_sM_sU_sW^s\preceq M_{s-1}\quad(2\le s\le H).
\end{equation}
Then every disturbance with $\|\eta\|_*\le r$ changes the output at this
input by at most $r\sqrt{1+B_H}$.

Indeed, induction gives
$\|\Delta h_s\|_{M_s}\le\sum_{j\le s}\|\eta_j\|_{M_j}
\le r\sqrt{B_s}$.
The preceding-layer bound places each preactivation displacement in its
specified interval $[-R_{s,i},R_{s,i}]$. The coordinatewise mean-value
slopes are consequently between zero and the corresponding entries of
$U_s$. Since $M_s$ is diagonal,~\eqref{eq:tube-metric} gives the same
contraction inequality as before. This closes the induction and proves the
finite-noise bound, including scalar output noise by Cauchy--Schwarz.
The construction is self-consistent: the interval widths are computed
before accepting the matrix inequalities, rather than assumed from the
clean Jacobian alone.

To certify the prior, this test must hold at every required input, with
the resulting output bound at most $1/16$. Metrics may be chosen separately
at different inputs. Checking only a finite test set supplies evidence on
those inputs, not a global certificate. Unlike the global test, optimizing
the finite-tube metrics is not in general a single convex semidefinite
program, because the slope bounds depend on the candidate metrics.

\subsection{Large output amplification does not automatically repair this condition}

Let $w$ be fixed with $E_v>0$, and let $w^{(A)}$ be obtained by multiplying
only its readout by $A>0$.
For an input $x$, let
\[
 G_{\rm hid}(w,x)^2
  =N\sum_{s=1}^H
       \left\|\frac{\partial z_w(x;\eta)}{\partial\eta_s}
                     \bigg|_{\eta=0}\right\|_2^2.
\]
This is the squared clean gain from RMS-normalized hidden disturbances;
it does not include the scalar output-noise coordinate. With $c=1/1000$
and $r_A=c/\sqrt{1+H A^2E_v}$,
\begin{equation}\label{eq:readout-scaling-limit}
 \lim_{A\to\infty}
   \sup_{\|\eta\|_*\le r_A}
     |z_{w^{(A)}}(x;\eta)-z_{w^{(A)}}(x)|
   =\frac{c\,G_{\rm hid}(w,x)}{\sqrt{H E_v}}.
\end{equation}
The convergence is uniform over the finite Boolean cube.
To prove it, write $\eta=r_Au$. Smoothness gives
\[
 z_{w^{(A)}}(x;r_Au)-z_{w^{(A)}}(x)
   =A r_A\,\ell_x(u_{\rm hid})+r_Au_{\rm out}
                                      +O(A r_A^2),
\]
uniformly on the unit disturbance ball, where $\ell_x$ is the clean hidden
linear response. Since $A r_A\to c/\sqrt{HE_v}$,
$r_A\to0$, and $A r_A^2\to0$, taking the supremum gives
\eqref{eq:readout-scaling-limit}. Finiteness of the cube makes the remainder
uniform over inputs.

Consequently, if
\[
 \max_x\frac{G_{\rm hid}(w,x)}{\sqrt{H E_v}}>62.5,
\]
then no sufficiently large readout amplification makes the network satisfy
the numerical support rule. For a nonexpansive hidden stack the ratio is
at most one. This distinguishes numerical stability from the earlier
calibrated sign condition, whose required radius could always be met by
sufficiently large amplification of a nonvanishing finite-cube network.
The limiting clean gain is used here only for an asymptotic scaling
obstruction; the positive certificates above control finite disturbances.

\paragraph{A limitation of the global metric certificate on the deeper trained models.}
The global diagonal-metric certificate cannot certify the depth-32 and
-depth-64 gradient-descent digit models in this experiment, even after
optimizing its metrics. This failure has an elementary exact certificate;
an SDP solver is unnecessary. In column-vector notation, define the
\emph{unactivated suffix readouts}
\[
 g_H=v,\qquad g_{s-1}=(W^s)^\top g_s\quad(2\le s\le H).
\]
For any positive semidefinite metric family satisfying
$M_H\succeq vv^\top$ and
$M_{s-1}\succeq(W^s)^\top M_sW^s$, induction gives
$M_s\succeq g_sg_s^\top$. Consequently,
\[
 Q_M=N\sum_{s=1}^H\lambda_{\max}(M_s)
       \ge N\sum_{s=1}^H\|g_s\|_2^2=:G_{\rm lin}^2.
\]
Thus if
\[
 \frac{1+G_{\rm lin}^2}{1+H E_v}>3906.25,
\]
no such metric family can meet the proposed numerical certificate.
The argument even permits arbitrary positive semidefinite metrics in this
necessary inequality; restricting to diagonal metrics cannot help.

We preselected the two successful Boolean-input GD models with seed zero,
302 ($H=32$) and 306 ($H=64$). Their respective ratios are approximately
$93\,243.76$ and $63\,136\,897.94$, both strictly above $3906.25$.
The implied lower bounds on the output-change \emph{certificate} are
$0.30536$ and $7.94587$, versus the required $1/16=0.0625$.
These inequalities were checked with exact integer dyadic matrix--vector
products and rational arithmetic on the saved weights. An exploratory
check of the other six successful deep GD models gives the same
impossibility, again with exact rational verification: at $H=32$ the
necessary output-change certificate lies between $0.24417$ and $0.32754$;
at $H=64$ it lies between $6.31046$ and $13.65547$.

This is an impossibility result for the \emph{global sufficient certificate},
not a violation of numerical stability itself. The unactivated suffixes
need not be derivatives at any actual Boolean input. The successful
pointwise certificates for these models use saturation, which the global
metric inequalities above omit. A positive global certificate for these
particular networks therefore needs additional information about the
reachable states, rather than merely a better choice of diagonal scales.
The exploratory source and exact rational records are kept separately
from the frozen experiment archive in
\texttt{diagonal-metric-exploration/}.

\section{Stability of a bounded prediction score}
\label{sec:clipped-stability}

Raw output stability may penalize changes between two very large outputs
that represent the same confident prediction. A bounded score avoids this.
Write $T(u)=\max\{-1,\min\{1,u\}\}$, fix $0<c\le1/100$, and put
\[
 r_{\rm clip}(w)=\frac{c}{\sqrt{1+(L-1)E_v(w)}}.
\]
Restrict the Gaussian to networks satisfying
\begin{equation}\label{eq:clipped-support}
 |T(z_w(x;\eta))-T(z_w(x))|\le1/16
 \quad\text{for all }x\text{ and }\|\eta\|_*\le r_{\rm clip}(w).
\end{equation}
Impose no all-input clean margin, and retain the training likelihood
$y_i z_w(x_i)>1/2$.

\begin{corollary}[Bounded-score stability]
\label{cor:clipped-stability}
This prior and likelihood satisfy the polynomial-depth conclusion of
Theorem~\ref{thm:numeric-prior}, including its test signed-output margin
$1/4$, every width above Boolean circuit width, and every $\gamma\ge1$.
The condition can also be imposed only on a fixed trusted region $G$,
at an added test-error cost $D(G^c)$.
\end{corollary}
\begin{proof}
Before rounding, the entropy compressor evaluates the original network at
a disturbance of norm less than the retained common radius divided by two.
Condition~\eqref{eq:clipped-support} bounds the resulting change of $T(z)$
by $1/16$. Since $T$ is $1$-Lipschitz, the grid-rounding error adds at most
$r/8$, so the total bounded-score error is at most $1/12$.
Apply the ramp proof to $T(h(x))$ instead of $h(x)$.
Oddness and monotonicity of $T$ ensure
\[
 yz>1/2\Longrightarrow yT(z)>1/2,\qquad
 yz\le1/4\Longrightarrow yT(z)\le1/4.
\]
Thus the same training and test implications, and the same finite-class
bound, hold.

For target mass, even the original robust digital cylinder of
Appendix~\ref{sec:cylinder} suffices: it has correct noisy signed output greater
than one at every input throughout per-layer RMS radius
$\tau_0=1/(64\sqrt{2C})$, and $E_v\ge C$.
Since $r_{\rm clip}\le(100\sqrt C)^{-1}<\tau_0$, its clean and noisy
scores are both exactly $f(x)$. Hence it satisfies the new support
condition with zero score change and retains the same Gaussian mass bound.
The entropy and readout-tail steps are unchanged. Masked ramp losses give
the trusted-region version.
\end{proof}

The raw numerical condition implies the clipped condition at the same
radius, but not conversely. Any finite-noise raw output enclosure
$z_w(x;\eta)\in[z-B,z+B]$ gives the sound clipped-score certificate
\[
 \max\{T(z+B)-T(z),\ T(z)-T(z-B)\}\le1/16.
\]
This is often much smaller than $B$ for large clean outputs. An attack
that changes only the raw magnitude while leaving both scores saturated
does not refute~\eqref{eq:clipped-support}.

\paragraph{Soft version.}
Let $\rho_{\rm clip}(w)$ be the smallest disturbance norm, minimized over
inputs, at which the clipped score changes by at least $1/16$.
Continuity makes it positive for every finite weight array; output noise
makes it finite. A density proportional to
\[
 \exp\{-(N/t^2)\lambda\min(1,\rho_{\rm clip}(w))^{-p}\}
\]
relative to the raw Gaussian therefore has full support. For fixed
$p,\lambda>0$, the soft-prior tail calculation and the bounded-score
compression argument give the same polynomial-depth theorem. After
truncating the penalty at $u$, use common radius
$\min\{1/6,\tfrac12u^{-1/p}\}$; this both lies strictly below the first
score-change radius and satisfies the rounding constants above.

\paragraph{Readout amplification.}
The bounded-score condition restores the scaling effect of calibrated
sign robustness. If a fixed finite-cube network has no clean zero outputs,
then sufficiently large positive readout amplification makes it satisfy
\eqref{eq:clipped-support}. Indeed there are a hidden-noise radius $\rho>0$
and signed unscaled output lower bound $m>0$ valid on the whole cube.
The calibrated radius tends to zero as the readout scale $A$ grows, while
$Am-r_{\rm clip}>1$ eventually. Both clean and noisy scores then equal
the clean sign. The required scale can be very large, and its Gaussian
readout-energy cost is $A^2E_v$; this observation supplies no uniform sample
bound by itself. The large-readout obstruction for raw numerical stability
in~\eqref{eq:readout-scaling-limit} concerns a different condition.

\section{Alternative priors and operational stability conditions}
\label{sec:alternative-priors}

This section gives two changes with complete statistical bridges. A soft
inverse-margin penalty retains the support of the raw Gaussian. A bounded
activation-fault condition weakens the previous hard support rule and has a
direct interpretation in terms of numerical errors and neuron failures.
Neither is a theorem for the \emph{unchanged} Gaussian prior.

\subsection{A soft prior with the same support as the raw Gaussian}

Let $\mu_{N,t}$ be the raw Gaussian of Section~\ref{sec:model}, $0<t\le1$,
and put $a=N/t^2\ge1$. Retain the postactivation disturbance convention
\[
 \|\eta\|_*^2=\sum_{s=1}^{L-1}\frac{\|\eta_s\|_2^2}{N}
                         +|\eta_{\rm out}|^2.
\]
Define the exact all-input distance to a zero output by
\[
 \rho(w)=\min_{x\in\{-1,1\}^d}
       \inf\{\|\eta\|_*:z_w(x;\eta)=0\},\qquad
 r(w)=\min\{1,\rho(w)\}.
\]
An output-only disturbance makes $\rho$ finite. If all clean outputs are
nonzero, continuity and finiteness of the input cube make $\rho>0$.
Every disturbance ball with radius strictly below $\rho$ preserves the clean
signs: a sign change would yield a zero along the segment from zero noise.
The function $\rho$ is measurable, since the noise balls are compact and
the input set is finite. Under the raw Gaussian, all clean outputs are
nonzero almost surely. To see this, condition successively on the preceding
layers: each nonzero state gives nondegenerate Gaussian preactivations,
and finally a nondegenerate Gaussian scalar readout. A finite union over
inputs completes the argument.

Fix constants $p>0$ and $\lambda>0$, and define
\begin{equation}\label{eq:soft-margin-prior}
 \Phi(w)=r(w)^{-p},\qquad
 d\nu_{N,t}(w)=Z^{-1}\exp[-a\lambda\Phi(w)]\,d\mu_{N,t}(w).
\end{equation}
On the null set $\rho=0$, set the density to zero. The normalizer is positive
and finite, and $\nu_{N,t}$ and $\mu_{N,t}$ are mutually absolutely
continuous. Thus no set of positive Gaussian probability is excluded.
In energy notation, the new density is proportional to
$\exp\{-a[\kappa_N(w)+2\lambda\Phi(w)]/2\}$.
One can optionally retain the original clean-margin restriction
$\min_x|z_w(x)|>1/2$; without that optional restriction, even the original
hard output-margin condition is unnecessary for the following result.
In both versions, the likelihood is noiseless \emph{sign consistency} with
the observed training labels.

\begin{theorem}[Soft robustness prior]
\label{thm:soft-margin-prior}
Suppose that, for a Boolean target $f$, there is a set $A_f$ of exact sign
representations such that
\[
 \mu_{N,t}(A_f)\ge e^{-aD},\qquad
 \rho(w)\ge r_*\quad(w\in A_f),
\]
where $D\ge1$ and $0<r_*\le1$ are independent of $N,t$.
For fixed $p,\lambda$, the consistency posterior under
\eqref{eq:soft-margin-prior} has a Gibbs-risk sample bound polynomial in
\[
 d,L,D,r_*^{-1},\varepsilon^{-1},\log(1/\delta).
\]
If the cylinder hypotheses hold at every $N\ge C$ and $0<t\le1$, the
guarantee holds on one sample event simultaneously for all those widths
and scales. The optional clean-margin version has the same conclusion
when $A_f$ also satisfies that clean margin.
\end{theorem}

\begin{proof}
Set $B=D+\lambda r_*^{-p}$. For every sample $S$ from $f$, its unnormalized
partition function satisfies
\[
 Z_S=\int \mathbf1_{\text{consistent with }S}
                 e^{-a\lambda\Phi}\,d\mu_{N,t}\ge e^{-aB}.
\]
The global prior normalizer cancels after conditioning. The posterior
$Q_S$ therefore obeys
\[
 D_{\rm KL}(Q_S\Vert\mu_{N,t})
  =-a\lambda\mathbb E_{Q_S}\Phi-\log Z_S\le aB,
 \qquad
 Q_S\{\Phi>u\}\le e^{a(B-\lambda u)}.
\]
Choose $0<\beta\le1/2$ and
$u=[B+\log(1/\beta)]/\lambda\ge1$. Since $a\ge1$, the displayed tail is
at most $\beta$. Let $Q_S'$ be $Q_S$ conditioned on $\Phi\le u$. Its
explicit density gives
\[
 D_{\rm KL}(Q_S'\Vert\mu_{N,t})
  \le aB+\log\frac1{1-\beta}\le a(B+1).
\]
Its support is strictly sign-robust at radius
$r_0=\tfrac12u^{-1/p}$. Apply the uniform robust posterior-compression
theorem in Appendix~\ref{sec:robust-entropy} with $K=B+1$ and $r=r_0$. It gives one
finite class $\mathcal H$, with polynomial log cardinality, and a common
candidate marginal whose prediction can be coupled to $Q_S'$ at each
input with disagreement at most $\alpha$.
Since $Q_S'$ is training-consistent, the finite-class relative Chernoff
bound from that report gives
\[
 R(Q_S)\le\beta+3\alpha+
       \frac{8[\log|\mathcal H|+\log(1/\delta)]}{n}
\]
with probability at least $1-\delta$. Choose $\beta,\alpha$ proportional
to $\varepsilon$. All class parameters are independent of $N,t$, which
also proves simultaneity. The same calculation with the extra clean-margin
indicator proves the optional version.
\end{proof}

The robust digital cylinder of Appendix~\ref{sec:cylinder} has
$D=\operatorname{poly}(d,L,C)$ and $r_*^{-1}=\operatorname{poly}(C)$.
Consequently Theorem~\ref{thm:soft-margin-prior} gives a polynomial-depth
bound for every $N\ge C$ and every $\gamma\ge1$ while retaining the support
of the raw Gaussian. The prior probabilities have changed. The result
therefore does not settle whether ordinary independent Gaussian weights
alone give the same guarantee.
An inverse-squared margin ($p=2$) is a particularly simple choice.
A penalty merely logarithmic in $1/\rho$ would only produce an exponentially
small retained radius through this proof, so would not give the same
polynomial conclusion.

\subsection{Robustness only within the activation range}

Call a disturbance \emph{bounded-state} if every perturbed hidden activation
belongs to $[-1,1]^N$. This allows replacing an activation by another legal
activation value, and leaves scalar output noise unrestricted. Require
sign stability only against bounded-state disturbances in the specified
RMS ball.

\begin{proposition}[Bounded-state compression]
\label{prop:bounded-state-compression}
The deterministic and posterior compression constructions in Appendix~\ref{sec:robust-entropy} remain valid under bounded-state robustness in place of unrestricted
postactivation robustness. The previously established target Gaussian-mass
estimates remain valid.
\end{proposition}

\begin{proof}
On every successful compression event, every virtual hidden state is a
$\tanh$ output. First-layer clipping occurs before the activation;
active-coordinate clipping, Gaussian replacement, and source reconstruction
also feed the activation. The exact postactivation errors comparing the
virtual computation to the original weights therefore produce states in
$[-1,1]^N$, and have the same total RMS bound as before. The margin argument
adds only output noise. Subsequent finite-grid rounding uses that positive
output margin and does not invoke robustness again. The small-width case
also uses only output noise. These are all uses of robustness in the two
proofs. Finally, the previous target cylinder satisfies the stronger,
unrestricted condition, so its probability lower bound is unchanged.
\end{proof}

This relaxation can be substantial at a fixed radius. With one hidden
layer, let $n=N-1$, $c=0.9$, and take a constant hidden state
$h=(c,\ldots,c,-c)$ and readout
$v=(2/(nc),\ldots,2/(nc),1/c)$. Its output is $1$.
Unrestricted hidden and output noise can erase it at norm
$[N\|v\|^2+1]^{-1/2}=O(N^{-1/2})$.
For bounded-state noise, the last coordinate can decrease by at most
$1-c$. The remaining coordinates and scalar output noise can decrease the
output by at most
$\sqrt{4N/(nc^2)+1}\,\|\eta\|_*$.
Thus a zero requires
\[
 \|\eta\|_*\ge\frac{2c-1}{\sqrt{4N/n+c^2}}>0.26
 \qquad(N\ge2).
\]
This separates the notions of robustness; its readout energy grows with
$N$, so it is not a sample lower bound or a separation of the calibrated
prior supports.

\subsection{Small numerical errors and rare neuron faults}

For $a>0$, $0\le p\le1$, and $b\ge0$, say a network tolerates $(a,p,b)$ if it preserves its own clean signs at
every input under all simultaneous bounded-state disturbances such that,
in each hidden layer, at most $pN$ coordinates have $|\eta_i|>a$, and
$|\eta_{\rm out}|\le b$. Exceptional coordinates may be replaced
arbitrarily within $[-1,1]$. This combines numerical errors at all neurons
with an adversarial fraction of neuron faults. Dropout and pruning are
particular faults; average performance under random dropout is a weaker
property than this adversarial condition.

\begin{lemma}[Fault and RMS bridges]
\label{lem:fault-rms}
Tolerance of $(a,p,b)$ implies bounded-state RMS robustness at
$r=\min\{a\sqrt p,b\}$. Conversely, unrestricted RMS robustness at $\tau$
implies tolerance of $(a,p,b)$ whenever
\[
 (L-1)(a+2\sqrt p)^2+b^2\le\tau^2.
\]
\end{lemma}
\begin{proof}
If the total RMS error is at most $r$, each layer has at most
$Nr^2/a^2\le pN$ coordinates with error larger than $a$; output noise is
at most $r\le b$. Conversely, the ordinary coordinates contribute at most
$a\sqrt N$ to a layer's Euclidean error. Each faulty coordinate differs
by at most $2$, so their contribution is at most $2\sqrt{pN}$.
Sum the squared layer RMS bounds and the scalar output bound.
\end{proof}

For a single rule independent of target and data, retain the original
clean margin and define
\begin{align*}
 E_v&=N\|v\|^2,&
 \tau(w)&=\frac{0.01}{\sqrt{1+(L-1)E_v}},\\
 a(w)&=\frac{\tau(w)}{4\sqrt L},&
 p(w)&=\frac{\tau(w)^2}{64L},&
 b(w)&=\frac{\tau(w)}2.
\end{align*}
Restrict the Gaussian to networks tolerating these faults, then condition
on noiseless training sign consistency. Every network in the previous
radius-$\tau(w)$ support belongs to this support: the second bound in
Lemma~\ref{lem:fault-rms} is at most $\tau(w)^2/2$.
In particular, the previous target cylinder still lies in the support.

\begin{corollary}[Polynomial-depth learning with activation faults]
\label{cor:fault-prior}
For this fixed fault-tolerance support rule, the same qualitative
polynomial-depth sample theorem holds for all $\gamma\ge1$ and every
width above the layered Boolean target circuit's width.
\end{corollary}
\begin{proof}
The same target partition bound and Gaussian readout tail as before allow
conditioning on $E_v\le B$ while discarding posterior mass at most
$\beta$, where $B=\operatorname{poly}(d,L,C,\log(1/\beta))$.
Every retained network has bounded-state robustness radius at least
\[
 r_B=\frac{10^{-4}}{32L[1+(L-1)B]},
\]
because $a(w)\sqrt{p(w)}=\tau(w)^2/(32L)\le b(w)$.
The conditional posterior has the same polynomial scaled-KL bound as in
the previous argument. Apply Proposition~\ref{prop:bounded-state-compression}
and the finite-class generalization step. All parameters are independent
of width and Gaussian scale, giving the simultaneous conclusion.
\end{proof}

This condition has an operational interpretation, but it is not established
as a universal outcome of ordinary gradient descent. Likewise, random-noise
stability alone cannot simply replace it. A network that reads a single
hidden coordinate can be stable under small independent per-coordinate
Gaussian noise while that coordinate can be erased with only
$O(N^{-1/2})$ RMS adversarial noise. Testing particular dropout masks or
one quantization grid similarly does not control the disturbances used by
the compression coupling.
The conservative constants may give $p(w)N<1$ at small widths; then the
allowed number of arbitrary faults is zero, and the condition concerns
bounded errors at all coordinates. No claim of tolerance to a practically
large dropout fraction follows from these constants.

\section{A trusted input region retains the original sign likelihood}
\label{sec:trusted-region}

The distributional robustness restriction can be combined with the original
sign-only likelihood if the exceptional region is shared across hypotheses.
This is relevant when non-data-like inputs, rather than ordinary test inputs,
are responsible for a failed all-cube certificate.

Fix an input region $G\subseteq\{-1,1\}^d$ before the labeled sample is
drawn. Suppose the test input distribution satisfies
$D(G^c)\le\rho$. The region may be selected using independent unlabeled
information; its coverage is an assumption to be justified separately.
It is independent of the unknown labels. Set
\[
 E_v(w)=N\|v\|^2,\qquad
 r_{\rm ref}(w)=\frac1{100\sqrt{1+(L-1)E_v(w)}}.
\]
Restrict the raw Gaussian only to networks whose own clean sign is preserved
under every disturbance of total RMS norm at most $r_{\rm ref}(w)$ at
every input in $G$. Impose no margin or robustness condition outside $G$,
and no additional global clean margin. Then condition on \emph{ordinary
sign consistency at every observed training example}, including examples
outside $G$.

\begin{theorem}[Trusted-region prior, original likelihood]
\label{thm:trusted-region}
For a target represented by a width-$C$ layered Boolean circuit, this
posterior has Gibbs error at most $\rho+\varepsilon$ after
\[
 n\ge\operatorname{poly}(d,L,C,\varepsilon^{-1},\log(1/\delta))
\]
examples, with probability at least $1-\delta$, simultaneously for every
$N\ge C$ and $0<t\le1$. The polynomial's degree is independent of width,
depth, and Gaussian scale. The likelihood is the original sign-only
likelihood, rather than a robust training likelihood.
\end{theorem}

\begin{proof}
The earlier global robust target cylinder belongs to the new support and
satisfies ordinary training consistency for every sample. Its raw Gaussian
mass is at least $\exp(-ND_*/t^2)$, where
$D_*=\operatorname{poly}(d,L,C)$. Therefore the posterior has scaled KL
at most $ND_*/t^2$. The same Gaussian readout tail bound permits discarding
posterior mass at most $\beta$ while retaining $E_v\le B$, where
$B=\operatorname{poly}(D_*,\log(1/\beta))$.
The retained law $Q'$ has KL at most $N(D_*+1)/t^2$ and is robust on $G$
at the common radius $r_0=[100\sqrt{1+(L-1)B}]^{-1}$.

The pointwise-defect compression proposition gives a common law $\Gamma$
on a finite class $\mathcal H$ and a prediction coupling with disagreement
at most $\alpha$ at each $x\in G$. No prediction approximation is needed
outside $G$. Apply the finite-class relative Chernoff bound to the
\emph{masked losses}
\[
 \ell_G(h,x)=\mathbf1_G(x)\mathbf1\{h(x)\ne f(x)\}.
\]
With probability at least $1-\delta$, simultaneously for every
$h\in\mathcal H$,
\[
 \mathbb E_D\ell_G(h,X)
 \le\frac2n\sum_{i=1}^n\ell_G(h,x_i)
       +\frac{8[\log|\mathcal H|+\log(1/\delta)]}{n}.
\]
The law $Q'$ is training-consistent, so averaging the empirical loss under
$\Gamma$ gives at most $\alpha$. The population comparison adds at most
another $\alpha$. Hence
\[
 \mathbb E_{w\sim Q',X\sim D}
    \mathbf1_G(X)\mathbf1\{\operatorname{sign}z_w(X)\ne f(X)\}
 \le3\alpha+
       \frac{8[\log|\mathcal H|+\log(1/\delta)]}{n}.
\]
Outside $G$ the error costs at most $\rho$, and the discarded posterior
mass costs at most $\beta$. Choosing $\alpha,\beta$ proportional to
$\varepsilon$ proves the theorem. The finite class is independent of
$N,t$, so its one sample event proves simultaneity.
\end{proof}

The fixed exceptional region is crucial. This theorem does not allow each
hypothesis to select its own exceptional inputs after seeing the sample.
It also does not justify choosing $G$ from that same labeled sample without
additional complexity control. A finite collection of independently
obtained unlabeled inputs need not cover most of $D$ merely because all its
points are certified. Coverage and certification over $G$ are separate
requirements.

The soft-prior construction also localizes to $G$. Replace the global
margin in Theorem~\ref{thm:soft-margin-prior} by
\[
 \rho_G(w)=\min_{x\in G}\inf\{\|\eta\|_*:z_w(x;\eta)=0\}.
\]
For nonempty $G$, this is positive and finite Gaussian-almost surely.
The soft density proportional to
$\exp\{-(N/t^2)\lambda\min(1,\rho_G(w))^{-p}\}$ therefore retains the
support of the raw Gaussian and does not penalize instability outside $G$.
The same target cylinder and posterior truncation as in
Theorem~\ref{thm:soft-margin-prior}, followed by the masked-loss proof above,
give error at most $D(G^c)+\varepsilon$ with the original sign likelihood
and polynomial sample complexity. This combines a soft prior with
restriction of the stability requirement to a trusted data region.

\begin{corollary}[A finite catalog of trusted regions]
\label{cor:trusted-catalog}
Let $G_1,\ldots,G_M$ be fixed before the labeled sample, with
$D(G_j^c)\le\rho$ for every $j$. Choose any prior distribution on the
index $j$, and conditional on $j$ use the Gaussian restricted to robustness
on $G_j$ as above. Condition the resulting joint prior on the original
training signs. Then the marginal network posterior satisfies the same
error guarantee with polynomial sample complexity also in $\log M$.
\end{corollary}
\begin{proof}
Conditional on any index $j$ and the observed sample, the network law is
the raw Gaussian conditioned on robustness on $G_j$ and sign consistency.
The global robust target cylinder belongs to every such conditioning event,
so its same raw mass lower bound applies separately to every component.
Consequently every component has the same entropy bound, readout tail
bound, and common finite compression class $\mathcal H$.
Apply the preceding masked-loss proof simultaneously to the $M|\mathcal H|$
loss functions $\mathbf1_{G_j}(x)\mathbf1\{h(x)\ne f(x)\}$.
Only $\log M$ is added to the finite-class bound. Every component posterior
then has error at most $\rho+\varepsilon$ on the same sample event, so
averaging under the posterior law of $j$ preserves the guarantee.
\end{proof}

Thus the exceptional region may adapt to the labels through posterior
selection from a predeclared catalog, with an explicit complexity cost.
The prior here is a mixture of region-conditioned Gaussians; it is not
silently identified with the Gaussian conditioned on their union.

\subsection{Why a per-hypothesis distributional exception is insufficient}
\label{sec:patch-obstruction}

The following abstract example shows why the compression-with-defect and
target-mass facts alone cannot remove the empirical defect term under
ordinary sign conditioning. It is \emph{not} a lower bound for the exact
Gaussian neural prior.

Let the input domain have $M$ points, uniformly distributed, and let the
target be $f\equiv+1$. Fix an integer $q$ with $\rho=q/M<1/2$. For each
size-$q$ set $S$, define the patch predictor
\[
 h_S(x)=\begin{cases}+1,&x\in S,\\-1,&x\notin S,\end{cases}
 \qquad b_{h_S}(x)=\mathbf1_S(x).
\]
Include the target predictor $h_+\equiv+1$ with bad-set indicator zero.
Every predictor has population bad-set mass at most $\rho$.
The two-element class $\{h_+,h_-\}$, where $h_-\equiv-1$, compresses these
predictors exactly outside their bad sets: map $h_S$ to $h_-$ and $h_+$
to itself. Thus the pointwise prediction-compression conclusion holds with
zero approximation error and $\log|\mathcal H|=\log2$.

Give $h_+$ prior mass $\eta=e^{-aD_*}$ and distribute the remaining mass
uniformly over the patch predictors. Suppose $\eta\le1/2$, $q\ge2n$, and
\[
 n\le\frac{aD_*}{2\log(2/\rho)}.
\]
If the sample has $k\le n$ distinct inputs, the patch component's probability
of consistency is
\[
 L_S=\frac{(q)_k}{(M)_k}\ge(\rho/2)^n\ge\sqrt\eta.
\]
Here $(u)_k=u(u-1)\cdots(u-k+1)$. Consequently the posterior mass on the
target is at most
$\eta/[(1-\eta)L_S]\le2\sqrt\eta$, and its Gibbs error is at least
\[
 (1-\rho)(1-2e^{-aD_*/2}).
\]
The posterior still has population bad-set mass at most $\rho$, but its
empirical bad-set mass is at least $1-2\sqrt\eta$: every training point
belongs to every consistent patch's bad set. Its KL to this abstract prior
is at most $\log(1/\eta)=aD_*$, because conditioning has a partition
function at least $\eta$.

Thus a target-mass exponent of size $aD_*$, bounded population defect,
scaled KL control, and a tiny compressed class do not by themselves yield
a bound independent of $a$. In our neural setting $a=N/t^2$, exactly the
quantity that the robust compression argument must eliminate. The example
isolates a logical obstruction; showing that the specified Gaussian prior
does or does not exhibit it would require additional analysis.

\section{Choosing and validating a trusted region from unlabeled data}
\label{sec:unlabeled-trusted-region}

A trusted region need not be specified before seeing any data. Independent
unlabeled sample splitting supplies a simple rigorous version of this idea.
This is a standard statistical construction, not a new sample-splitting
principle. Its role here is to make the coverage assumption in the preceding
robust-prior results independently checkable.

Let $D$ be the test input distribution on $\{-1,1\}^d$. Draw three mutually
independent samples: an unlabeled pilot $P$, an unlabeled validation sample
$V=(X'_1,\ldots,X'_m)$, and the labeled training sample
$S=((X_i,f(X_i)))_{i=1}^n$. Using only $P$, choose any measurable region
$G=G(P)$. The selection rule can be arbitrarily complicated. For example,
$P$ may determine a density level set, a union of neighborhoods of observed
inputs, or a support estimate in a fixed representation. The theorem does
not assert that such regions admit an easy all-input stability certificate.
Define
\[
 \widehat\rho_G=\frac1m\sum_{j=1}^m\mathbf1\{X'_j\notin G\},\qquad
 \overline\rho_G=\min\!\left\{1,\widehat\rho_G+
            \sqrt{\frac{\log(1/\delta_u)}{2m}}\right\}.
\]

\begin{theorem}[An independently learned trusted region]
\label{thm:unlabeled-trusted-region}
Fix $0<\delta_u,\delta_l<1$ and $0<\varepsilon<1$.
Suppose the target has a width-$C$ layered Boolean circuit. For the selected
$G(P)$, form any one of the following priors by conditioning the raw Gaussian
on the indicated condition at every $x\in G(P)$:
\begin{enumerate}
\item The calibrated sign stability condition of
Theorem~\ref{thm:trusted-region}, followed by ordinary training sign
consistency. Let $R(Q)$ denote Gibbs sign error.
\item The numerical stability condition of
Theorem~\ref{thm:numeric-prior}, with signed training outputs greater
than $1/2$. Let $R(Q)$ denote
$\Pr_{w\sim Q,X\sim D}\{f(X)z_w(X)\le1/4\}$.
\item The clipped-score stability condition of
Corollary~\ref{cor:clipped-stability}, with the same training-margin
likelihood and the same test loss as in item 2.
\end{enumerate}
No stability condition is imposed outside $G(P)$, and no additional
all-input clean-margin restriction is imposed. For the clipped condition, fix $0<c\le1/100$ once and for all;
constants in the following bound may depend on $c$. The numerical
condition retains its stated fixed calibration. There is a
polynomial sample bound
\[
 n\ge\operatorname{poly}
       (d,L,C,\varepsilon^{-1},\log(1/\delta_l))
\]
such that, with probability at least $1-\delta_u-\delta_l$ over $P,V,S$,
\[
 R(Q)\le\overline\rho_G+\varepsilon,
\]
simultaneously for all widths $N\ge C$ and scales $0<t\le1$.
The complexity of the pilot selection rule does not appear in this bound.
The Gaussian is conditioned exactly on the selected support condition;
no robust training likelihood or soft reweighting is introduced.
\end{theorem}

\begin{proof}
Condition on the pilot. The set $G$ is now fixed and independent of both
$V$ and $S$. Hoeffding's inequality implies
$D(G^c)\le\overline\rho_G$ with probability at least $1-\delta_u$ over
$V$. The corresponding fixed-region theorem gives
$R(Q)\le D(G^c)+\varepsilon$ with probability at least $1-\delta_l$ over
$S$. Its constants and common finite compression class are uniform over
$G$: in particular, the globally stable target cylinder belongs to the
conditioned support for every possible $G$. A union bound, followed by
averaging over the pilot, proves the result. The sample event supplied by
the fixed-region theorem is simultaneous in width and Gaussian scale.
\end{proof}

For any of these losses, a pointwise posterior correct-answer probability
of at most $2/3$ entails pointwise Gibbs error at least $1/3$.
Consequently the measure of such inputs is at most
$3(\overline\rho_G+\varepsilon)$. For items 2 and 3, ``correct answer''
can mean the stronger event $f(x)z_w(x)>1/4$.
An exact one-sided binomial upper confidence limit can replace Hoeffding.
For example, when the validation sample has no misses, the upper limit is
$1-\delta_u^{1/m}\le\log(1/\delta_u)/m$.

\paragraph{Validation candidates.}
The pilot may produce a catalog of $M$ regions, denoted $G_j(P)$.
Using
\[
 \overline\rho_j=\min\!\left\{1,\widehat\rho_j+
              \sqrt{\frac{\log(M/\delta_u)}{2m}}\right\}
\]
certifies all $M$ coverages simultaneously. One may then choose any single
index $J=J(P,V)$ before seeing $S$, for example the index with the best
coverage among regions admitting a convenient certificate. The theorem
holds with $G_J$ and $\overline\rho_J$. There is no extra $\log M$ cost in
the labeled sample bound: conditional on $P,V$, just one region has been
chosen independently of $S$. The $\log M$ cost is paid in the unlabeled
validation confidence bound. If the region index is instead chosen using
$S$, the additional complexity control of the finite-catalog theorem is
needed. Reusing validation to generate an unbounded number of new regions
also requires further control; the single-region guarantee cannot simply
be applied after that adaptation.

\paragraph{What is still required.}
Coverage validation checks only how often an input lies in $G$. The prior
must still impose stability throughout $G$, simultaneously for every
admissible disturbance. Certifying stability merely on the finite pilot or
validation inputs is not sufficient. In particular, high average certified
robustness of one trained network does not by itself prove a theorem about
the Bayesian posterior. A network-dependent exceptional region can adapt
to the training labels, which is precisely the difficulty exhibited by the
sparse-patch example. The present construction avoids that difficulty by
freezing a shared region before the labeled sample. It modifies the support
of the Gaussian prior; it does not establish the same conclusion for the
unrestricted Gaussian prior.

\section{Network-dependent exceptions with controlled input-set complexity}
\label{sec:low-complexity-exceptions}

The exceptional inputs need not be shared by all networks if the family of
permitted exceptional sets has controlled complexity. This provides a
second way to retain the ordinary training likelihood. The restriction is
on where instability may occur, rather than on an altered likelihood.

Fix a family $\mathcal B$ of subsets of the Boolean cube, including the
empty set. Write $M$ for its number of distinct sets and suppose
$\log M\le H_B$. A family learned from an independent unlabeled pilot is
also allowed: all definitions below are then conditional on that pilot,
and $H_B$ must be a uniform bound. Draw an independent unlabeled sample
$V$ of size $m$ and a labeled sample $S$ of size $n$. Fix $0\le\rho<1$.
For a network $w$, permit any exceptional set $B\in\mathcal B$ satisfying
\[
 \widehat D_V(B)\le\rho.
\]
Condition the raw Gaussian on the existence of such a $B$ for which the
calibrated sign-stability condition holds at every $x\notin B$.
Here the radius is $[100\sqrt{1+(L-1)N\|v\|^2}]^{-1}$, as in the
trusted-region theorem. Then condition on ordinary training sign
consistency at every example, including examples in the exceptional set.
The exceptional set may depend on the network and hence on the labels
through the posterior. This prior is the Gaussian conditioned on a union
of admissible support events; it is not a mixture of separately normalized
conditioned Gaussians.

\begin{theorem}[Low-complexity distributional exceptions]
\label{thm:low-complexity-exceptions}
For every $0<\varepsilon,\delta<1$ there are sample bounds
\[
 m\ge O\!\left(\frac{H_B+\log(1/\delta)}{\varepsilon^2}\right),
 \qquad
 n\ge\operatorname{poly}
       (d,L,C,H_B,\varepsilon^{-1},\log(1/\delta))
\]
such that a width-$C$ layered Boolean target circuit has posterior Gibbs
sign error at most $3\rho+\varepsilon$, with probability at least
$1-\delta$. The guarantee holds simultaneously for every $N\ge C$ and
$0<t\le1$. There is no all-input clean-margin assumption.

The same result holds when sign stability is replaced by numerical or
clipped-score stability, the training likelihood is
$f(X_i)z_w(X_i)>1/2$, and the test loss is
$\mathbf1\{f(X)z_w(X)\le1/4\}$. The numerical radius is unchanged.
In the clipped case fix $0<c\le1/100$ once and for all, with bound
constants allowed to depend on $c$.
\end{theorem}

\begin{proof}
Set $u=\varepsilon/20$. Uniform one-sided Hoeffding bounds, one for each
sample, imply that, except on events of probability $\delta/3$ each,
\begin{align*}
 D(B)&\le\widehat D_V(B)+u&&\text{for every }B\in\mathcal B,\\
 \widehat D_S(B)&\le D(B)+u&&\text{for every }B\in\mathcal B,
\end{align*}
provided $m,n\ge[H_B+\log(3/\delta)]/(2u^2)$.
Every admissible exceptional set therefore has population mass at most
$\rho+u$ and empirical training mass at most $\rho+2u$.
These statements hold uniformly and hence also for sets chosen by a
posterior network.

The global robust target cylinder qualifies using $B=\varnothing$, for
every realization of $V$. It has the same raw Gaussian mass lower bound
$\exp(-ND_*/t^2)$ as before and is consistent with all labels. Consequently
the full posterior---including the union in its support event---has
scaled KL at most $ND_*/t^2$. The usual readout-energy truncation discards
mass at most $\beta$ and leaves a law $Q'$ with a common inverse-polynomial
stability radius and the same polynomial scaled entropy bound.

For each retained network choose one witnessing exceptional set; the
selection can be measurable because there are only finitely many distinct
sets. Let $b_D$ and $b_S$ denote its population and empirical masses,
averaged under $Q'$. The preceding uniform events give
$b_D\le\rho+u$ and $b_S\le\rho+2u$.
The pointwise-defect compression argument and its finite-class relative
Chernoff bound give, on one further event of probability at least
$1-\delta/3$,
\[
 R(Q)\le\beta+3\alpha+b_D+2b_S+
       \frac{8[\log|\mathcal H|+\log(3/\delta)]}{n}
 \le3\rho+5u+\beta+3\alpha+
       \frac{8[\log|\mathcal H|+\log(3/\delta)]}{n}.
\]
Take $\beta=\varepsilon/8$, $\alpha=\varepsilon/24$, and choose $n$ so
the last term is at most $\varepsilon/2$. The numerical and clipped cases
use the same bounded ramp loss as before; an exceptional input contributes
at most one to its population or empirical comparison, giving the same
inequality. The common finite class is independent of $N,t$, and the
uniform set-mass events do not involve either parameter. This proves the
simultaneous statement.
\end{proof}

\paragraph{A VC-dimension formulation.}
If $\mathcal B$ has VC dimension at most $h\ge1$, its restriction to the
$2^d$ Boolean inputs has at most
\[
 \sum_{j=0}^{\min(h,2^d)}\binom{2^d}{j}
 \le(2^d+1)^h
\]
distinct sets. Thus $H_B=h\log(2^d+1)\le h(d+1)\log2$ suffices.
No finiteness of a continuous parameterization is required. The theorem
therefore gives polynomial depth bounds whenever the exceptional-set VC
dimension is polynomial in the stated parameters and independent of
network width. This use of the finite Boolean domain is intentionally
simple; ordinary VC uniform-convergence estimates can also be used.

For example, one can permit exceptional sets that are halfspaces or balls
in a fixed $p$-dimensional input representation, or unions of a bounded
number of such sets. The corresponding VC dimensions are polynomial in
$p$ and the number of components. The representation may be learned from
the independent pilot. These are candidate restrictions, not a claim that
the actual failure sets of unconstrained trained networks have small VC
dimension. Whole-region stability outside a witnessing exceptional set
must still hold. A certificate proving that all failures lie inside such
a region is enough; high empirical robustness alone is not.

The theorem gives a genuine distributional relaxation with
network-dependent exceptions and the original sign likelihood. It still
changes the Gaussian prior's support. It neither proves an unrestricted
Gaussian result nor evades the sparse-patch obstruction without an added
assumption: the family of arbitrary sparse patches has large input-set
complexity, precisely what is controlled here.



\section{The remaining research questions}

The general unrestricted theorem still carries
$Q_*=[(K+2)/L]^{L/2}$. The shallow-target construction bounds this
quantity polynomially for several substantial target classes. This report does not remove
its exponential dependence for general small deep circuits. The later
conditional obstruction cited above rules out a universal joint-polynomial
bound under a circuit-complexity assumption, but leaves the rate of
sample-complexity growth unresolved.
At $\gamma=1$, the typical Gaussian bulk has order-one operator norm,
so the small-bulk proof does not apply. The finite-width result and the
robustness-based results cover this exponent under their stated hypotheses.
No new sample lower bound for the exact unrestricted prior is established.

The empirical question is also sharper than universal robustness of every
trained network. One should test finite-radius numerical stability on
relevant data and on a justified region around that data, retaining
optimization failures and boundary cases. Testing one trained network
does not establish a property of the Bayesian posterior. Testing finitely
many inputs does not establish a worst-case condition on an exponentially
large cube. The trusted-region and posterior-defect statements show which
additional facts would turn such evidence into a learning guarantee.

For designing a prior, stability of the numerical output or of a clipped
prediction score, with positive training margins, is a direct candidate: it expresses resistance to computation
errors without requiring every unobserved input to have a confident label.
The clipped-score version additionally ignores changes that leave the
prediction confidently saturated. Diagonal-metric certificates permit
compensating expansion and contraction between layers, rather than requiring
every layer to be nonexpansive in Euclidean norm.
A soft version retains the Gaussian support, while a trusted-region version
avoids demanding stability on arbitrary non-data-like inputs. Both change
the learning rule in explicit, measurable ways. Whether ordinary gradient
descent selects networks satisfying an appropriate version with high
probability is not proved by the present finite experiments or by the
counterexamples to a universal assertion.

\appendix
\section{An exact variational principle at the critical exponent}
\label{sec:critical-variational}

For completeness we reproduce the finite-domain variational foundation.
Here $\phi=\tanh$. For arbitrary fixed depth, the critical prior admits a direct finite-domain
large-deviation description. This result requires neither a propagation-of-chaos
assumption nor differentiability of the activation. It identifies the functional used in the critical-rate capacity constructions
with the actual Gaussian probability rate.
Covariance-process large deviations for Gaussian networks have been studied
in a more general functional setting by Andreis, Bassetti, and
Hirsch~\cite{andreis2026}. No novelty claim is made for the general
large-deviation framework; the finite-domain proof below fixes the precise
normalizations and strict-margin convention used here.

Fix $\gamma=1$, enumerate $\mathcal X=\{x^1,\ldots,x^m\}$, where $m=2^d$,
and put $r=L-1$. All limits in this section hold with $d$ and $L$ fixed.
Write $E=[-1,1]^m$ and let $\mathcal P(E)$ carry its weak topology.
For $\nu\in\mathcal P(E)$ define
\[
 K(\nu)=\int_E hh^{\mathsf T}\,d\nu(h),\qquad
 Q_0(a,b)=\frac{\bar x^a\cdot\bar x^b}{d}.
\]
For any positive semidefinite $m\times m$ matrix $Q$, let
\[
 \mathsf T_Q(u)=\bigl(\phi((Q^{1/2}u)_a)\bigr)_{a=1}^m,
 \qquad
 \mathsf M_Q=(\mathsf T_Q)_\#\mathcal N(0,I_m).
\]
Here $T_\#\rho$ is the distribution of $T(U)$ when $U$ has distribution $\rho$,
and $Q^{1/2}$ is the positive semidefinite square root. Relative entropy is
\[
 D(\nu\Vert\rho)=
 \begin{cases}
 \displaystyle\int\log\frac{d\nu}{d\rho}\,d\nu,&\nu\ll\rho,\\
 +\infty,&\text{otherwise}.
 \end{cases}
\]
Finally define the Gaussian energy
\[
 \mathcal G_Q(z)=
 \begin{cases}
 \frac12 z^{\mathsf T}Q^+z,&z\in\operatorname{range}(Q),\\
 +\infty,&z\notin\operatorname{range}(Q),
 \end{cases}
\]
where $Q^+$ is the Moore--Penrose inverse.

\begin{theorem}[Critical finite-domain large deviations]
\label{thm:critical-ldp}
Let $Z_N=(z_w(x^a))_{a=1}^m$ under the unconditioned prior
$\mu_{N,1}$. The sequence $Z_N$ satisfies a large-deviation principle
with speed $N$ and good rate function
\begin{equation}\label{eq:critical-rate}
 \begin{split}
 I(z)=\inf_{\nu_1,\ldots,\nu_r\in\mathcal P(E)}
 \bigg\{&D(\nu_1\Vert\mathsf M_{Q_0})
  +\sum_{s=2}^{r}D\bigl(\nu_s\Vert\mathsf M_{K(\nu_{s-1})}\bigr)
  +\mathcal G_{K(\nu_r)}(z)\bigg\}.
 \end{split}
\end{equation}
In particular, for every open $O\subseteq\mathbb R^m$ and closed
$F\subseteq\mathbb R^m$,
\[
 -\inf_{z\in O}I(z)
 \leq\liminf_{N\to\infty}\frac1N\log\Prb(Z_N\in O),
 \qquad
 \limsup_{N\to\infty}\frac1N\log\Prb(Z_N\in F)
 \leq-\inf_{z\in F}I(z).
\]
The infimum in \eqref{eq:critical-rate} is attained whenever it is finite.
No covariance matrix in the statement is required to be nonsingular.
\end{theorem}

\begin{proof}
We give an exact representation of the finite-width law and then apply
Sanov's theorem and continuous contraction~\cite{dembozeitouni}.

Let $\Gamma_m=\mathcal N(0,I_m)$. Independently for each
$s=1,\ldots,r$, let $U_{s,1},\ldots,U_{s,N}$ have law $\Gamma_m$, and set
\[
 \eta_{s,N}=\frac1N\sum_{i=1}^N\delta_{U_{s,i}}.
\]
Also let $V_N=G/\sqrt N$, where $G\sim\Gamma_m$ is independent of these
empirical measures. Recursively define
\[
 \nu_{1,N}=(\mathsf T_{Q_0})_\#\eta_{1,N},\qquad
 \nu_{s,N}=(\mathsf T_{K(\nu_{s-1,N})})_\#\eta_{s,N}\quad(2\leq s\leq r).
\]
These are exactly the joint laws of the network's empirical hidden-activation
vectors on $\mathcal X$. Indeed, first-layer preactivation vectors are
independent Gaussians with covariance $Q_0$. Conditional on the preceding
layer, each subsequent row of weights gives an independent Gaussian vector
with covariance
\[
 \frac1N\sum_{i=1}^N
 \bigl(h_i^{s-1}(x^a)h_i^{s-1}(x^b)\bigr)_{a,b}
 =K(\nu_{s-1,N}).
\]
The conditional distribution depends on the previous layer only through its
empirical measure. Finally, readout weights have variance $N^{-2}$, so
\begin{equation}\label{eq:critical-output-representation}
 Z_N\ \stackrel{\mathrm{law}}{=}\ K(\nu_{r,N})^{1/2}V_N.
\end{equation}
This representation is exact for every $N$.

Sanov's theorem on $\mathcal P(\mathbb R^m)$ gives the good rate
$D(\eta\Vert\Gamma_m)$ for each $\eta_{s,N}$. The elementary Gaussian
large-deviation principle gives rate $\|v\|_2^2/2$ for $V_N$.
Independence and the finite-product principle therefore give the good rate
\[
 \sum_{s=1}^r D(\eta_s\Vert\Gamma_m)+\frac12\|v\|_2^2
\]
for $(\eta_{1,N},\ldots,\eta_{r,N},V_N)$.

The recursive map from these variables to $(\nu_1,\ldots,\nu_r,z)$ is
continuous. To verify the only potentially delicate point, suppose
$Q_j\to Q$ and $\eta_j\Rightarrow\eta$ weakly. For every continuous
$b:E\to\mathbb R$, the functions $b\circ\mathsf T_{Q_j}$ converge
uniformly on each compact subset of $\mathbb R^m$ to
$b\circ\mathsf T_Q$, since matrix square roots are continuous and $\phi$
is continuous. Their absolute values are uniformly bounded. Tightness of a
weakly convergent sequence of probability measures then gives
\[
 \int b\circ\mathsf T_{Q_j}\,d\eta_j
 \longrightarrow\int b\circ\mathsf T_Q\,d\eta.
\]
Also $\nu\mapsto K(\nu)$ is continuous because every coordinate
$h_ah_b$ is continuous and bounded on $E$. This proves continuity of the
whole recursive map, including \eqref{eq:critical-output-representation}.

Continuous contraction now gives the rate by minimizing over its preimages.
For a measurable map $T$ and a reference probability $\rho$, the entropy
contraction identity is
\begin{equation}\label{eq:entropy-contraction-identity}
 \inf_{\eta:T_\#\eta=\nu}D(\eta\Vert\rho)
 =D(\nu\Vert T_\#\rho).
\end{equation}
The lower bound is data processing. When the right side is finite and
$q=d\nu/d(T_\#\rho)$, the measure $d\eta=q(T(u))\,d\rho(u)$ attains
equality. If the right side is infinite, data processing already proves
the identity in the extended real sense. Apply
\eqref{eq:entropy-contraction-identity} successively with
$T=\mathsf T_{Q_0}$ and $T=\mathsf T_{K(\nu_{s-1})}$.
The last minimization is
\[
 \inf_{v:Q^{1/2}v=z}\frac12\|v\|_2^2=\mathcal G_Q(z),
\]
including singular $Q$. This proves \eqref{eq:critical-rate}.

The contraction theorem supplies a good rate and exponential tightness.
For completeness, the displayed minimization is also directly well behaved:
relative entropy is jointly lower semicontinuous, the map
$\nu\mapsto\mathsf M_{K(\nu)}$ is weakly continuous, and
\[
 \mathcal G_Q(z)=\sup_{t\in\mathbb R^m}
 \left\{t^{\mathsf T}z-\frac12t^{\mathsf T}Qt\right\}
\]
is jointly lower semicontinuous in $(Q,z)$. Since $\mathcal P(E)^r$ is
compact, every finite infimum in \eqref{eq:critical-rate} is attained.
Moreover $K(\nu)\preceq mI_m$, so $I(z)\geq\|z\|_2^2/(2m)$.
\end{proof}

\begin{corollary}[Full-margin masses and posterior rates]
\label{cor:critical-margin-rates}
For a Boolean function $f$, regarded as a vector in $\{-1,1\}^m$, set
\[
 J(f)=\inf_{z:\ f_a z_a\geq1/2\ \forall a}I(z).
\]
Then the strict-margin event used throughout this paper has the exact rate
\begin{equation}\label{eq:critical-full-margin-rate}
 \lim_{N\to\infty}-\frac1N\log\mu_{N,1}(A_{N,f})=J(f).
\end{equation}
For a fixed noiseless training set $S$, let $\mathcal B_S$ be the Boolean
functions agreeing with its labels. If
$J_S=\min_{g\in\mathcal B_S}J(g)<\infty$, the actual globally conditioned
posterior satisfies, for every $f\in\mathcal B_S$,
\[
 \lim_{N\to\infty}-\frac1N\log\Pi_{N,1}(f\mid S)=J(f)-J_S.
\]
In particular, posterior mass concentrates on the minimizers of $J$ within
$\mathcal B_S$. These logarithmic asymptotics do not determine how mass is
divided between tied minimizers.
\end{corollary}

\begin{proof}
For $c\geq1$, retaining any feature laws in \eqref{eq:critical-rate} and
scaling only the output gives
\[
 I(cz)\leq c^2 I(z),
\]
because the entropy terms are nonnegative and the Gaussian energy scales
quadratically. If $f_a z_a\geq1/2$ for all $a$, then $cz$ has strict
margin for every $c>1$. Letting $c\downarrow1$ shows that the infimum of
$I$ over the open strict-margin orthant equals its infimum over the
closed orthant. The lower and upper large-deviation bounds therefore agree,
proving \eqref{eq:critical-full-margin-rate}.

Conditioning the prior on the global margin event and subsequently on the
labels gives exactly
\[
 \Pi_{N,1}(f\mid S)=
 \frac{\mu_{N,1}(A_{N,f})}
 {\sum_{g\in\mathcal B_S}\mu_{N,1}(A_{N,g})}.
\]
For fixed $d$ this is a finite sum, whose exponential rate is the smallest
$J(g)$ among its terms. This proves the ratio formula and concentration
claim.
\end{proof}


\section{Robust Gaussian target mass}\label{sec:cylinder}

This appendix supplies the target probability bound used by the robustness-based theorems, for every $N\ge C$ and every $0<t\le1$.
Take balanced clone groups of sizes $k_a$, so $N/k_a\le2C$, and let $P$
project onto group-constant vectors. Every unused circuit slot is filled by
a Boolean constant, so the ideal wire vector $b^s(x)$ lies in the range of
$P$ and has norm $\sqrt N$.
Put
\[
 e_*=(16\sqrt{2C})^{-1},\quad \tau_0=e_*/4,\quad
 b_L=\sqrt{16\log9+8\log(4L)},
\]
\[
 A=\max\{32b_L,2\log(8/e_*)\},\qquad
 \rho=\min\{1/4,A/[4(d+1)],Ae_*/(64\sqrt C)\}.
\]
Use tanh gate gain $A$ and readout coefficient $2$. The gain need not be
bounded by logical circuit width; it can be accommodated by polynomial
padding if invoking the ``weights bounded by width'' complexity convention.

Allow any $0<t\le1$. The standardized coordinates
$\sqrt d W^1$, $\sqrt N W^s$, and $Nv$ have common standard deviation $t$;
the exponent-$\gamma$ case is $t=N^{(1-\gamma)/2}$.
Restrict every first-layer coordinate, every
internal row's projection onto each normalized clone-group direction, and
every readout coordinate to an interval of radius $\rho t$ around its clone
value. Then
\[
 W^s=W_*^s+R^s+E^s,\qquad R^s=R^sP,
 \quad\|R^s\|_\op\le\rho t\sqrt C,\quad E^sP=0.
\]
The free Gaussian bulk $E^s$ is independent of the cylinder. A $1/4$-net
Gaussian operator bound makes $\|E^s\|_\op\le b_Lt$ simultaneously with
conditional probability at least $1/2$.

We check robustness to a separate per-layer RMS budget $\tau_0$, which is
stronger than the total-budget requirement for $r\le\tau_0$. Write
$e_s=\|h^s_\eta-b^s\|_2/\sqrt N$. If $e_{s-1}\le e_*$, each clone-group
mean has error at most $1/16$, so the teacher part of every intended signed
preactivation is at least $A/2$. For every real $u$ and $\epsilon\in\{-1,1\}$,
\[
 \epsilon a\ge A/2\quad\Longrightarrow\quad
 |\tanh(a+u)-\epsilon|\le2e^{-A/2}+8|u|/A.
\]
This is a global secant bound: for $|u|<A/4$ use saturation, and otherwise
use the trivial bound $2$. It handles concentrated adversarial noises.
Since $E^sb^{s-1}=0$ and noisy states have norm at most
$(1+\tau_0)\sqrt N$,
\[
 e_s\le2e^{-A/2}+(8/A)\rho\sqrt C(1+\tau_0)
 +(8b_L/A)e_{s-1}+\tau_0\le e_*.
\]
The first layer starts the invariant because its coordinatewise
preactivation error is at most $\rho(d+1)\le A/4$.
The readout has noisy signed value at least
$2-1/8-\rho(1+\tau_0)-\tau_0>1$.

There are $Nq$ constrained independent scalar coordinates, where
$q=d+2+(L-2)C$. Their squared centers sum to at most $NB_0'$, with
\[
 B_0'=9dA^2+18(L-2)CA^2+8C.
\]
The elementary Gaussian interval bound and the independent operator event give
\[
 -\log\mu_{N,t}(\mathcal R_{r,N,f})
 \le \frac{N}{t^2}\left[B_0'+q\log\frac{e\sqrt{2\pi}}{2\rho}+\log2\right]
 =DN/t^2.
\]
Here $D$ is jointly polynomially bounded in $d,C,L$. Exponentially small
coordinate precision would in any case cost only its logarithm in this
negative log-mass estimate; it should not be replaced unnecessarily by its
reciprocal.


\section{Robust posterior compression uniformly for \texorpdfstring{$\gamma\ge1$}{gamma at least one}}
\label{sec:robust-entropy}

This section gives a different statistical bridge. It uses robustness of the
original network while constructing a finite candidate, and tracks entropy
only in the original weights that have not yet been processed. Consequently
the intermediate compressed prefix need not have a density under a Gaussian
prior. The result applies to a posterior containing many different functions.

Let $\mu_{N,t}$, $0<t\le1$, have independent weight-entry variances
$t^2/d$, $t^2/N$, and $t^2/N^2$ in the first, internal, and readout layers.
The exponent-$\gamma$ prior is obtained by $t=N^{(1-\gamma)/2}$.
All-layer sign robustness refers to the total noise norm
in~\eqref{eq:noisenorm}; the clean $1/2$ margin may additionally be imposed.

\begin{theorem}[Uniform robust posterior compression]
\label{thm:robust-entropy}
Fix a scalar $1$-Lipschitz activation $|\phi|\le1$, and parameters
$d\ge1$, $L\ge2$, $K\ge1$, $0<r\le1$, and $0<\alpha\le1/8$.
There is a finite class $\mathcal H=\mathcal H(d,L,K,r,\alpha)$, consisting of
signs of bounded-weight depth-$L$ networks, such that
\[
 \log|\mathcal H|\le\poly(d,L,K,r^{-1},\alpha^{-1}).
\]
The polynomial is universal and does not depend on $N$ or $t$.
If $Q\ll\mu_{N,t}$ is supported on all-layer sign-$r$ robust networks and
\[
 D_{\rm KL}(Q\Vert\mu_{N,t})\le NK/t^2,
\]
then there is a distribution $\Gamma_Q$ on $\mathcal H$ such that, for every
$x\in\X$, a coupling of $w\sim Q$ and $h\sim\Gamma_Q$ satisfies
\[
 \Pp\{h(x)\ne\sgn z_w(x)\}\le\alpha.
\]
The coupling may depend on $x$; its two marginals do not. Neither the support
of $Q$ nor its samples need represent a single Boolean function.
\end{theorem}

We give the proof, including deliberately loose polynomial choices.
Two elementary entropy facts organize the construction. First, for a standard
$D$-dimensional Gaussian $G$, the entropy variational inequality gives
\begin{equation}\label{eq:entropy-moment}
 \E_P\|Z\|^2\le4D_{\rm KL}(P\Vert G)+2D\log2.
\end{equation}
Second, Gaussian Talagrand~\cite{talagrand} gives
$W_2^2(P,G)\le2D_{\rm KL}(P\Vert G)$.

\begin{lemma}[Entropy of the unprocessed weights]
\label{lem:future-entropy}
Let $A,X,Y$ be random variables, with fixed reference laws $P_X,P_Y$, and put
\[
 B=\E_A D_{\rm KL}(Q_{X,Y\mid A}\Vert P_X\otimes P_Y).
\]
Let $J$ be a finite message, possibly depending on $(A,X,Y)$, and let
$H=f(A,J,X)$. If $A'$ is a randomized function of $(A,J,H)$ whose fresh
randomness is conditionally independent of $Y$, then
\[
 \E_{A'}D_{\rm KL}(Q_{Y\mid A'}\Vert P_Y)\le B+H(J\mid A).
\]
\end{lemma}
\begin{proof}
The left side equals $D_{\rm KL}(Q_Y\Vert P_Y)+I(Y;A')$.
Data processing and
\[
 I(Y;J,H\mid A)\le I(Y;X,J\mid A)
 \le I(Y;X\mid A)+H(J\mid A)
\]
bound it by $\E_{A,X}D_{\rm KL}(Q_{Y\mid A,X}\Vert P_Y)+H(J\mid A)$,
which is at most the claimed bound by the KL chain rule. These expressions
can equivalently be justified by finite partitions and lower semicontinuity.
\end{proof}

Here $A$ will be a constructed prefix, $X$ the current original weight block,
and $Y$ all later original blocks. The lemma imposes no reference law on $A$.

\begin{lemma}[Whole-layer active directions]
\label{lem:block-entropy}
Condition on $(A,Y)$ and let $X$ be a standardized Gaussian-reference matrix,
with conditional relative entropy $R(A,Y)$. For vectors $b_x(A)$ of norm at
most one, one can select a list $J$ of inputs and its span
$U=\operatorname{span}\{b_x:x\in J\}$ with
\[
 \dim U\le\frac{t^2R(A,Y)}{N\delta}
\]
such that preserving $H=tXU$ and replacing the orthogonal coordinates by fresh
Gaussian coordinates changes the physical evaluation vectors $tXb_x$, under
a suitable coupling, by expected squared RMS at most $2\delta$ at each $x$.
The candidate matrix law preserves $(A,Y,H)$ and does not depend on $x$.
\end{lemma}
\begin{proof}
Use a fixed ordering of the finite input set. Whenever the averaged conditional
KL of the entire residual evaluation vector $Xv$ exceeds $N\delta/t^2$, add
the unit residual direction $v$ of that input. Each query is selected from
the conditional law given $(A,Y)$, averaging over active-coordinate values;
thus $J$ is a function of $(A,Y)$. Orthogonal Gaussian coordinates
are independent, so the active-coordinate entropy increases by precisely this
conditional KL. The displayed rank bound follows. At termination, Gaussian
Talagrand on the entire evaluation vector, multiplication by the squared
residual-vector length, and division by $N/t^2$ give the error bound.
Scalar-vector couplings extend to full matrices by disintegration.
The basis of $U$ is a measurable function of $A,J$ under fixed Gram--Schmidt;
it has no further dependence on $Y$ once those variables are supplied.
\end{proof}

\paragraph{Parameter choices.}
Write
\begin{align*}
 e&=\alpha r^2/(1000L),& \eta&=\alpha/(100L),&
 q&=\lceil2K/(e\eta)\rceil,\\
 D_0&=[(d+1)q+2]\log2,& V&=8K+2q,& B_H^2&=4V/e,\\
 T&=128(q+1)(B_H^2+1)/e,& a&=e/[128(B_H^2+1)],\\
 K_0&=(4K+2\log2)/\eta,& B_v&=16K_0/r.
\end{align*}
Choose $m$ to be the least positive integer no smaller than both
\begin{align*}
 &10000T^4(q+1)^2a^{-2}[d+\log(10(q+1)/\eta)+1],\\
 &10000B_v^2r^{-2}[d+\log(10/\eta)+1],
\end{align*}
and put
\[
 N_0=\max\{m,\lceil LD_0/K\rceil,1\},\quad
 B_G=\max\{2,2\sqrt{\log(3m/e)}\},\quad
 B_1^2=4[4K+2(d+1)\log2]/e.
\]
All quantities have a fixed polynomial upper bound in the theorem parameters.

\paragraph{The small-width case.}
If $N\le N_0$, let
$P_0=N_0(d+1)+(L-2)N_0^2+N_0$. Applying~\eqref{eq:entropy-moment} to all
standardized weights, and then multiplying by $t^2$, gives
\[
 \E_Q\|tZ_{\rm all}\|^2\le4N_0K+2P_0\log2.
\]
With probability at least $1-\alpha/4$, all actual weights are bounded by
\[
 B_{\rm small}=[4(4N_0K+2P_0\log2)/\alpha]^{1/2}.
\]
Keep the original network on this event and use a default classifier otherwise.
Padding with nodes having zero outgoing weights preserves the original output.
The successful candidate has output margin greater than $r$.
It remains to treat $N>N_0$.

\paragraph{The virtual prefix and its entropy budget.}
The constructed prefix $A$ has $N$ virtual nodes in its current layer. Each
virtual row acts on $m$ stored nodes of the preceding layer. The auxiliary
description may keep all $N$ rows; the final network will keep just the rows
selected at the following stage. Track the conditional KL of the unprocessed
\emph{original} weights. By Lemma~\ref{lem:future-entropy}, it is at most
\[
 B_s\le NK/t^2+sD_0\le2NK/t^2.                              \tag{*}
\]
This is not a KL bound for the singular compressed prefix itself.

For the first layer, zero every physical standardized row $tZ_{1,i}$ of norm
greater than $B_1$, and retain all other rows. Bounded activations and
\eqref{eq:entropy-moment} bound the expected squared RMS activation change
at each input by $e$. This deterministic prefix channel does not increase the
future entropy budget.

\paragraph{An internal layer: retain active data.}
Condition on the current prefix $A$ and the future original weights $Y$.
Apply Lemma~\ref{lem:block-entropy} to
$b_x=h_A(x)/\sqrt N$ and $\delta=e$. Abort if more than $q$ inputs are needed.
By (*) and Markov, this failure has probability at most $\eta$.
The successful query list, together with an abort flag, has a code of at most
$[(d+1)q+2]$ bits. Its entropy is at most $D_0$.

Let $H=tXU$ be the physical active matrix. Averaging the Gaussian entropy
moment bound, with the success indicator rather than conditioning on success, gives
\[
 \E\left[\frac{\|H\|_F^2}{N};\ \text{success}\right]
 \le4t^2B_s/N+2t^2q\log2\le V.
\]
Zero every row $H_i$ whose norm exceeds $B_H$. The expected squared RMS
activation change from this operation is at most $4V/B_H^2=e$.
All rows are clipped before any later selection, including selections that
may depend on future weights.

\paragraph{An internal layer: reconstruct using $m$ source nodes.}
For the orthonormal active basis let $u(j)=\sqrt N\,U_{j,:}$; uniform source
sampling gives $\E_j u(j)u(j)^\top=I$.
Clip the \emph{whole vector} at radius $T$:
$c(j)=u(j)\mathbf1_{\{\|u(j)\|\le T\}}$. Then
\[
 \E cc^\top\preceq I,\qquad \E\|u-c\|\le q/T.
\]
Sample $m$ uniform source indices $S_j$. Entrywise Hoeffding and a union bound
give, with probability at least $1-\eta$,
\begin{align*}
 \|\widehat G-\E cc^\top\|_{\op}&\le a,\quad
 \sup_x\|\widehat v(x)-\E c h_A(x)\|\le a,\\
 \sup_{x,y}|\widehat K(x,y)-K(x,y)|&\le a,
\end{align*}
where $\widehat G=m^{-1}\sum_jc(S_j)c(S_j)^\top$,
$\widehat v(x)=m^{-1}\sum_jc(S_j)h_{A,S_j}(x)$, and $\widehat K$ is the
empirical feature covariance; its population counterpart is
$K(x,y)=N^{-1}\sum_i h_{A,i}(x)h_{A,i}(y)$.
Abort on failure; this event depends only on
$A,J,S$, so it does not reveal further future information.

Set $\widetilde c_j=c(S_j)/\sqrt{1+a}$, and let $D$ have rows
$\widetilde c_j^\top/\sqrt m$. Then $D^\top D\preceq I$ and
$I-DD^\top\succeq0$. Independently for every target virtual row draw
$\xi_i\sim N(0,I-DD^\top)$ and use coefficients
\begin{equation}\label{eq:finite-row}
 \widehat w_{ij}=m^{-1}H_i^{\rm clip}\cdot\widetilde c_j
                         +t\xi_{ij}/\sqrt m.
\end{equation}
These are actual neural-network rows on the $m$ sampled nodes.
Their active feature vector is
$\widehat\psi(x)=m^{-1}\sum_j\widetilde c_jh_{A,S_j}(x)$, and the lazy
covariance is exactly $t^2(\widehat K-\widehat\psi^\top\widehat\psi)$.
For the ideal field, write $\psi(x)=U^\top b_x$ and
$C(x,y)=K(x,y)-\psi(x)^\top\psi(y)$. Direct calculation gives
\[
 \sup_x\|\widehat\psi(x)-\psi(x)\|\le q/T+2a=:\zeta,
 \qquad \sup_{x,y}|\widehat C(x,y)-C(x,y)|\le a+2\zeta.
\]
Here $\widehat C=\widehat K-\widehat\psi^\top\widehat\psi$.
Bessel's inequality and $D^\top D\preceq I$ give
$\|\psi(x)\|,\|\widehat\psi(x)\|\le1$, which justify the covariance bound.
At a fixed input, coupling scalar Gaussian evaluations with a common standard
normal bounds the squared preactivation transport error, per row, by
\[
 B_H^2\zeta^2+t^2(a+2\zeta)\le e/8.
\]
Extend these scalar couplings to full row fields. The candidate weight law is
the same at every input; there is no claim of a single Gaussian-field coupling
that is close at all inputs simultaneously.

Finally clip \emph{all} coordinates $\xi_{ij}$ to $[-B_G,B_G]$ before later
row selection. Since each has variance at most one,
\[
 \E\left|\frac{t}{\sqrt m}\sum_j
   (\xi_{ij}-\operatorname{clip}_{B_G}\xi_{ij})h_{A,S_j}(x)\right|^2
 \le3m e^{-B_G^2/4}\le e.
\]
Thus every virtual coefficient is bounded by $B_HT/m+B_G/\sqrt m$,
without a union bound over $N$.

All newly constructed data are a randomized function of $A,J,H$.
Lemma~\ref{lem:future-entropy} therefore proves (*) for the next step.
Combining entropy replacement, active clipping, finite reconstruction, and
Gaussian clipping by $(a_1+\cdots+a_4)^2\le4\sum_i a_i^2$, the expected
squared RMS activation injection at this layer is less than $20e$.

\paragraph{The readout and final network.}
The original readout marginal has never changed. Its physical moment satisfies
$\E_Q[N\|v\|^2]\le4K+2\log2$.
Abort if $N\|v\|^2>K_0$, an event of probability at most $\eta$.
Zero coefficients with $|Nv_i|>B_v$. This changes the output uniformly on
bounded virtual feature vectors by at most $K_0/B_v=r/16$.
Sampling $m$ uniform last-layer sources and using readout coefficients
$(Nv_i\mathbf1_{|Nv_i|\le B_v})/m$ adds at most $r/16$ uniformly on $\X$,
except with probability $\eta$, by Hoeffding and the chosen $m$.

At each stage, retain in the final network just the $m$ source rows selected
from the preceding virtual layer. Each such row already acts on the previously
stored $m$ nodes. The result is an ordinary width-$m$ network; repeated indices
are duplicate nodes. Its weights are bounded by
\[
 B_{\rm large}=\max\{B_1,B_HT/m+B_G/\sqrt m,B_v/m,1\}.
\]
The larger virtual construction is not part of the final predictor.

\paragraph{The original network receives all perturbations.}
Fix an input $x$ and glue the conditional couplings, retaining the original
weight history. The current and future original blocks are untouched until
their own stages. Define
\[
 \eta^s=h^s_{\rm virtual}(x)
           -\phi(W^s_{\rm original}h^{s-1}_{\rm virtual}(x)).
\]
The analogous readout difference is $\eta^L$. Later stages leave earlier
computational prefixes unchanged, so the final candidate output is exactly
the \emph{original} network's output under these postactivation injections.
No intermediate robustness assumption is made.

For precision, the canonical construction draws its new Gaussian coefficients
independently of the discarded current block given $(A,J,H)$. The comparison
coupling may correlate them with that discarded block. It preserves the same
joint marginal of the new prefix and the future original blocks, as well as
the original full-network marginal, but need not preserve the canonical joint
law with all discarded weights. The next entropy step uses the first marginal;
original-network robustness uses the second. There is no conflicting demand
for independence in the comparison coupling itself.

More explicitly, the induction invariant is that the comparison law has
full original-weight marginal $Q$ and the canonical marginal of
$(A,X,Y)$. The conditional transports supply a joint law of
$(A,X,Y,A')$ with this incoming marginal and the prescribed outgoing
$(A',Y)$ marginal. Extend it by the previous conditional law of the past
original weights given $(A,X,Y)$. The extension keeps the complete original
weight law unchanged, preserves all earlier computational prefixes, and
establishes the same invariant for the next layer. This is ordinary
disintegration on finite-dimensional standard Borel spaces.

Abort probability is at most $2L\eta$. The unconditional moment bound gives
\[
 \Pp\left\{\text{no abort and }\sum_{s<L}\|\eta^s\|^2/N>r^2/8\right\}
 \le160Le/r^2=0.16\alpha.
\]
Otherwise the readout injection is at most $r/8$, and total squared noise is
at most $9r^2/64<r^2/4$. Robustness then gives candidate margin greater than
$r/2$ relative to the original network's label: a smaller margin could be
cancelled by an additional output shift within the radius-$r$ noise ball.
The exceptional probability is less than $\alpha/4$.

\paragraph{A fixed finite class.}
Pad candidates to width $M=N_0$ and take
$B=\max(B_{\rm small},B_{\rm large},1)$. Round every parameter to a fixed
grid on $[-B,B]$ of mesh at most
\[
 h=\frac{r}{16(d+L+1)(2M(B+1))^L}.
\]
The bounded-activation, Lipschitz layer recursion gives uniform output change
at most $r/8$. Thus rounding preserves every successful coupled prediction.
Writing $P=M(d+1)+(L-2)M^2+M$, the grid class satisfies
\[
 \log|\mathcal H|\le
 P\log\left[3+\frac{32B(d+L+1)(2M(B+1))^L}{r}\right].
\]
Although grid precision is exponentially small in $L$, its number of bits is
polynomial in $L$. All other bounds above are polynomial too. Rounding is a
fixed map on candidate networks, giving the common marginal $\Gamma_Q$ and
proving Theorem~\ref{thm:robust-entropy}.

\begin{corollary}[Polynomial-depth learning at every width]
\label{cor:robust-entropy-pac}
Suppose the raw Gaussian target event, consisting of networks with all-layer
sign margin $r$ and clean margin $1/2$ for $f$, has mass at least
$\exp(-NK/t^2)$. Restrict the prior to all sign-$r$ robust networks having
global clean margin $1/2$, regardless of their Boolean function, then condition
on $n$ noiseless iid training examples. With probability at least $1-\beta$,
the resulting posterior satisfies
\[
 \operatorname{GibbsErr}(Q)
 \le2\alpha+\sqrt{\frac{\log(2|\mathcal H|/\beta)}{2n}}.
\]
Consequently the sample bound is jointly polynomial in
$d,L,K,r^{-1}$ and inverse desired error, uniformly in $N$ and $0<t\le1$.
For posterior-confidence error at threshold $2/3$, multiply Gibbs error by
three; for threshold $1/2$, multiply it by two.
\end{corollary}
\begin{proof}
The posterior is the raw Gaussian conditioned on an event containing the target
event, so its KL is at most $NK/t^2$. Its training error is zero. The compression
couplings imply that $\Gamma_Q$ has expected empirical error at most $\alpha$.
Finite-class uniform convergence holds for every member of $\mathcal H$, hence
also for this sample-dependent mixture. Its true error is at most $\alpha$
plus the displayed square-root term. The pointwise coupling adds at most one
more $\alpha$ to recover original posterior Gibbs error. This argument does
not require a global coupling over inputs or a common function for all draws.
The confidence-error statements follow from Markov on posterior error
probability at a random test input.
\end{proof}

\section{A finite-radius certificate using clean Jacobian products}
\label{sec:clean-transport-certificate}

Products of individual spectral norms discard cancellations between layers.
The following sufficient test retains the exact clean transition products and
bounds the nonlinear remainder separately. It is a finite-radius statement,
not an inference from the infinitesimal Jacobian alone.

Write the $H=L-1$ hidden layers in column-vector convention as
$h_s=\tanh(W_sh_{s-1})$, with output $z=v^\top h_H$ and width $N$.
Fix an input and its clean preactivations $a_s=W_sh_{s-1}$. Define
\[
 A_s=\operatorname{diag}(\operatorname{sech}^2 a_s)W_s,
 \qquad P_{s,t}=A_s\cdots A_{t+1}\ (t<s),\quad P_{s,s}=I.
\]
The Gramian for normalized additive hidden disturbances is
\[
 G_1=NI,\qquad G_s=A_sG_{s-1}A_s^\top+NI.
\]
For a proposed total-noise radius $r$, recursively put $u_1=0$ and, for $s\ge2$,
\begin{align}
 b_s&=r\sqrt{\operatorname{diag}(W_sG_{s-1}W_s^\top)}
       +\sum_{t=2}^{s-1}|W_sP_{s-1,t}|u_t,\label{eq:transport-b}\\
 u_s&=\frac{2}{3\sqrt3}\,b_s^{\odot2}.\label{eq:transport-u}
\end{align}
Absolute values, the square root, and the square in these expressions are
coordinatewise. Finally define
\begin{equation}\label{eq:transport-output}
 B(r,x)=r\sqrt{v^\top G_Hv+1}
                  +\sum_{t=2}^H|v^\top P_{H,t}|u_t.
\end{equation}

\begin{proposition}[Finite-radius transport certificate]
Every simultaneous postactivation and readout disturbance satisfying
$\sum_{s=1}^H\|\eta_s\|^2/N+\eta_{\rm out}^2\le r^2$ changes the output
at $x$ by at most $B(r,x)$. Consequently, if $B(r,x)<|z(x)|$ on the full
input cube, the network preserves its own clean signs throughout this ball.
\end{proposition}
\begin{proof}
Let $\Delta h_s$ be the hidden displacement and write its exact recursion as
\[
 \Delta h_s=A_s\Delta h_{s-1}+\eta_s+R_s,
 \qquad
 \Delta h_s=\sum_{t=1}^sP_{s,t}\eta_t+\sum_{t=2}^sP_{s,t}R_t.
\]
There is no input perturbation, so $R_1=0$. The global scalar bound
$|\tanh''|\le4/(3\sqrt3)$ implies
$|R_s|\le[2/(3\sqrt3)]|W_s\Delta h_{s-1}|^{\odot2}$.
Assuming $|R_t|\le u_t$ at preceding layers, Cauchy--Schwarz applied to the
linear disturbance map, followed by the coordinatewise triangle inequality
for the remainder terms, gives $|W_s\Delta h_{s-1}|\le b_s$.
This proves $|R_s|\le u_s$ inductively. At the readout, Cauchy--Schwarz
including scalar output noise gives the first term of~\eqref{eq:transport-output};
the remainder sum gives the second. No statistical independence or
small-noise approximation is used.
\end{proof}

The test can be sharpened by replacing the global curvature bound in
\eqref{eq:transport-u} with the supremum of $|\tanh''|$ on
$[a_{s,j}-b_{s,j},a_{s,j}+b_{s,j}]$. An alternative exact scalar remainder
envelope evaluates
$\tanh(a+q)-\tanh(a)-\operatorname{sech}^2(a)q$
at $q=\pm b$ and at $q=-2a$ when $|2a|\le b$: these and $q=0$ exhaust
its extrema on the interval. The global-curvature version above already
suffices for the numerical improvement below.


The numerical condition in Section~\ref{sec:numerical-stability} is certified at an input whenever $B(r_{\rm num},x)\le1/16$. Floating-point evaluations of this sufficient test are distinguished from interval-verified statements in the experimental section.


\begin{thebibliography}{9}
\bibitem{jacot2023} A. Jacot.
Implicit bias of large depth networks: a notion of rank for nonlinear functions.
\emph{ICLR}, 2023. \url{https://arxiv.org/abs/2209.15055}.
\bibitem{parkinson2024} S. Parkinson, G. Ongie, R. Willett, O. Shamir,
and N. Srebro. Depth separation in norm-bounded infinite-width neural networks.
\emph{COLT}, 2024. \url{https://proceedings.mlr.press/v247/parkinson24a.html}.
\bibitem{andreis2026} L. Andreis, F. Bassetti, and C. Hirsch.
LDP for the covariance process in fully connected Gaussian neural networks.
\emph{Electronic Journal of Probability} 31, 2026.
\url{https://doi.org/10.1214/26-ejp1477}.
\bibitem{dembozeitouni} A. Dembo and O. Zeitouni.
\emph{Large Deviations Techniques and Applications}. Second edition,
Springer, 1998.
\bibitem{borell1975} C. Borell.
The Brunn--Minkowski inequality in Gauss space.
\emph{Inventiones Mathematicae} 30, 207--216, 1975.
\url{https://doi.org/10.1007/BF01425510}.
\bibitem{previous} ChatGPT-6 Astra and Claude (Fable 5 and Opus 5). \emph{Depth dependence, internal robustness,
and Bayesian neural-network learning}. Research report for K. H\"anni,
17 October 2026.
\url{https://github.com/kaarelh/AI-works/blob/main/neural-network-generalization/works-in-progress/depth-supporting-notes/round1-depth-report/paper.pdf}.
The compression and target-cylinder proofs are reproduced in this report.
\bibitem{fable} Claude Fable 5. \emph{Robustness kills the exp(depth):
Compression, depth, and neural-network priors}. User-supplied research note,
\texttt{robust-depth.pdf}.
\bibitem{golowich2018} N. Golowich, A. Rakhlin, and O. Shamir.
Size-independent sample complexity of neural networks.
\emph{COLT}, 2018. \url{https://proceedings.mlr.press/v75/golowich18a.html}.
\bibitem{nakhleh2026} J. Nakhleh and R. D. Nowak.
Deep Neural Variation Spaces: A Unifying Perspective on Depth and Complexity.
2026. \url{https://arxiv.org/abs/2607.05546}.
\bibitem{arora} S. Arora, R. Ge, B. Neyshabur, and Y. Zhang.
Stronger generalization bounds for deep nets via a compression approach.
\emph{ICML}, 2018. \url{https://proceedings.mlr.press/v80/arora18b.html}.
\bibitem{weima2020} C. Wei and T. Ma.
Improved sample complexities for deep neural networks and robust classification
via an all-layer margin. \emph{ICLR}, 2020.
\url{https://arxiv.org/abs/1910.04284}.
\bibitem{talagrand} M. Talagrand.
Transportation cost for Gaussian and other product measures.
\emph{Geometric and Functional Analysis} 6, 587--600, 1996.
\end{thebibliography}
\end{document}
